%% Question bank chapter file - generated by bank_tools.py build (do not edit while the bank is still being built:
%% it is regenerated from the per-test data in Question Bank\data). Every question is one bankq block:
%%   \begin{bankq}[id=..., source={...}, paper=.., qno=.., type=..., subtopic={...}, answer={...}, answer_src=key|solved,
%%                 working={...}, difficulty=1-10, diff_parts={K S I C F}, diff_note={...},
%%                 group=<passage id or empty>, also_in={duplicate copies in other tests}]
%% difficulty = min(10, K+S+I+C+F-1): K concepts 1-3, S steps 1-4, I insight 0-2, C computation 0-2, F format 0-1
%%   ...question LaTeX, options, figures (\bfig / \bfigi -> img/)...
%%   \end{bankq}
%% Passages shared by several questions are bankpassage blocks; their questions carry group=<passage id>.
\documentclass[10pt]{article}
\usepackage[a4paper,margin=1.8cm]{geometry}
\usepackage[T1]{fontenc}
\usepackage{amsmath,amssymb,mathtools}
\usepackage{graphicx}
\usepackage[export]{adjustbox}
\usepackage{xcolor,needspace,keyval}
\usepackage{enumitem}
\usepackage{booktabs,array,tabularx}
\usepackage[hidelinks]{hyperref}
\setlength{\parindent}{0pt}
\setlength{\parskip}{3pt}
\makeatletter
\define@key{bq}{id}{\def\bq@id{#1}}\define@key{bq}{source}{\def\bq@src{#1}}\define@key{bq}{paper}{\def\bq@paper{#1}}
\define@key{bq}{qno}{\def\bq@qno{#1}}\define@key{bq}{type}{\def\bq@type{#1}}\define@key{bq}{subtopic}{\def\bq@sub{#1}}
\define@key{bq}{answer}{\def\bq@ans{#1}}\define@key{bq}{answer_src}{\def\bq@asrc{#1}}\define@key{bq}{working}{\def\bq@work{#1}}
\define@key{bq}{group}{\def\bq@group{#1}}\define@key{bq}{also_in}{\def\bq@also{#1}}
\define@key{bq}{difficulty}{\def\bq@diff{#1}}\define@key{bq}{diff_parts}{\def\bq@dparts{#1}}\define@key{bq}{diff_note}{\def\bq@dnote{#1}}
\newenvironment{bankq}[1][]{\def\bq@work{}\def\bq@also{}\def\bq@group{}\def\bq@diff{?}\def\bq@dparts{}\def\bq@dnote{}\setkeys{bq}{#1}%
  \par\needspace{4\baselineskip}\medskip\noindent{\small\textbf{\bq@id}\hfill\color{gray}[\bq@src, Paper \bq@paper, Q\bq@qno]
  \bq@type\ $\cdot$ \bq@sub\ $\cdot$ \textbf{D\bq@diff}}\par\nopagebreak\relax}
  {\par{\small\color{teal}Answer: \textbf{\bq@ans}\ifx\bq@work\@empty\else\ (solved: \bq@work)\fi
  \ifx\bq@also\@empty\else\par\color{gray}Also in: \bq@also\fi}\par}
\define@key{bp}{id}{\def\bp@id{#1}}\define@key{bp}{source}{\def\bp@src{#1}}\define@key{bp}{questions}{\def\bp@qs{#1}}
\newenvironment{bankpassage}[1][]{\setkeys{bp}{#1}\par\needspace{4\baselineskip}\bigskip\noindent
  {\small\textbf{Passage \bp@id}\hfill\color{gray}[\bp@src; questions \bp@qs]}\par\nopagebreak\relax\itshape}{\par}
\makeatother
\newcommand{\bfig}[2]{\par\smallskip{\centering\includegraphics[width=#1pt,max width=0.62\linewidth,max height=0.32\textheight]{img/#2}\par}\smallskip}
\newcommand{\bfigi}[2]{\includegraphics[width=#1pt,max width=\linewidth,max height=0.32\textheight,valign=c]{img/#2}}
\newcommand{\bmissing}[1]{\textit{[figure missing in the source conversion: \detokenize{#1}]}}
\newenvironment{qoptions}{\begin{itemize}[leftmargin=1.8em,itemsep=1pt,topsep=2pt,label={}]}{\end{itemize}}
\newcommand{\opt}[2]{\item[(#1)] #2}
\providecommand{\fwstart}{}\providecommand{\fwend}{}\providecommand{\qfullsplit}{}
\newcolumntype{L}[1]{>{\raggedright\arraybackslash}p{#1}}
\newcolumntype{C}[1]{>{\centering\arraybackslash}p{#1}}
\newcolumntype{R}[1]{>{\raggedleft\arraybackslash}p{#1}}
\providecommand{\sectionhdr}[1]{}\providecommand{\subsectionhdr}[1]{}


\begin{document}
{\Large\bfseries Ray Optics --- Question Bank}\par\smallskip
{\small Physics chapter P19. Questions from all converted test papers, grouped by source series. Answers: printed key where it matches the paper, otherwise solved (answer\_src=solved, with working). Exact repeats across papers appear once, with also\_in listing the other copies.}\par\medskip\hrule

\section*{Score-I papers}
\begin{bankq}[id=SI12-P1-PHY-S1-Q05, source={Score-I - Test 12 (09 Mar 2022)}, paper=1, qno=5, type=multi, subtopic={Lenses}, answer={A,C}, answer_src=key, difficulty=4, diff_parts={K2 S2 I0 C0 F1}, diff_note={concepts: lens maker formula, power of a lens-mirror combination | steps: focal length of the lens from the radii; equivalent power P = 2PL + PM; nature and magnitude of the system}]
In the figure shown the radius of curvature of the left \& right surface of the concave lens are 10 cm \& 15 cm respectively. The radius of curvature of the mirror is 15 cm:
\bfig{126.6}{SI12_p1_phy_s1_q05_1.png}
\begin{qoptions}
\opt{A}{modulus of equivalent focal length of the combination is 18 cm}
\opt{B}{modulus of equivalent focal length of the combination is 36 cm}
\opt{C}{the system behaves like a concave mirror}
\opt{D}{the system behaves like a convex mirror}
\end{qoptions}
\end{bankq}
\begin{bankq}[id=SI12-P2-PHY-S3-Q10, source={Score-I - Test 12 (09 Mar 2022)}, paper=2, qno=10, type=matching, subtopic={Reflection}, answer={P$\to$3, Q$\to$2, R$\to$1}, answer_src=key, difficulty=4, diff_parts={K2 S2 I0 C0 F1}, diff_note={concepts: reflection, refraction deviation | steps: deviation graph for a mirror; rarer to denser; denser to rarer including total internal reflection}]
Match the entries of list-I with the entries of list-II, which describes the angle of deviation of ray with angle of incidence.
\begin{center}\footnotesize
\begin{tabular}{|p{0.6cm}|p{2.9cm}|p{0.6cm}|p{2.4cm}|}
\hline
 & \centering List-I & \multicolumn{2}{c|}{List-II}\\
\hline
(P) & A ray is falling on a plane smooth mirror & (1) & \bfigi{63.4}{SI12_p2_phy_s3_q10_1.png}\rule{0pt}{32pt}\\
\hline
(Q) & A ray is going from a rarer to denser medium & (2) & \bfigi{67.1}{SI12_p2_phy_s3_q10_2.png}\rule{0pt}{32pt}\\
\hline
(R) & A ray is going from a denser to rarer medium & (3) & \bfigi{67.1}{SI12_p2_phy_s3_q10_3.png}\rule{0pt}{32pt}\\
\hline
\end{tabular}
\end{center}
\end{bankq}

\section*{Score Advanced papers}
\begin{bankq}[id=SA01-P1-PHY-S3-Q08, source={Score Advanced - Test 01 (22 Sep 2020)}, paper=1, qno=8, type=numerical, subtopic={Reflection}, answer={13.335}, answer_src=key, difficulty=7, diff_parts={K2 S2 I2 C1 F1}, diff_note={concepts: convex mirror magnification and displacement, reflection | steps: mirror formula; two positions; focal length}]
The image of an object formed by a convex spherical mirror is magnified one third in comparison with the object. If the mirror is displaced by 40 cm along its principal optical axis in the direction of light, the image is again found one-third in comparison with the object. Find the radius of curvature of the mirror (in cm).
\end{bankq}
\begin{bankq}[id=SA02-P1-PHY-S2-Q04, source={Score Advanced - Test 02 (13 May 2021)}, paper=1, qno=4, type=single, subtopic={Lenses}, answer={D}, answer_src=solved, working={Lens gives real inverted magnified image; mirror gives virtual erect magnified; refracting surface gives virtual erect magnified; diverging lens gives virtual erect diminished}, difficulty=5, diff_parts={K2 S2 I1 C1 F0}, diff_note={concepts: images by lenses mirrors and surface, image position for each device | steps: image position for each device; magnification; nature}]
Match the list--I with nature of image in list--II (Consider all optical devices are kept in air)\par
\noindent{\footnotesize\setlength{\tabcolsep}{2pt}\begin{tabular}{|L{.7cm}|L{3.8cm}|L{.7cm}|L{1.8cm}|}\hline
 & List-I & & List-II\\\hline
(P) & \bfigi{103.8}{SA02_p1_phy_s2_q04_1.png} & (1) & Erect image\\\hline
(Q) & \bfigi{105.5}{SA02_p1_phy_s2_q04_2.png} & (2) & Inverted image\\\hline
(R) & \bfigi{104.1}{SA02_p1_phy_s2_q04_3.png} & (3) & magnified image\\\hline
(S) & \bfigi{106.6}{SA02_p1_phy_s2_q04_4.png} & (4) & diminished image\\\hline
 & & (5) & real image\\\hline
 & & (6) & virtual image\\\hline
\end{tabular}}\par
\begin{qoptions}
\opt{A}{P $\to$ 2,3,5;Q $\to$ 1,4,6;R $\to$ 1,3,6;S $\to$ 2,3,5}
\opt{B}{P $\to$ 1,4;Q $\to$ 2,3,5;R $\to$ 1,4,6;S $\to$ 1,3,6}
\opt{C}{P $\to$ 1,4,6;Q $\to$ 1,3,6;R $\to$ 2,3,5;S $\to$ 1,4}
\opt{D}{P $\to$ 2,3,5;Q $\to$ 1,3,6;R $\to$ 1,3,6;S $\to$ 1,4,6}
\end{qoptions}
\end{bankq}
\begin{bankq}[id=SA02-P2-PHY-S3-Q05, source={Score Advanced - Test 02 (13 May 2021)}, paper=2, qno=5, type=numerical, subtopic={Lenses}, answer={2.75}, answer_src=solved, working={Lens image 6 mm below the lens axis; virtual object 20 cm behind the mirror, magnification 0.5; measured from the mirror axis 2.75 mm}, difficulty=8, diff_parts={K2 S2 I2 C2 F1}, diff_note={concepts: object imaged by a lens and a mirror, lenses | steps: lens image; virtual object; mirror image; offset of axes}]
A concave mirror of focal length 20 cm and a convex lens of focal length 10 cm are kept with their optic axis parallel but separated by 0.5 mm as shown in figure. The distance between lens and mirror is 10 cm. An object of height 3 mm is placed on the optic axis of lens at distance 15 cm from the lens. Find height of image formed by mirror after refraction from lens (in mm).\par
\bfig{157.3}{SA02_p2_phy_s3_q05_1.png}
\end{bankq}
\begin{bankq}[id=SA05-P1-PHY-S1-Q09, source={Score Advanced - Test 05 (15 Jul 2021)}, paper=1, qno=9, type=single, subtopic={Refraction}, answer={C}, answer_src=key, difficulty=3, diff_parts={K1 S1 I1 C1 F0}, diff_note={concepts: statements on reflection and refraction | steps: total internal reflection needs dense to rare; focal length of a mirror does not depend on the medium}]
Select the correct statement :-
\begin{qoptions}
\opt{A}{Total internal reflection takes place when light travel from rarer to denser medium.}
\opt{B}{Focal length of a convex lens is positive irrespective of medium in which it is placed.}
\opt{C}{Focal length of spherical mirror depends on its radius of curvature not on the surrounding medium.}
\opt{D}{Blue colour of sky appear due to total internal reflection}
\end{qoptions}
\end{bankq}
\begin{bankq}[id=SA05-P1-PHY-S1-Q10, source={Score Advanced - Test 05 (15 Jul 2021)}, paper=1, qno=10, type=single, subtopic={Reflection}, answer={D}, answer_src=key, difficulty=7, diff_parts={K3 S2 I2 C1 F0}, diff_note={concepts: image in a silvered glass hemisphere, reflection, first refraction gives parallel rays for object at 2R | steps: first refraction gives parallel rays for object at 2R; reflect from the plane face; refraction out of the curved surface}]
A glass hemisphere of radius R and material having refractive index 1.5 is silvered on its flat face as shown in the figure. A small object of height h is located at a distance 2R from surface of hemisphere. The final image will form :-\par
\bfig{123.5}{SA05_p1_phy_s1_q10_1.png}
\begin{qoptions}
\opt{A}{at a distance of R from silvered surface, on right side}
\opt{B}{on the hemispherical surface}
\opt{C}{at the plane surface}
\opt{D}{at a distance of 2 R from the curved surface on left side}
\end{qoptions}
\end{bankq}
\begin{bankq}[id=SA06-P1-PHY-S1-Q01, source={Score Advanced - Test 06 (04 Sep 2021)}, paper=1, qno=1, type=multi, subtopic={Refraction}, answer={A,C}, answer_src=key, difficulty=7, diff_parts={K2 S2 I2 C1 F1}, diff_note={concepts: dispersion at a spherical refracting surface, refraction | steps: n/v=(n-1)/R for a parallel beam; compute v for violet and red; difference}]
A spherical surface separates air \& medium for which $\mu$ = 1.615 for violet and $\mu$ = 1.600 for red color. A paraxial beam parallel to optic axis is incident on the surface as shown. The distance between point of convergence for violet and red color is $\Delta$f.\par
\bfig{102.7}{SA06_p1_phy_s1_q01_1.png}
\begin{qoptions}
\opt{A}{$\Delta$f = 0.41 cm}
\opt{B}{Point of convergence for red is closer to optical centre than that for violet.}
\opt{C}{Point of convergence for violet is closer to optical centre than that for red.}
\opt{D}{$\Delta$f = 0.84 cm.}
\end{qoptions}
\end{bankq}
\begin{bankq}[id=SA06-P2-PHY-S3-Q07, source={Score Advanced - Test 06 (04 Sep 2021)}, paper=2, qno=7, type=numerical, subtopic={Refraction}, answer={2}, answer_src=solved, working={x+n sqrt((f-x)\^{}2+y\^{}2)=nf gives (n\^{}2-1)x\^{}2+n\^{}2y\^{}2-2n(n-1)fx=0, B=2}, difficulty=8, diff_parts={K3 S1 I2 C2 F1}, diff_note={concepts: surface that focuses a parallel beam, refraction, equal optical path for all rays | steps: equal optical path for all rays}]
Section of a refracting surface is shown in figure. The section is symmetrical about x axis. Refractive index of medium on right of the surface is n with respect to the medium on left of the surface. Parallel monochromatic light rays are incident on the surface as shown in figure and the refracted rays are focused at F at a distance f from origin O of the co-ordinate system shown in figure. The equation y(x) of shape of the section shown is found to be (n$^2$ -- 1) x$^2$ + n$^2$y$^2$ -- Bn (n--1) fx = 0. Find B.\par
\bfig{156.0}{SA06_p2_phy_s3_q07_1.png}
\end{bankq}
\begin{bankpassage}[id=SA07-P1-PHY-S2-G9-10, source={Score Advanced - Test 07 (16 Apr 2023)}, questions={9-10}]
A long rectangular slab of transparent medium of thickness d is placed on table with length parallel to x-axis and width parallel to the y-axis. A ray of light traveling nearly along y-axis at origin. The refractive index $\mu$ of the medium varies as $\mu = \sqrt{1+\text{e}^{\text{x/d}}}$. The refractive index of the air is 1.\par
\bfig{97.7}{SA07_p1_phy_s2_para9_10_1.png}
\end{bankpassage}
\begin{bankq}[id=SA07-P1-PHY-S2-Q09, source={Score Advanced - Test 07 (16 Apr 2023)}, paper=1, qno=9, type=single, subtopic={Refraction}, answer={B}, answer_src=key, difficulty=7, diff_parts={K3 S1 I2 C2 F0}, diff_note={concepts: ray in a medium of varying refractive index, refraction, Snell's law invariant mu sin(theta) in the slab | steps: Snell's law invariant mu sin(theta) in the slab; turning point at the upper surface}, group=SA07-P1-PHY-S2-G9-10]
Determine the x-coordinate of the point A, where the ray intersect the upper surface of the slab-air boundary.
\begin{qoptions}
\opt{A}{d(ln2)}
\opt{B}{2d(ln2)}
\opt{C}{$\dfrac{\text{d(ln 2)}}{2}$}
\opt{D}{None of these}
\end{qoptions}
\end{bankq}
\begin{bankq}[id=SA07-P1-PHY-S2-Q10, source={Score Advanced - Test 07 (16 Apr 2023)}, paper=1, qno=10, type=single, subtopic={Refraction}, answer={A}, answer_src=key, difficulty=5, diff_parts={K2 S1 I2 C1 F0}, diff_note={concepts: refractive index at the exit point, refraction | steps: mu\^{}2=1+e\^{}(x/d) at x=2d ln2}, group=SA07-P1-PHY-S2-G9-10]
Write down the refractive index of the medium at A.
\begin{qoptions}
\opt{A}{$\sqrt{5}$}
\opt{B}{$\sqrt{6}$}
\opt{C}{$\sqrt{7}$}
\opt{D}{$\sqrt{8}$}
\end{qoptions}
\end{bankq}
\begin{bankq}[id=SA07-P1-PHY-S3-Q06, source={Score Advanced - Test 07 (16 Apr 2023)}, paper=1, qno=6, type=numerical, subtopic={Prisms}, answer={0.75}, answer_src=key, difficulty=8, diff_parts={K3 S1 I2 C2 F1}, diff_note={concepts: total internal reflection in a right-angled prism, prisms, critical angle conditions at the two faces | steps: critical angle conditions at the two faces; angle theta=15 degrees}]
A right angled prism is placed with air all around it. A horizontal incident ray strikes the face PQ of prism normally, then undergoes total internal reflection twice as shown and exists the prism from face PR horizontally. The refractive index of prism is $\mu$ and angle $\theta$ = 15º. If the value of $\mu$ is $\sqrt{\text{n}}+\dfrac{1}{2}$ then find n.\par
\bfig{154.6}{SA07_p1_phy_s3_q06_1.png}
\end{bankq}
\begin{bankq}[id=SA07-P2-PHY-S1-Q06, source={Score Advanced - Test 07 (16 Apr 2023)}, paper=2, qno=6, type=multi, subtopic={Prisms}, answer={A}, answer_src=key, difficulty=8, diff_parts={K3 S2 I2 C1 F1}, diff_note={concepts: polarisation by reflection at a prism, prisms, Brewster angle tan(ip)=4/3 | steps: Brewster angle tan(ip)=4/3; refraction angle; deviation}]
Unpolarized light is incident on a prism of angle 45$^\circ$ and refractive index 4/3. The ray partially reflected at the first surface is completely polarized. Choose the correct statement(s).
\begin{qoptions}
\opt{A}{The angle of deviation is approximately 18.7$^\circ$.}
\opt{B}{The transmitted ray at the second face is completely polarized.}
\opt{C}{The transmitted ray at the first face is completely polarized.}
\opt{D}{There is no partial reflection at the second face.}
\end{qoptions}
\end{bankq}
\begin{bankq}[id=SA07-P2-PHY-S3-Q02, source={Score Advanced - Test 07 (16 Apr 2023)}, paper=2, qno=2, type=numerical, subtopic={Lenses}, answer={2}, answer_src=key, difficulty=8, diff_parts={K3 S1 I2 C2 F1}, diff_note={concepts: lens combination producing no deviation, lenses, focal lengths of the half lenses | steps: focal lengths of the half lenses; effect of the liquid between them}]
An equiconvex lens of focal length 20cm is placed in air medium. It is divided into two parts about an axis $\perp$ to principal axis. One part is inverted about axis of division of lens. Both parts are in contact with each other. Space between lenses is filled by a liquid of unknown refractive index. When light rays parallel to principal axis fall on this combination no deflection is produced in light rays by this combination. Find refractive index of liquid. \{Refractive index of lens is $\dfrac{3}{2}$\}\par
\bfig{169.7}{SA07_p2_phy_s3_q02_1.png}
\end{bankq}
\begin{bankq}[id=SA08-P2-PHY-S1-Q06, source={Score Advanced - Test 08 (22 Apr 2023)}, paper=2, qno=6, type=multi, subtopic={Reflection}, answer={A,C}, answer_src=key, difficulty=8, diff_parts={K3 S1 I2 C2 F1}, diff_note={concepts: maximum speed of the image of a falling particle, reflection, mirror formula | steps: mirror formula; differentiate image position}]
A particle is dropped along the axis from a height f/2 on a concave mirror of focal length f as shown in figure. The acceleration due to gravity is g. The maximum speed of image is :-\par
\bfig{105.5}{SA08_p2_phy_s1_q06_1.png}
\begin{qoptions}
\opt{A}{$\dfrac{3}{8}\sqrt{12\text{fg}}$}
\opt{B}{$\dfrac{3}{4}\sqrt{5\text{fg}}$}
\opt{C}{$\dfrac{3}{4}\sqrt{3\text{fg}}$}
\opt{D}{$\dfrac{3}{16}\sqrt{7\text{fg}}$}
\end{qoptions}
\end{bankq}
\begin{bankq}[id=SA08-P2-PHY-S3-Q05, source={Score Advanced - Test 08 (22 Apr 2023)}, paper=2, qno=5, type=numerical, subtopic={Reflection}, answer={3}, answer_src=key, difficulty=8, diff_parts={K3 S1 I2 C2 F1}, diff_note={concepts: slope of the image line for a concave mirror, reflection, mirror formula applied to two points of the object line | steps: mirror formula applied to two points of the object line}]
A concave mirror of radius of curvature 50 cm lies at origin of the coordinate system such that its axis coincides with x axis. A small straight object lies on the x-axis along the line 5y + 4x $-$ 25 = 0, where x and y are in cm. Assuming mirror formula to be valid for the image formation of this object, calculate the value of 5m where m is slope of the line along which the image lies.
\end{bankq}
\begin{bankq}[id=SA09-P1-PHY-S3-Q06, source={Score Advanced - Test 09 (29 Apr 2023)}, paper=1, qno=6, type=numerical, subtopic={Reflection}, answer={1}, answer_src=key, difficulty=8, diff_parts={K3 S1 I2 C2 F1}, diff_note={concepts: distance from a tilted mirror to see the full image, reflection, geometry of reflection in the tilted plane mirror | steps: geometry of reflection in the tilted plane mirror}]
A large plane mirror with its bottom on the floor is tilted at an angle 30\textdegree{} to the vertical. A boy whose eyes are at height $\sqrt{3}$ m above the floor is standing in front of the mirror. At what maximum distance (in m) from mirror should the boy be to see his full image in mirror?\par
\bfig{174.9}{SA09_p1_phy_s3_q06_1.png}
\end{bankq}
\begin{bankq}[id=SA09-P2-PHY-S1-Q04, source={Score Advanced - Test 09 (29 Apr 2023)}, paper=2, qno=4, type=multi, subtopic={Refraction}, answer={A,B}, answer_src=key, difficulty=8, diff_parts={K3 S1 I2 C2 F1}, diff_note={concepts: images in a cabin made of glass plates, refraction, apparent depth in glass | steps: apparent depth in glass; reflection from each plate}]
There is an insect inside a cabin eyeing towards a thick glass plate P$_1$. Insect sees the images of light source across glass plate P$_1$ outside the cabin. Cabin is made of thick glass plate of refractive index $\mu$ = $\dfrac{3}{2}$ and thickness 3 cm. Insect is eyeing from the middle of the cabin and there is an identical glass plate P$_2$ on the other side as shown in the figure. (glass plates are partially reflective and consider only paraxial rays)\par
\bfig{192.8}{SA09_p2_phy_s1_q04_1.png}
\begin{qoptions}
\opt{A}{The first image is seen at a distance of 14 cm from eye of insect.}
\opt{B}{The second image is seen at a distance of 18 cm from eye of insect.}
\opt{C}{Number of images seen by insect is 2.}
\opt{D}{Number of images seen by insect is 4.}
\end{qoptions}
\end{bankq}
\begin{bankq}[id=SA10-P1-PHY-S1-Q05, source={Score Advanced - Test 10 (06 May 2023)}, paper=1, qno=5, type=multi, subtopic={Prisms}, answer={C,D}, answer_src=key, difficulty=8, diff_parts={K3 S1 I2 C2 F1}, diff_note={concepts: deviation from vernier readings of a spectrometer, prisms, least count of the vernier | steps: least count of the vernier; difference of telescope positions}]
The common prism spectroscope is shown in figure is generally made up of three parts, viz., a prism and two telescopes. One of these telescopes is used to render the rays which fall on the prism parallel to each other, and is therefore called the collimator, the other telescope is made movable so that it can be turned into the proper direction for observing any desired color, and is, therefore, called the view-telescope. A circular scale is also provided to measure angular position of telescope. On the circular scale 29 divisions of the main scale exactly coincide with the 30 divisions of the vernier scale. The smallest division of the main scale is half-a-degree (0.5\textdegree). Initially view-telescope is focused to collect rays directly from collimator and position of telescope is measured. This is position (1). Now a prism is placed as shown in the figure and telescope is focused for red, yellow and violet. The respective positions of telescope are (2), (3) and (4).\par
{\footnotesize\setlength{\tabcolsep}{2pt}\noindent\begin{tabular}{|l|c|c|}
\hline
 & Reading of main scale & Reading of veriner scale\\ \hline
Position (1) & 20.0\textdegree & 20\\ \hline
Position (2) & 30.0\textdegree & 16\\ \hline
Position (3) & 32.5\textdegree & 9\\ \hline
Position (4) & 33.5\textdegree & 11\\ \hline
\end{tabular}}\par
\bfig{243.0}{SA10_p1_phy_s1_q05_1.png}
The deviation for:
\begin{qoptions}
\opt{A}{red light 10\textdegree 4'}
\opt{B}{violet light 13\textdegree 19'}
\opt{C}{yellow light 12\textdegree 19'}
\opt{D}{violet light is 13\textdegree 21'}
\end{qoptions}
\end{bankq}
\begin{bankq}[id=SA10-P1-PHY-S3-Q04, source={Score Advanced - Test 10 (06 May 2023)}, paper=1, qno=4, type=numerical, subtopic={Refraction}, answer={1.33}, answer_src=key, difficulty=8, diff_parts={K3 S1 I2 C2 F1}, diff_note={concepts: focused position of a parallel beam through a glass hemisphere, refraction, refraction at the plane surface | steps: refraction at the plane surface; then at the curved surface for the paraxial rays}]
Figure shows a solid hemispherical glass of refractive index $\mu$ = $\dfrac{3}{2}$ and radius R = $\dfrac{30}{\pi}$cm. A parallel beam of light is incident on the plane surface at angle of incidence 6\textdegree. At what distance (in cm) from principal axis beam will get focused after refraction from both surfaces?\par
\bfig{119.0}{SA10_p1_phy_s3_q04_1.png}
\end{bankq}
\begin{bankq}[id=SA11-P1-PHY-S3-Q06, source={Score Advanced - Test 11 (09 May 2023)}, paper=1, qno=6, type=numerical, subtopic={Refraction}, answer={6.055}, answer_src=key, difficulty=8, diff_parts={K3 S1 I2 C2 F1}, diff_note={concepts: ray through a block with a spherical surface, refraction, Snell's law at the plane surface | steps: Snell's law at the plane surface; geometry at the curved face}, also_in={Score Advanced - Test 29 (01 May 2024) P1 Q6}]
Figure shows an irregular block of material of refractive index $\sqrt{2}$. A ray of light strikes the face AB as shown in the figure. After refraction it is incident on a spherical surface CD of radius of curvature 0.4m and enters a medium of refractive index 1.514 to meet PQ at E. Find the distance OE in meters (upto two places of decimal)\par
\bfig{112.8}{SA11_p1_phy_s3_q06_1.png}
\end{bankq}
\begin{bankq}[id=SA12-P1-PHY-S2-Q09, source={Score Advanced - Test 12 (11 May 2023)}, paper=1, qno=9, type=multi, subtopic={Refraction}, answer={A,B,C}, answer_src=key, difficulty=8, diff_parts={K3 S1 I2 C2 F1}, diff_note={concepts: point source inside a glass block and its images, refraction, refraction at the plane and spherical surfaces | steps: refraction at the plane and spherical surfaces}, also_in={Score Advanced - Test 29 (01 May 2024) P2 Q9}]
As shown below, a point light source (S) is inside a glass block of refractive index n$_\text{G}$ = $\dfrac{3}{2}$, at a distance 2 mm from the optical axis and distance x = 3 cm from the flat surface. A thick plano-convex lens made of transparent plastic material of refractive index n$_\text{P}$ = 2, is placed near the source as shown. The thickness of the lens is d, and the radius of convex surface is R = 15 cm. The width of air gap between the glass block and the lens is `t'. The width of air gap t and thickness of lens d is adjusted so that light (paraxial rays) from point source forms always a parallel beam after passing through the lens. Select \textbf{CORRET} option(s).\par
\bfig{133.8}{SA12_p1_phy_s2_q09_1.png}
\begin{qoptions}
\opt{A}{For d = 16cm, width of air gap t = 5 cm}
\opt{B}{For thickness d = 10 cm, width of air gap t = 8 cm}
\opt{C}{The maximum adjustable value of d = 26 cm.}
\opt{D}{The maximum adjustable value of d = 16 cm.}
\end{qoptions}
\end{bankq}
\begin{bankq}[id=SA13-P1-PHY-S2-Q07, source={Score Advanced - Test 13 (13 May 2023)}, paper=1, qno=7, type=numerical, subtopic={Lenses}, answer={3}, answer_src=key, difficulty=8, diff_parts={K3 S1 I2 C2 F1}, diff_note={concepts: plano-convex lens on water in a large concave mirror, lenses, focal length of the combination | steps: focal length of the combination; image at the object position}]
A large concave mirror with a radius of curvature of 0.5 m is placed on a horizontal plate. We pour a small amount of water into the mirror and place on it a plano-convex lens with a radius of curvature of 1.0 m whose plane side just covers the water. Then we place a point light source above this setup at a height of 0.3 m. The distance (in meters) from the object where the final image is formed will be $\dfrac{\text{x}}{90}$. Find x. (Given that the refractive index of water is $\dfrac{4}{3}$ and the refractive index of glass is 1.50)\par
\end{bankq}
\begin{bankpassage}[id=SA13-P2-PHY-S2-G7-8, source={Score Advanced - Test 13 (13 May 2023)}, questions={7-8}]
A rectangular wire frame is kept on the principal axis of a concave mirror as shown in figure. F \& C are the focus \& centre of curvature of mirror respectively.\par
\bfig{108.3}{SA13_p2_phy_s2_para7_8_1.png}
f is the focal length of mirror. A small particle is moving along the length of wire frame with constant velocity. Based on above information answer the following questions.\par
\end{bankpassage}
\begin{bankq}[id=SA13-P2-PHY-S2-Q07, source={Score Advanced - Test 13 (13 May 2023)}, paper=2, qno=7, type=single, subtopic={Lenses}, answer={A}, answer_src=key, difficulty=7, diff_parts={K3 S1 I2 C2 F0}, diff_note={concepts: image of a wire frame by a concave mirror, lenses, mirror formula for different positions of the sides | steps: mirror formula for different positions of the sides}, group=SA13-P2-PHY-S2-G7-8]
The image of the given object is best shown by :-
\begin{qoptions}
\opt{A}{\bfigi{108.2}{SA13_p2_phy_s2_q07_A.png}}
\opt{B}{\bfigi{108.3}{SA13_p2_phy_s2_q07_B.png}}
\opt{C}{\bfigi{108.2}{SA13_p2_phy_s2_q07_C.png}}
\opt{D}{\bfigi{108.3}{SA13_p2_phy_s2_q07_D.png}}
\end{qoptions}
\end{bankq}
\begin{bankq}[id=SA13-P2-PHY-S2-Q08, source={Score Advanced - Test 13 (13 May 2023)}, paper=2, qno=8, type=single, subtopic={Lenses}, answer={D}, answer_src=key, difficulty=7, diff_parts={K3 S1 I2 C2 F0}, diff_note={concepts: statements on a concave mirror image, reflection, check each statement | steps: check each statement}, group=SA13-P2-PHY-S2-G7-8]
Mark the correct statement(s) :-
\begin{qoptions}
\opt{A}{The velocity of image of particle is remaining constant.}
\opt{B}{The velocity of image of particle is varying.}
\opt{C}{The total length of image of wire frame is $\left(\dfrac{7+\sqrt{17}}{4}\right)$f.}
\opt{D}{Both (B) and (C) are correct.}
\end{qoptions}
\end{bankq}
\begin{bankq}[id=SA14-P1-PHY-S3-Q05, source={Score Advanced - Test 14 (16 May 2023)}, paper=1, qno=5, type=numerical, subtopic={Prisms}, answer={2.12}, answer_src=key, difficulty=8, diff_parts={K3 S1 I2 C2 F1}, diff_note={concepts: prism made of two halves with different refractive indices, prisms, deviation through each half | steps: deviation through each half}, also_in={Score Advanced - Test 20 (06 Apr 2024) P1 Q5}]
For a prism of prism angle $\theta$ = 60$^\circ$, the refractive indices of the left half and the right half are, respectively, n$_1$ and n$_2$(n$_2\geq$ n$_1$) as shown in the figure. The angle of incidence `i' is chosen such that the incident light rays will have minimum deviation if n$_1$ = n$_2$ = n = 1.5. For the case of unequal refractive indices, n$_1$ = n and n$_2$ = n + $\Delta$n (where $\Delta$n $<<$ n), the angle of emergence e = i + $\Delta$e. If $\Delta$n = 2.8 x 10$^{-3}$ rad then the value of $\Delta$e in milliradians is equal to (Round off to nearest decimal up to 2 digits)\par
\bfig{106.6}{SA14_p1_phy_s3_q05_1.png}
\end{bankq}
\begin{bankq}[id=SA15-P1-PHY-S1-Q04, source={Score Advanced - Test 15 (18 May 2023)}, paper=1, qno=4, type=single, subtopic={Lenses}, answer={C}, answer_src=key, difficulty=7, diff_parts={K3 S1 I2 C2 F0}, diff_note={concepts: velocity of the image of an observer in a moving mirror, lenses, lens formula applied twice with the moving plane mirror | steps: lens formula applied twice with the moving plane mirror}]
An observer at O is at 100 cm from a convex lens of focal length 50 cm. A plane mirror is placed 150 cm behind the lens and starts moving towards the right at 10 cm/s. Find the velocity of her own image as seen by the observer when the mirror just starts moving.\par
\bfig{99.7}{SA15_p1_phy_s1_q04_1.png}
\begin{qoptions}
\opt{A}{$\dfrac{200}{9}$ cm/s}
\opt{B}{$\dfrac{200}{3}$ cm/s}
\opt{C}{$\dfrac{20}{9}$ cm/s}
\opt{D}{$\dfrac{10}{9}$ cm/s}
\end{qoptions}
\end{bankq}
\begin{bankq}[id=SA16-P1-PHY-S1-Q04, source={Score Advanced - Test 16 (20 May 2023)}, paper=1, qno=4, type=single, subtopic={Refraction}, answer={B}, answer_src=key, difficulty=7, diff_parts={K3 S1 I2 C2 F0}, diff_note={concepts: ray through a transparent ball intersecting the axis, refraction, refraction at two surfaces of the sphere | steps: refraction at two surfaces of the sphere; small geometry relation between alpha and beta}]
Rays of light emitted by a point source S are incident on a uniform ball of transparent material with a refractive index $n$. The ray leaving the source S at an angle $\alpha$ to the straight line MN, on which the source and the center of the ball lie, after two refractions at the boundary of the ball, intersects MN at an angle $\beta$ at point K (see figure). The distance SK is equal to $\ell$. The value of radius $R$ of the ball is given by :- (Take : $\alpha=53^\circ$, $\beta=37^\circ$ and $n=\sqrt2$)\par
\bfig{177.0}{SA16_p1_phy_s1_q04_B.png}
\begin{qoptions}
\opt{A}{$\dfrac{3\sqrt2\,\ell}{35}$}
\opt{B}{$\dfrac{12\ell}{35}$}
\opt{C}{$\dfrac{6\sqrt2\,\ell}{175}$}
\opt{D}{$\dfrac{\ell}{5}$}
\end{qoptions}
\end{bankq}
\begin{bankpassage}[id=SA16-P1-PHY-S3-G1-2, source={Score Advanced - Test 16 (20 May 2023)}, questions={1-2}]
\textbf{Question Stem}\par
An experiment is performed on thin concavo-convex lens in order to determine the radii $R_1$ (convex) and $R_2$ (concave) of curvature of its two surfaces and the refractive index `$\mu$' of its glass.\par
\textbf{Experiment I:} A real object is placed 36 cm to the left of lens and a real image is formed 72 cm to the right of lens.\par
\textbf{Experiment II:} The lens is floated on the surface of mercury with surface having radius of curvature $R_2$ facing air. A real object is placed at 60 cm above it, and a real image is formed at 30 cm above it (assume mercury to be perfectly reflecting).\par
\textbf{Experiment III:} The lens is floated on the surface of mercury with surface having radius of curvature $R_1$ facing air. A real object is placed at 7.5 cm above it and real image is formed at $\dfrac{20}{3}$ cm above it. On the bases of Experiments I, II and III, answer the following questions.\par
\end{bankpassage}
\begin{bankq}[id=SA16-P1-PHY-S3-Q01, source={Score Advanced - Test 16 (20 May 2023)}, paper=1, qno=1, type=numerical, subtopic={Lenses}, answer={1.50}, answer_src=key, difficulty=8, diff_parts={K3 S1 I2 C2 F1}, diff_note={concepts: refractive index of a concavo-convex lens, lenses, three experiments give two equations in the focal powers | steps: three experiments give two equations in the focal powers; solve for mu}, group=SA16-P1-PHY-S3-G1-2]
Find the refractive index of material of lens.
\end{bankq}
\begin{bankq}[id=SA16-P1-PHY-S3-Q02, source={Score Advanced - Test 16 (20 May 2023)}, paper=1, qno=2, type=numerical, subtopic={Lenses}, answer={60}, answer_src=key, difficulty=8, diff_parts={K3 S1 I2 C2 F1}, diff_note={concepts: radius of curvature of the convex surface, lenses, lens plus mercury mirror combination | steps: lens plus mercury mirror combination; lens formula}, group=SA16-P1-PHY-S3-G1-2]
Find the magnitude of $R_1$ (in cm)
\end{bankq}
\begin{bankq}[id=SA16-P2-PHY-S1-Q04, source={Score Advanced - Test 16 (20 May 2023)}, paper=2, qno=4, type=multi, subtopic={Reflection}, answer={A,B,C}, answer_src=key, difficulty=8, diff_parts={K3 S1 I2 C2 F1}, diff_note={concepts: image coinciding with the object between two mirrors, reflection, multiple reflections | steps: multiple reflections; mirror formula}]
A small object O is placed symmetrically between two small concave mirrors of same radius of curvature $R$. Find the value(s) of the distance $d$ between the mirrors for which an image formed between the mirrors coincides with O.\par
\bfig{71.1}{SA16_p2_phy_s1_q04_1.png}
\begin{qoptions}
\opt{A}{$R$}
\opt{B}{1.5$R$}
\opt{C}{2$R$}
\opt{D}{3$R$}
\end{qoptions}
\end{bankq}
\begin{bankq}[id=SA16-P2-PHY-S4-Q09, source={Score Advanced - Test 16 (20 May 2023)}, paper=2, qno=9, type=numerical, subtopic={Lenses}, answer={35}, answer_src=key, difficulty=8, diff_parts={K3 S1 I2 C2 F1}, diff_note={concepts: ray through a chain of lenses at the vertices of a polygon, lens formula repeatedly, exterior angle symmetry | steps: lens formula repeatedly; exterior angle symmetry}]
At the vertices of the regular 20-gon, there are 20 identical lenses. The optical centers of the lenses are located exactly at the vertices of the polygon, the planes of all lenses are perpendicular to one of the sides adjacent to the lens. The focal lengths of the lenses are all equal to $F=10$ cm and coincide with the length of the side of the 20-gon. One of the lenses is illuminated by a parallel light flux directed along its optical axis. It turns out that one of the rays has a closed trajectory. Determine the radius of the circle (in cm) inscribed in this trajectory. Consider all of the lenses are converging and all angles small such that $\sin\alpha\approx\tan\alpha\approx\alpha$.\par
\bfig{87.5}{SA16_p2_phy_s4_q09_1.png}
\end{bankq}
\begin{bankq}[id=SA17-P1-PHY-S3-Q03, source={Score Advanced - Test 17 (25 May 2023)}, paper=1, qno=3, type=numerical, subtopic={Prisms}, answer={1.51}, answer_src=key, difficulty=7, diff_parts={K3 S1 I2 C1 F1}, diff_note={concepts: minimum refractive index for total internal reflection at the second face, prisms, r1+r2=A | steps: r1+r2=A; critical angle condition}, also_in={Score Advanced - Test 35 (15 May 2024) P1 Q3}]
Light falls on a prism of angle $\cos^{-1}\left(\dfrac18\right)$. If the minimum value of refractive index of the prism-material, so that light entering one face may not exit through the other face
\end{bankq}
\begin{bankq}[id=SA17-P2-PHY-S3-Q02, source={Score Advanced - Test 17 (25 May 2023)}, paper=2, qno=2, type=numerical, subtopic={Prisms}, answer={2.00}, answer_src=key, difficulty=7, diff_parts={K3 S1 I2 C1 F1}, diff_note={concepts: refractive index from grazing incidence and emergence, prisms, r1=r2=30 | steps: r1=r2=30; mu=sin 90/sin 30}, also_in={Score Advanced - Test 35 (15 May 2024) P2 Q2}]
For an equilateral prism, it is observed that when a ray strikes grazingly at one face it emerges grazingly at the other. Find the refractive index of the prism.
\end{bankq}
\begin{bankq}[id=SA18-P1-PHY-S2-Q08, source={Score Advanced - Test 18 (27 May 2023)}, paper=1, qno=8, type=single, subtopic={Lenses}, answer={A}, answer_src=key, difficulty=7, diff_parts={K3 S1 I2 C2 F0}, diff_note={concepts: focal lengths of lens and mirror combinations, lenses, lens maker formula | steps: lens maker formula; combination formula}]
Four combination of two thin lenses/lens and mirror are given in list-I. The radius of curvature of all curved surfaces is same and the refractive index of all the lenses is 1.5. Focal length of a equiconvex lens kept in air with same radius and same refractive index as that of surface here is $f$. Match lens combination in list-I with their modulus of equivalent focal length in air in list-II and select the correct answer using the code given below the lists.\par
\bfig{140.4}{SA18_p1_phy_s2_q08_1.png}
\begin{qoptions}
\opt{A}{I $\rightarrow$ R;II $\rightarrow$ Q;III $\rightarrow$ S;IV $\rightarrow$ P}
\opt{B}{I $\rightarrow$ R;II $\rightarrow$ Q;III $\rightarrow$ R;IV $\rightarrow$ P}
\opt{C}{I $\rightarrow$ S;II $\rightarrow$ Q;III $\rightarrow$ R;IV $\rightarrow$ S}
\opt{D}{I $\rightarrow$ Q;II $\rightarrow$ S;III $\rightarrow$ P;IV $\rightarrow$ R}
\end{qoptions}
\end{bankq}
\begin{bankq}[id=SA19-P1-PHY-S1-Q01, source={Score Advanced - Test 19 (04 Apr 2024)}, paper=1, qno=1, type=single, subtopic={Reflection}, answer={D}, answer_src=key, difficulty=7, diff_parts={K3 S1 I2 C2 F0}, diff_note={concepts: ratio of intensities from multiple reflections between two glass plates, reflection and transmission at each surface, sum the geometric series | steps: reflection and transmission at each surface; sum the geometric series}]
A narrow monochromatic beam of light of intensity $I$ is incident on a glass plate labeled 1 as shown in figure (not drawn to scale). Another identical glass plate labeled 2 is kept close to the first one and parallel to it. Each glass plate reflects 25\% of the light incident on it and transmits the remaining. Find the ratio $\sqrt{\dfrac{I_{max}}{I_{min}}}$ of the interference pattern formed by two beams obtained after one reflection at each plate. ($I_{max}$ and $I_{min}$ represent the maximum value and minimum value of intensity respectively. Ignore the distance between plates 1 and 2 in the diagram. It is drawn for the sake of readability. Consider only the first reflection from the glass plates)\par
\bfig{74.9}{SA19_p1_phy_s1_q01_1.png}
\begin{qoptions}
\opt{A}{3}
\opt{B}{5}
\opt{C}{6}
\opt{D}{7}
\end{qoptions}
\end{bankq}
\begin{bankq}[id=SA21-P1-PHY-S1-Q04, source={Score Advanced - Test 21 (08 Apr 2024)}, paper=1, qno=4, type=multi, subtopic={Prisms}, answer={A,B}, answer_src=key, difficulty=8, diff_parts={K3 S1 I2 C2 F1}, diff_note={concepts: identical prisms in contact, deviation by successive prisms | steps: deviation by successive prisms}]
$n$ number of identical equilateral prism are kept in contact as shown in figure. If deviation through a single prism is $\delta$. Then ($n$, $m$ are integers)(Assume surrounding medium be air)\par
\bfig{101.1}{SA21_p1_phy_s1_q04_1.png}
\begin{qoptions}
\opt{A}{If $n=2m$, deviation through $n$ prism is zero.}
\opt{B}{If $n=2m+1$, deviation through system of $n$ prisms is $\delta$.}
\opt{C}{If $n=2m$, deviation through system of $n$ prisms is $\delta$.}
\opt{D}{If $n=2m+1$, deviation through system of $n$ prism is zero.}
\end{qoptions}
\end{bankq}
\begin{bankq}[id=SA23-P1-PHY-S1-Q03, source={Score Advanced - Test 23 (14 Apr 2024)}, paper=1, qno=3, type=single, subtopic={Lenses}, answer={B}, answer_src=key, difficulty=7, diff_parts={K3 S1 I2 C2 F0}, diff_note={concepts: images of objects formed by a converging lens, lenses, image orientation and magnification for each object position | steps: image orientation and magnification for each object position}]
In the figures below, images of the two objects on the left, created by a converging lens, are given on the right. Assuming that the image given as K is correct, which of the following options is correct?\par
\bfig{213.9}{SA23_p1_phy_s1_q03_1.png}
\begin{qoptions}
\opt{A}{L, M and N}
\opt{B}{L and M}
\opt{C}{L and O}
\opt{D}{M, N and O}
\end{qoptions}
\end{bankq}
\begin{bankq}[id=SA23-P1-PHY-S4-Q07, source={Score Advanced - Test 23 (14 Apr 2024)}, paper=1, qno=7, type=numerical, subtopic={Prisms}, answer={3}, answer_src=key, difficulty=7, diff_parts={K3 S1 I2 C1 F1}, diff_note={concepts: rotation of a mirror after a thin prism deviation, prisms, thin prism deviation (n-1)A=6 degrees | steps: thin prism deviation (n-1)A=6 degrees; mirror rotation is half}]
A horizontal ray of light passes through a prism of refractive index 2.0 and apex angle 6$^\circ$ and then strikes a vertical mirror as shown in the figure. Through what angle (in degree) must the mirror be rotated if after reflection the ray is to be horizontal?\par
\bfig{113.4}{SA23_p1_phy_s4_q07_1.png}
\end{bankq}
\begin{bankq}[id=SA23-P2-PHY-S4-Q07, source={Score Advanced - Test 23 (14 Apr 2024)}, paper=2, qno=7, type=numerical, subtopic={Reflection}, answer={3}, answer_src=key, difficulty=8, diff_parts={K3 S1 I2 C2 F1}, diff_note={concepts: velocity of the image in a concave mirror on a moving platform, reflection, mirror formula derivative | steps: mirror formula derivative; platform recoil by momentum conservation}]
A man of mass m starts moving w.r.t a platform of mass 2 m with a velocity u = 9/13 m/s towards right as shown in the figure. The platform is fitted with a concave mirror of focal length f. The velocity of image (in m/s) at the moment is :-\par
\bfig{116.8}{SA23_p2_phy_s4_q07_1.png}
\end{bankq}
\begin{bankq}[id=SA23-P2-PHY-S4-Q08, source={Score Advanced - Test 23 (14 Apr 2024)}, paper=2, qno=8, type=numerical, subtopic={Refraction}, answer={3}, answer_src=key, difficulty=8, diff_parts={K3 S1 I2 C2 F1}, diff_note={concepts: apparent depth of a stone for an oblique line of sight, refraction, apparent depth for non paraxial viewing | steps: apparent depth for non paraxial viewing}]
A person is observing a stone at the bottom of the pool. The water depth is h = 14 m. The angle between the line of sight and the normal line of the water surface is 60$^\circ$, the refractive index of water is 1.33. The depth of the image of the stone (in m) is :-
\end{bankq}
\begin{bankq}[id=SA24-P1-PHY-S3-Q05, source={Score Advanced - Test 24 (17 Apr 2024)}, paper=1, qno=5, type=numerical, subtopic={Lenses}, answer={4}, answer_src=key, difficulty=9, diff_parts={K3 S2 I2 C2 F1}, diff_note={concepts: final image after two refractions and one reflection in a lens, lenses, refraction at the first surface | steps: refraction at the first surface; reflection at the second; refraction again}]
A set of parallel rays incident on the bi-convex lens. If after two refractions and one reflection final image formed at 15 cm left of lens, the value of $\dfrac X9$ is\par
\bfig{147.4}{SA24_p1_phy_s3_q05_1.png}
\end{bankq}
\begin{bankq}[id=SA24-P2-PHY-S1-Q01, source={Score Advanced - Test 24 (17 Apr 2024)}, paper=2, qno=1, type=single, subtopic={Refraction}, answer={D}, answer_src=key, difficulty=7, diff_parts={K3 S1 I2 C2 F0}, diff_note={concepts: light ray in a medium with a refractive index varying with y, refraction, Snell law in the form n(y) sin theta=constant | steps: Snell law in the form n(y) sin theta=constant}]
The refractive index of the medium within a certain region, x $>$ 0, y $>$ 0, changes with y. A thin light ray travelling in the x-direction strikes the medium at right angles and moves through the medium along a circular arc. How does the refractive index depend on y ?\par
[R : Radius of path \& $\eta_0$ : Initial refractive index]\par
\bfig{115.1}{SA24_p2_phy_s1_q01_1.png}
\begin{qoptions}
\opt{A}{$\dfrac{2R}{R-y}\eta_0$}
\opt{B}{$\dfrac{R}{R-y}\dfrac{\eta_0}{2}$}
\opt{C}{$\dfrac{R+y}{R}\eta_0$}
\opt{D}{$\dfrac{R}{R-y}\eta_0$}
\end{qoptions}
\end{bankq}
\begin{bankq}[id=SA25-P1-PHY-S2-Q06, source={Score Advanced - Test 25 (21 Apr 2024)}, paper=1, qno=6, type=multi, subtopic={Reflection}, answer={B,C}, answer_src=key, difficulty=8, diff_parts={K3 S1 I2 C2 F1}, diff_note={concepts: images of a lamp in mirrors on three walls of a hexagonal room, reflection, image positions by repeated reflection | steps: image positions by repeated reflection}, also_in={Score Advanced - Test 26 (21 Apr 2024) P1 Q6}]
Three adjacent walls of a regular hexagon-shaped room are covered with flat mirror (see the figure). A lamp is lit at the centre of the room.\par
\bfig{59.3}{SA25_p1_phy_s2_q06_1.png}
\begin{qoptions}
\opt{A}{The number of images seen by an observer at A is 4.}
\opt{B}{The number of images seen by an observer at B is 4.}
\opt{C}{The number of images seen by an observer at C is 5.}
\opt{D}{The number of images seen by an observer at D is 5.}
\end{qoptions}
\end{bankq}
\begin{bankq}[id=SA25-P1-PHY-S4-Q05, source={Score Advanced - Test 25 (21 Apr 2024)}, paper=1, qno=5, type=numerical, subtopic={Lenses}, answer={1}, answer_src=key, difficulty=8, diff_parts={K3 S1 I2 C2 F1}, diff_note={concepts: two lens system that forms an image of the same size, lenses, lens formula with magnification condition | steps: lens formula with magnification condition}, also_in={Score Advanced - Test 26 (21 Apr 2024) P1 Q5}]
Consider two convex lens with focal length $f_a$ and $f_b$ respectively. If placed appropriately, a small object perpendicular to the optic axis can create an inverted image, the size and position of which matches the size and position of object as shown in the figure. The distance between the lens is d and distance between the object and the left lens is u. The ratio $\dfrac ud$ is (Take : $f_a$ = 1 cm, $f_b$ = 4 cm)\par
\bfig{135.7}{SA25_p1_phy_s4_q05_1.png}
\end{bankq}
\begin{bankq}[id=SA27-P1-PHY-S1-Q02, source={Score Advanced - Test 27 (24 Apr 2024)}, paper=1, qno=2, type=single, subtopic={Prisms}, answer={C}, answer_src=key, difficulty=7, diff_parts={K3 S1 I2 C2 F0}, diff_note={concepts: ray through two right angled prisms, deviation, refraction at each face | steps: deviation and refraction at each face}]
Two right angled prisms having the same refracting angle 60$^\circ$ are placed as shown in the figure. A ray of light incident on the first prism is finally deviated by an angles 90$^\circ$ grazingly as shown. The refractive index of second prism is $\sqrt2$. The refractive index of the first prism is equal to\par
\bfig{108.2}{SA27_p1_phy_s1_q02_1.png}
\begin{qoptions}
\opt{A}{$\sqrt2\,\dfrac{\sin30^\circ}{\sin60^\circ}$}
\opt{B}{$\sqrt2\,\dfrac{\sin45^\circ}{\sin60^\circ}$}
\opt{C}{$\sqrt2\,\dfrac{\sin15^\circ}{\sin60^\circ}$}
\opt{D}{$\sqrt2\,\dfrac{\sin60^\circ}{\sin30^\circ}$}
\end{qoptions}
\end{bankq}
\begin{bankpassage}[id=SA27-P2-PHY-S2-G11-12, source={Score Advanced - Test 27 (24 Apr 2024)}, questions={11-12}]
An object O (real) is placed at focus of an equi-biconvex lens as shown in figure. The refractive index of the lens is $\mu$ = 1.5 and the radius of curvature of either surface of lens is R. The lens is surrounded by air. In each statement of List-I, some changes are made to situation given above and information regarding final image formed as a result is given in List-II. The distance between lens and object is unchanged in all statements of List-I. Match the statements in List-I with resulting image in List-II.\par
\bfig{93.0}{SA27_p2_phy_s2_para11_12_1.png}
{\footnotesize\setlength{\tabcolsep}{1.5pt}\noindent\begin{tabular}{|c|L{3.1cm}|c|L{3.7cm}|}
\hline
 & \multicolumn{1}{c|}{List-I} & & \multicolumn{1}{c|}{List-II} \\ \hline
(I) & If the refractive index of the lens is doubled (that is, made 2$\mu$), then & (P) & Final image is real \\ \hline
(II) & If the radius of curvature is doubled (that is, made 2R), then & (Q) & Final image is virtual \\ \hline
(III) & If a glass slab of refractive index $\mu$ = 1.5 is introduced between the object and lens as shown, then \bfigi{59.0}{SA27_p2_phy_s2_para11_12_III.png} & (R) & Final image is always enlarged. \\ \hline
(IV) & If the left side of lens is filled with a medium of refractive index $\mu$ = 1.5 as shown, then \bfigi{62.1}{SA27_p2_phy_s2_para11_12_IV.png} & (S) & Final image is of same size as the object \\ \hline
 & & (T) & Final image becomes smaller in size in comparison to size of image before the change was made \\ \hline
 & & (U) & Image formation is not possible. \\ \hline
\end{tabular}}\par
\end{bankpassage}
\begin{bankq}[id=SA27-P2-PHY-S2-Q11, source={Score Advanced - Test 27 (24 Apr 2024)}, paper=2, qno=11, type=single, subtopic={Lenses}, answer={C}, answer_src=solved, working={With f = R and the object at R: (I) doubling mu gives f=R/4, so a real image at R/3 (smaller than before: image earlier at infinity), P and T. (II) 2R gives f=2R, object inside focus, virtual enlarged image: Q and T. Option C.}, difficulty=7, diff_parts={K3 S1 I2 C2 F0}, diff_note={concepts: image changes of a lens when index or radius or medium is changed, lenses, object at the focal distance R | steps: object at the focal distance R; lens formula for each change}, group=SA27-P2-PHY-S2-G11-12]
Choose correct options for conditions in I and II
\begin{qoptions}
\opt{A}{I-P; II-P,Q,R}
\opt{B}{I-P,T; II-P,Q,T}
\opt{C}{I-P,T; II- Q,T}
\opt{D}{I-P; II-P,Q,S}
\end{qoptions}
\end{bankq}
\begin{bankq}[id=SA27-P2-PHY-S2-Q12, source={Score Advanced - Test 27 (24 Apr 2024)}, paper=2, qno=12, type=single, subtopic={Lenses}, answer={A}, answer_src=solved, working={(III) The slab brings the object nearer by t/3 so it lies inside the focus: virtual image, smaller than the earlier image at infinity: Q,T. (IV) Only the right surface refracts: 1/v - 1.5/(-R) = -0.5/(-R) gives v=-R, virtual and smaller than before: Q,T. Option A.}, difficulty=7, diff_parts={K3 S1 I2 C2 F0}, diff_note={concepts: image changes with a slab or medium before the lens, lenses, apparent shift of the object by a slab | steps: apparent shift of the object by a slab; refraction at the second surface only}, group=SA27-P2-PHY-S2-G11-12]
Choose correct options for conditions in III and IV
\begin{qoptions}
\opt{A}{III-Q,T; IV-Q,T}
\opt{B}{III-P,T; IV-Q,T}
\opt{C}{III-P; IV- P,Q,R}
\opt{D}{III-Q,T; IV-P,Q,T}
\end{qoptions}
\end{bankq}
\begin{bankq}[id=SA28-P1-PHY-S1-Q03, source={Score Advanced - Test 28 (27 Apr 2024)}, paper=1, qno=3, type=single, subtopic={Refraction}, answer={0.06}, answer_src=key, difficulty=7, diff_parts={K3 S1 I2 C2 F0}, diff_note={concepts: parallel beam on a quarter cylinder, refraction, refraction at the curved surface and the critical angle | steps: refraction at the curved surface and the critical angle}]
A parallel beam of light falls on a quarter cylinder of radius 9 cm as shown in figure (A). Refractive index of the material of the cylinder is $\sqrt3$. Maximum value of OP in meter shown in figure (B), so that rays don't suffer TIR at the curved surface is (consider refractive index of surrounding to be 1).\par
\bfig{108.0}{SA28_p1_phy_s1_q03_1.png}
\end{bankq}
\begin{bankq}[id=SA28-P2-PHY-S3-Q02, source={Score Advanced - Test 28 (27 Apr 2024)}, paper=2, qno=2, type=matching, subtopic={Prisms}, answer={P$\to$1; Q$\to$3; R$\to$2; S$\to$4}, answer_src=key, difficulty=8, diff_parts={K3 S1 I2 C2 F1}, diff_note={concepts: colours of rays emerging from a prism with total internal reflection, prisms, dispersion of white light | steps: dispersion of white light}]
A white light ray is incident on a glass prism, and it create four refracted rays A, B, C and D. Match the refracted rays with the colours given (1 \& D are rays due to total internal reflection):\par
\bfig{90.1}{SA28_p2_phy_s3_q02_1.png}
{\footnotesize\setlength{\tabcolsep}{1.5pt}\noindent\begin{tabular}{|c|L{2.2cm}|c|L{2.2cm}|}
\hline
 & \multicolumn{1}{c|}{List-I (Ray)} & & \multicolumn{1}{c|}{List-II (Colour)} \\ \hline
(P) & A & (1) & Red \\ \hline
(Q) & B & (2) & Green \\ \hline
(R) & C & (3) & Yellow \\ \hline
(S) & D & (4) & Blue \\ \hline
\end{tabular}}\par
\end{bankq}
\begin{bankq}[id=SA28-P2-PHY-S3-Q06, source={Score Advanced - Test 28 (27 Apr 2024)}, paper=2, qno=6, type=matching, subtopic={Reflection}, answer={P$\to$2; Q$\to$1; R$\to$3; S$\to$3}, answer_src=key, difficulty=8, diff_parts={K3 S1 I2 C2 F1}, diff_note={concepts: speed of the image for a moving mirror, reflection, image velocity relative to mirror and ground | steps: image velocity relative to mirror and ground}]
A plane mirror is tied to the free end of an ideal spring. The other end of the spring is attached to a wall. The spring with mirror is held vertically to the floor can slide along it smoothly. When the spring is at its natural length, the mirror is found to move at speed v with respect to ground. An object moving towards mirror with speed 2v with respect to ground. Match the following : (Assume no collision before mirror comes to rest for 1$^{\mathrm{st}}$ time)\par
\bfig{133.0}{SA28_p2_phy_s3_q06_1.png}
{\footnotesize\setlength{\tabcolsep}{1.5pt}\noindent\begin{tabular}{|c|L{3.8cm}|c|L{1.2cm}|}
\hline
 & \multicolumn{1}{c|}{List-I} & & \multicolumn{1}{c|}{List-II} \\ \hline
(P) & Speed of image with respect to ground when spring at natural length & (1) & v \\ \hline
(Q) & Speed of image with respect to mirror when spring at natural length & (2) & 0 \\ \hline
(R) & Speed of image with respect to object when spring is at natural length & (3) & 2v \\ \hline
(S) & Speed of image with respect to ground frame when spring at maximum compressed state. & (4) & 3v \\ \hline
\end{tabular}}\par
\end{bankq}
\begin{bankq}[id=SA30-P1-PHY-S2-Q08, source={Score Advanced - Test 30 (04 May 2024)}, paper=1, qno=8, type=multi, subtopic={Lenses}, answer={A,D}, answer_src=key, difficulty=8, diff_parts={K3 S1 I2 C2 F1}, diff_note={concepts: standard problem in this chapter, lenses, apply the relevant results | steps: apply the relevant results; compute}, also_in={Score Advanced - Test 31 (04 May 2024) P1 Q8}]
As shown in figure, a plano-convex lens has a refractive index n and is placed in the air. The radius of the lens aperture is R. Take a point F on the principal optical axis of the lens, OF = f. When parallel light is incident along the principal optical axis, all the light rays converge at point F.\par
\bfig{126.7}{SA30_p1_phy_s2_q08_1.png}
\begin{qoptions}
\opt{A}{The shape of the convex surface of the lens should be hyperbola.}
\opt{B}{The shape of the convex surface of the lens should be parabola.}
\opt{C}{The maximum value of aperture R of the lens is nf.}
\opt{D}{The maximum value of aperture R of the lens is $\sqrt{n^2-1}\,$f.}
\end{qoptions}
\end{bankq}
\begin{bankq}[id=SA30-P1-PHY-S2-Q09, source={Score Advanced - Test 30 (04 May 2024)}, paper=1, qno=9, type=multi, subtopic={Refraction}, answer={A,D}, answer_src=key, difficulty=8, diff_parts={K3 S1 I2 C2 F1}, diff_note={concepts: standard problem in this chapter, refraction, apply the relevant results | steps: apply the relevant results; compute}]
An ice hemisphere with radius R and refractive index n lies on a warm flat table and melts slowly. The rate of heat transfer from the table to the ice is proportional to the area of contact between them. It is known that the ice hemisphere completely melts in time T$_0$. Throughout the process, a laser beam is incident on the ice from above. The beam is vertically incident at a distance of R/2 from the axis of symmetry (see figure).\par
Assume that the temperature of the ice and the surrounding atmosphere are 0$^\circ$C and remains constant during the melting process. The laser beam does not transfer energy to the ice. The melting water immediately flows off the table, and the ice does not move along the table.\par
\bfig{101.5}{SA30_p1_phy_s2_q09_1.png}
\bfig{103.2}{SA30_p1_phy_s2_q09_2.png}
\begin{qoptions}
\opt{A}{The value of x$_0$ (point on the table where the laser beam strikes at t = 0) is $\dfrac{2\mathrm R}{\left(\sqrt{12n^2-3}+1\right)}$.}
\opt{B}{The value of x$_0$ (point on the table where the laser beam strikes at t = 0) is $\dfrac{2\mathrm R}{\left(\sqrt{4n^2-3}+1\right)}$.}
\opt{C}{The height of ice z at any time t is z = R$\left(1-\dfrac{\mathrm t^2}{\mathrm T_0^2}\right)$.}
\opt{D}{The time at which the value of x (point on the table where the laser beam strikes) becomes $\dfrac{\mathrm R}2$ is $\dfrac{\sqrt3}2$T$_0$.}
\end{qoptions}
\end{bankq}
\begin{bankq}[id=SA31-P1-PHY-S2-Q09, source={Score Advanced - Test 31 (04 May 2024)}, paper=1, qno=9, type=multi, subtopic={Refraction}, answer={A,D}, answer_src=key, difficulty=8, diff_parts={K3 S1 I2 C2 F1}, diff_note={concepts: standard problem in this chapter, refraction, apply the relevant results | steps: apply the relevant results; compute}]
An ice hemisphere with radius R and refractive index n lies on a warm flat table and melts slowly. The rate of heat transfer from the table to the ice is proportional to the area of contact between them. It is known that the ice hemisphere completely melts in time T$_0$. Throughout the process, a laser beam is incident on the ice from above. The beam is vertically incident at a distance of R/2 from the axis of symmetry (see figure).\par
Assume that the temperature of the ice and the surrounding atmosphere are 0$^\circ$C and remains constant during the melting process. The laser beam does not transfer energy to the ice. The melting water immediately flows off the table, and the ice does not move along the table.\par
\bfig{227.2}{SA31_p1_phy_s2_q09_1.png}
\begin{qoptions}
\opt{A}{The value of x$_0$ (point on the table where the laser beam strikes at t = 0) is $\dfrac{2\mathrm R}{\left(\sqrt{12n^2-3}+1\right)}$.}
\opt{B}{The value of x$_0$ (point on the table where the laser beam strikes at t = 0) is $\dfrac{2\mathrm R}{\left(\sqrt{4n^2-3}+1\right)}$.}
\opt{C}{The height of ice z at any time t is z = R$\left(1-\dfrac{\mathrm t^2}{\mathrm T_0^2}\right)$.}
\opt{D}{The time at which the value of x (point on the table where the laser beam strikes) becomes $\dfrac{\mathrm R}2$ is $\dfrac{\sqrt3}2$T$_0$.}
\end{qoptions}
\end{bankq}
\begin{bankq}[id=SA32-P1-PHY-S1-Q04, source={Score Advanced - Test 32 (08 May 2024)}, paper=1, qno=4, type=single, subtopic={Lenses}, answer={C}, answer_src=key, difficulty=4, diff_parts={K2 S1 I1 C1 F0}, diff_note={concepts: standard problem in this chapter, lenses | steps: apply the relevant results; compute}]
An observer is standing at a point O, at a distance of 100 cm from a convex lens of focal length 50 cm. A plane mirror is placed behind the lens at a distance of 150 cm from the lens. The mirror now starts moving towards right with a velocity of 10 cm/s. What will be the magnitude of velocity of her own image as seen by the observer, at the moment when the mirror just starts moving ?\par
\bfig{94.5}{SA32_p1_phy_s1_q04_1.png}
\begin{qoptions}
\opt{A}{$\dfrac{200}9$ cm/s}
\opt{B}{$\dfrac{200}3$ cm/s}
\opt{C}{$\dfrac{20}9$ cm/s}
\opt{D}{$\dfrac{10}9$ cm/s}
\end{qoptions}
\end{bankq}
\begin{bankq}[id=SA32-P2-PHY-S4-Q09, source={Score Advanced - Test 32 (08 May 2024)}, paper=2, qno=9, type=numerical, subtopic={Lenses}, answer={7}, answer_src=key, difficulty=8, diff_parts={K3 S1 I2 C2 F1}, diff_note={concepts: standard problem in this chapter, lenses, apply the relevant results | steps: apply the relevant results; compute}]
The given thin lenses have R$_1$ = R$_2$ = 10 cm and $\mu_1$ = 1, $\mu_2$ = 2 and $\mu_3$ = 4. If final image formed at object and value of `d' is x $\times$ 10 cm . Find `x'.\par
\bfig{128.8}{SA32_p2_phy_s4_q09_1.png}
\end{bankq}
\begin{bankq}[id=SA33-P2-PHY-S3-Q01, source={Score Advanced - Test 33 (12 May 2024)}, paper=2, qno=1, type=numerical, subtopic={Lenses}, answer={3}, answer_src=key, difficulty=8, diff_parts={K3 S1 I2 C2 F1}, diff_note={concepts: standard problem in this chapter, lenses, apply the relevant results | steps: apply the relevant results; compute}, also_in={Score Advanced - Test 34 (12 May 2024) P2 Q1}]
Two thin lenses with lens powers D$_1$ and D$_2$ diopter are located at distance L = 25 cm from each other and their optical axes coincide. This system creates an erect real image of the object, located at the optical axis closer to lens D$_1$ with the net magnification M$_1$ = 1. If the position of the two lenses are exchanged, the system again produces an erect real image with the net magnification M$_2$ = 4. What is the difference between lens powers $\Delta$D = D$_1$ $-$ D$_2$ (in diopter)?
\end{bankq}
\begin{bankq}[id=SA34-P1-PHY-S1-Q03, source={Score Advanced - Test 34 (12 May 2024)}, paper=1, qno=3, type=multi, subtopic={Prisms}, answer={A,B,C}, answer_src=key, difficulty=8, diff_parts={K3 S1 I2 C2 F1}, diff_note={concepts: standard problem in this chapter, prisms, apply the relevant results | steps: apply the relevant results; compute}, also_in={Score Advanced - Test 36 (19 May 2024) P1 Q3}]
By properly combining two prisms made of different materials, it is possible to :-
\begin{qoptions}
\opt{A}{have dispersion without average deviation.}
\opt{B}{have deviation without dispersion.}
\opt{C}{have both dispersion and average deviation.}
\opt{D}{have neither dispersion nor average deviation.}
\end{qoptions}
\end{bankq}
\begin{bankq}[id=SA36-P2-PHY-S4-Q07, source={Score Advanced - Test 36 (19 May 2024)}, paper=2, qno=7, type=numerical, subtopic={Prisms}, answer={5}, answer_src=key, difficulty=8, diff_parts={K3 S1 I2 C2 F1}, diff_note={concepts: standard problem in this chapter, prisms, apply the relevant results | steps: apply the relevant results; compute}, also_in={Score Advanced - Test 37 (19 May 2024) P2 Q7}]
A parallel beam of light is incident from air at an angle `$\alpha$' on the side PQ of a right angled triangular prism of refractive index n = $\sqrt2$. Light undergoes total internal reflection in the prism at the face PR when $\alpha$ has a minimum value of 45$^\circ$. The angle `$\theta$' of the prism is `n' times 3$^\circ$. Find `n'.\par
\bfig{76.0}{SA36_p2_phy_s4_q07_1.png}
\end{bankq}
\begin{bankq}[id=SA38-P1-PHY-S2-Q07, source={Score Advanced - Test 38 (13 Apr 2025)}, paper=1, qno=7, type=multi, subtopic={Refraction}, answer={C,D}, answer_src=key, difficulty=8, diff_parts={K3 S1 I2 C2 F1}, diff_note={concepts: standard problem in this chapter, refraction, apply the relevant results | steps: apply the relevant results; compute}]
A glass beaker has a solid, plano-convex base of refractive index 1.60, as shown in the figure. The radius of curvature of the convex surface (SPU) is 9 cm, while the planar surface (STU) acts as a mirror. This beaker is filled with a liquid of refractive index \textit n up to the level QPR. If the image of a point object O at a height of \textit h (OT in the figure) is formed onto itself, then, which of the following option(s) is(are) correct ?\par
\bfig{125.3}{SA38_p1_phy_s2_q07_1.png}
\begin{qoptions}
\opt{A}{For \textit n = 1.42, \textit h = 40 cm.}
\opt{B}{For \textit n = 1.35, \textit h = 52 cm.}
\opt{C}{For \textit n = 1.45, \textit h = 60 cm.}
\opt{D}{For \textit n = 1.48, \textit h = 75 cm.}
\end{qoptions}
\end{bankq}
\begin{bankq}[id=SA38-P1-PHY-S3-Q10, source={Score Advanced - Test 38 (13 Apr 2025)}, paper=1, qno=10, type=single, subtopic={Refraction}, answer={C}, answer_src=key, difficulty=4, diff_parts={K2 S1 I1 C1 F0}, diff_note={concepts: standard problem in this chapter, refraction | steps: apply the relevant results; compute}]
A light ray is incident on the surface of a sphere of refractive index n at an angle of incidence $\theta_0$. The ray partially refracts into the sphere with angle of refraction $\phi_0$ and then partly reflects from the back surface. The reflected ray then emerges out of the sphere after a partial refraction. The total angle of deviation of the emergent ray with respect to the incident ray is $\alpha$. Match the quantities mentioned in List-I with their values in List-II and choose the correct option.\par
{\footnotesize\setlength{\tabcolsep}{1.5pt}\noindent\begin{tabular}{|c|L{5.0cm}|c|L{2.0cm}|}
\hline
 & \multicolumn{1}{>{\centering\arraybackslash}p{5.0cm}|}{List-I} & & \multicolumn{1}{>{\centering\arraybackslash}p{2.0cm}|}{List-II} \\ \hline
(P) & If \textit n = 2 and $\alpha$ = 180$^\circ$, then all the possible values of $\theta_0$ will be & (1) & 60$^\circ$ and 0$^\circ$ \\ \hline
(Q) & If n = $\sqrt3$ and $\alpha$ = 180$^\circ$, then all the possible values of $\theta_0$ will be & (2) & 150$^\circ$ \\ \hline
(R) & If n = $\sqrt3$ and $\alpha$ = 180$^\circ$, then all the possible values of $\phi_0$ will be & (3) & 0$^\circ$ \\ \hline
(S) & If n = $\sqrt2$ and $\theta_0$ = 45$^\circ$, then all the possible values of $\alpha$ will be & (4) & 30$^\circ$ and 0$^\circ$ \\ \hline
 & & (5) & 45$^\circ$ and 0$^\circ$ \\ \hline
\end{tabular}}\par
\begin{qoptions}
\opt{A}{P $\to$ 3;Q $\to$ 2;R $\to$ 4;S $\to$ 3}
\opt{B}{P $\to$ 5;Q $\to$ 1;R $\to$ 2;S $\to$ 4}
\opt{C}{P $\to$ 3;Q $\to$ 1;R $\to$ 4;S $\to$ 2}
\opt{D}{P $\to$ 3;Q $\to$ 1;R $\to$ 2;S $\to$ 5}
\end{qoptions}
\end{bankq}
\begin{bankq}[id=SA38-P2-PHY-S4-Q08, source={Score Advanced - Test 38 (13 Apr 2025)}, paper=2, qno=8, type=numerical, subtopic={Prisms}, answer={23}, answer_src=key, difficulty=8, diff_parts={K3 S1 I2 C2 F1}, diff_note={concepts: standard problem in this chapter, prisms, apply the relevant results | steps: apply the relevant results; compute}]
Two equilateral-triangular prisms P$_1$ and P$_2$ are kept with their sides parallel to each other and very close in vacuum, as shown in the figure. A light ray enters prism P$_1$ at an angle of incidence $\theta$ such that the outgoing ray undergoes maximum deviation in prism P$_2$. If the respective refractive indices of P$_1$ and P$_2$ are $\dfrac53$ and 2, then $\theta$ = sin$^{-1}\left(\dfrac53\sin(\beta)\right)$, where the value of $\beta$ (in degree) is \underline{\hspace{1.5cm}}.\par
\bfig{113.8}{SA38_p2_phy_s4_q08_1.png}
\end{bankq}
\begin{bankq}[id=SA39-P1-PHY-S4-Q09, source={Score Advanced - Test 39 (20 Apr 2025)}, paper=1, qno=9, type=numerical, subtopic={Refraction}, answer={4}, answer_src=key, difficulty=8, diff_parts={K3 S1 I2 C2 F1}, diff_note={concepts: standard problem in this chapter, refraction, apply the relevant results | steps: apply the relevant results; compute}]
An ornament is lying at the bed of a pond of depth 8$\sqrt3$m. Pond is full of water whose refractive index is $\sqrt{\dfrac32}$. A diver at the surface of pond is trying to look at ornament such that her line of vision makes 30$^\circ$ with the surface of pond. Perpendicular distance of image of ornament from surface of water in metres is\underline{\hspace{1.5cm}}.
\end{bankq}
\begin{bankpassage}[id=SA39-P2-PHY-S3-G1-2, source={Score Advanced - Test 39 (20 Apr 2025)}, questions={1-2}]
In the window of the physics laboratory, an inverted image of the neighboring school was photographed, created by a spherical glass flask with a radius of R = 60 mm filled with water. Take $\mu_{\mathrm{water}}$ = $\dfrac43$.\par
\bfig{76.0}{SA39_p2_phy_s3_para1_2_1.png}
\end{bankpassage}
\begin{bankq}[id=SA39-P2-PHY-S3-Q01, source={Score Advanced - Test 39 (20 Apr 2025)}, paper=2, qno=1, type=numerical, subtopic={Refraction}, answer={Bonus}, answer_src=solved, working={Printed key for this paper is incomplete and the figure/data needed to solve it was not recoverable; banked as Bonus.}, difficulty=5, diff_parts={K2 S1 I1 C1 F1}, diff_note={concepts: standard problem in this chapter, refraction | steps: apply the relevant results; compute}, group=SA39-P2-PHY-S3-G1-2]
At what distance d$_1$ (in cm) from the center of the flask is the center of this image located?
\end{bankq}
\begin{bankq}[id=SA39-P2-PHY-S3-Q02, source={Score Advanced - Test 39 (20 Apr 2025)}, paper=2, qno=2, type=numerical, subtopic={Refraction}, answer={Bonus}, answer_src=solved, working={Printed key for this paper is incomplete and the figure/data needed to solve it was not recoverable; banked as Bonus.}, difficulty=5, diff_parts={K2 S1 I1 C1 F1}, diff_note={concepts: standard problem in this chapter, refraction | steps: apply the relevant results; compute}, group=SA39-P2-PHY-S3-G1-2]
Since the camera is focused on the image of the building created by the flask, the actual school building is captured only as a blur. However, you can quite accurately compare the viewing angle under which you see the window of the building from the location where the center of the lens is located with the viewing angle under which you see, from the same location, the image of the window created by the flask. In the photograph, the height of the window appears as H = 50 mm, and the height of the image of the window created by the flask as h = 6 mm. What is the distance d$_2$ (in cm) from the center of the lens to the center of the image created by the flask ?
\end{bankq}
\begin{bankq}[id=SA40-P1-PHY-S3-Q07, source={Score Advanced - Test 40 (27 Apr 2025)}, paper=1, qno=7, type=matching, subtopic={Reflection}, answer={Bonus}, answer_src=solved, working={Figure/data needed was not recoverable and no key given; banked as Bonus.}, difficulty=5, diff_parts={K2 S1 I1 C1 F1}, diff_note={concepts: standard problem in this chapter, reflection | steps: apply the relevant results; compute}]
A very small circular object is kept infront of an optical device as shown in figure. The plane of object is parallel to the optical device. One quadrant of circle is hatched as shown. Match the images as seen by the observer (ignoring magnification).\par
\bfig{95.3}{SA40_p1_phy_s3_q07_1.png}
{\footnotesize\setlength{\tabcolsep}{1.5pt}\noindent\begin{tabular}{|c|L{4.0cm}|c|c|}
\hline
 & \multicolumn{1}{>{\centering\arraybackslash}p{4.0cm}|}{List-I (Device)} & & \multicolumn{1}{c|}{List-II (Probable Image)} \\ \hline
(P) & Plane mirror (The observer is at A) & (1) & \bfigi{51.9}{SA40_p1_phy_s3_q07_2.png} \\ \hline
(Q) & Concave mirror (The observer is at A) object between F and 2F & (2) & \bfigi{67.3}{SA40_p1_phy_s3_q07_3.png} \\ \hline
(R) & Convex mirror (The observer is at A) & (3) & \bfigi{51.4}{SA40_p1_phy_s3_q07_4.png} \\ \hline
(S) & Convex lens (The observer is at B) object between F and optical centre & (4) & \bfigi{52.9}{SA40_p1_phy_s3_q07_5.png} \\ \hline
\end{tabular}}\par
\end{bankq}
\begin{bankq}[id=SA40-P1-PHY-S3-Q08, source={Score Advanced - Test 40 (27 Apr 2025)}, paper=1, qno=8, type=matching, subtopic={Reflection}, answer={P$\to$3; Q$\to$1; R$\to$4; S$\to$3}, answer_src=key, difficulty=8, diff_parts={K3 S1 I2 C2 F1}, diff_note={concepts: standard problem in this chapter, reflection, apply the relevant results | steps: apply the relevant results; compute}]
Match statement of List-I with List-II and select the correct answer using the code given below the list.\par
{\footnotesize\setlength{\tabcolsep}{1.5pt}\noindent\begin{tabular}{|c|L{5.2cm}|c|L{1cm}|}
\hline
 & \multicolumn{1}{>{\centering\arraybackslash}p{5.2cm}|}{List-I} & & \multicolumn{1}{c|}{List-II} \\ \hline
(P) & \bfigi{172.2}{SA40_p1_phy_s3_q08_1.png}\par Graph-1 and Graph-2 represents the displacement time graph of two identical particles, performing SHM with same total energy. The value of a$_1$/a$_2$ is \underline{\hspace{0.8cm}}. & (1) & 0 \\ \hline
(Q) & \bfigi{84.7}{SA40_p1_phy_s3_q08_2.png}\par A circuit contains two capacitors as shown in the figure. Initially switch S$_1$ is closed and S$_2$ is open. At t = 0, S$_2$ is closed keeping S$_1$ closed, and a charge Q$_1$ is supplied by the battery to the circuit. After long time, S$_1$ is opened keeping S$_2$ closed. Now battery supplied charge Q$_2$. The value of Q$_2$/Q$_1$ is & (2) & 1 \\ \hline
(R) & \bfigi{148.8}{SA40_p1_phy_s3_q08_3.png}\par A massive rope of uniform mass density and length $\ell$ is suspended from points A and B in vertical plane as shown in the figure. Point C is the lowest point of the chord. Given length of rope from A to C = $\ell_1$ and length of rope from B to C = $\ell_2$. The value of $\ell_2/\ell_1$ is & (3) & 2 \\ \hline
(S) & \bfigi{98.1}{SA40_p1_phy_s3_q08_4.png}\par A particle performs SHM along a straight line between the points M and N about O(mean position) as shown in figure. t$_1$ and t$_2$ are the time taken by the particle from M to P and P to N respectively. Find t$_1$/t$_2$ (given OP = PN) & (4) & 3 \\ \hline
\end{tabular}}\par
\end{bankq}
\begin{bankq}[id=SA40-P2-PHY-S4-Q03, source={Score Advanced - Test 40 (27 Apr 2025)}, paper=2, qno=3, type=numerical, subtopic={Prisms}, answer={6}, answer_src=key, difficulty=8, diff_parts={K3 S1 I2 C2 F1}, diff_note={concepts: standard problem in this chapter, prisms, apply the relevant results | steps: apply the relevant results; compute}]
A narrow beam of light is incident on a 30$^\circ-$60$^\circ-$90$^\circ$ prism kept in air as shown figure. The refractive index of the prism is $\mu$ = 2.1. The total clockwise deviation of the beam is $\dfrac{2\pi}{\mathrm N}$ radians, where N ($\neq$ 2) is:\par
\bfig{87.9}{SA40_p2_phy_s4_q03_1.png}
\end{bankq}
\begin{bankq}[id=SA41-P1-PHY-S1-Q02, source={Score Advanced - Test 41}, paper=1, qno=2, type=single, subtopic={Refraction}, answer={A}, answer_src=key, difficulty=4, diff_parts={K2 S1 I1 C1 F0}, diff_note={concepts: standard problem in this chapter, refraction | steps: apply the relevant results; compute}]
A light ray is incident on a transparent sphere of index = $\sqrt2$, at an angle of incidence = 45$^\circ$. What is the deviation of a tiny fraction of the ray, which enters the sphere, undergoes two internal reflections and then refracts out into air ?
\begin{qoptions}
\opt{A}{270$^\circ$}
\opt{B}{240$^\circ$}
\opt{C}{120$^\circ$}
\opt{D}{180$^\circ$}
\end{qoptions}
\end{bankq}
\begin{bankq}[id=SA41-P1-PHY-S3-Q10, source={Score Advanced - Test 41}, paper=1, qno=10, type=single, subtopic={Lenses}, answer={C}, answer_src=key, difficulty=4, diff_parts={K2 S1 I1 C1 F0}, diff_note={concepts: standard problem in this chapter, lenses | steps: apply the relevant results; compute}]
Two transparent media of refractive indices $\mu_1$ and $\mu_3$ have a solid lens shaped transparent material of refractive index $\mu_2$ between them as shown in figures in List$-$II. A ray traversing these media is also shown in the figures. In List$-$I different relationships between $\mu_1$, $\mu_2$ and $\mu_3$ are given. Match them to the ray diagrams shown in List$-$II.\par
{\footnotesize\setlength{\tabcolsep}{1.5pt}\noindent\begin{tabular}{|c|L{1.6cm}|c|c|}\hline & \multicolumn{1}{>{\centering\arraybackslash}p{1.6cm}|}{List-I} & & \multicolumn{1}{c|}{List-II} \\ \hline
(P) & $\mu_1<\mu_2$ & (1) & \bfigi{71.3}{SA41_p1_phy_s3_q10_1.png} \\ \hline
(Q) & $\mu_1>\mu_2$ & (2) & \bfigi{70.6}{SA41_p1_phy_s3_q10_2.png} \\ \hline
(R) & $\mu_2=\mu_3$ & (3) & \bfigi{71.3}{SA41_p1_phy_s3_q10_3.png} \\ \hline
(S) & $\mu_2>\mu_3$ & (4) & \bfigi{71.3}{SA41_p1_phy_s3_q10_4.png} \\ \hline
 & & (5) & \bfigi{71.3}{SA41_p1_phy_s3_q10_5.png} \\ \hline
\end{tabular}}\par
\begin{qoptions}
\opt{A}{P $\to$ 1,4;Q $\to$ 2,3;R $\to$ 3,4;S $\to$ 1,5}
\opt{B}{P $\to$ 1,4,5;Q $\to$ 1,3,5;R $\to$ 1,2;S $\to$ 1,5}
\opt{C}{P $\to$ 1,3;Q $\to$ 2,4,5;R $\to$ 1,3,5;S $\to$ 2,4}
\opt{D}{P $\to$ 1,3,5;Q $\to$ 2,4;R $\to$ 2,3;S $\to$ 4,5}
\end{qoptions}
\end{bankq}
\begin{bankq}[id=SA41-P2-PHY-S1-Q04, source={Score Advanced - Test 41}, paper=2, qno=4, type=single, subtopic={Reflection}, answer={B}, answer_src=key, difficulty=4, diff_parts={K2 S1 I1 C1 F0}, diff_note={concepts: standard problem in this chapter, reflection | steps: apply the relevant results; compute}]
If $\theta$ =110$^\circ$ then total number of images formed by the mirror system will be :-\par
\bfig{84.6}{SA41_p2_phy_s1_q04_1.png}
\begin{qoptions}
\opt{A}{2}
\opt{B}{3}
\opt{C}{4}
\opt{D}{5}
\end{qoptions}
\end{bankq}
\begin{bankpassage}[id=SA41-P2-PHY-S3-G1-2, source={Score Advanced - Test 41}, questions={1-2}]
A semi cylinder made of a transparent plastic has a refraction index of n = $\sqrt2$ and a radius of R. There is a narrow incident light ray perpendicular to the flat side of the semi cylinder at d distance from the axis of symmetry.\par
\bfig{108.5}{SA41_p2_phy_s3_para1_2_1.png}
\end{bankpassage}
\begin{bankq}[id=SA41-P2-PHY-S3-Q01, source={Score Advanced - Test 41}, paper=2, qno=1, type=numerical, subtopic={Refraction}, answer={7.04}, answer_src=key, difficulty=8, diff_parts={K3 S1 I2 C2 F1}, diff_note={concepts: standard problem in this chapter, refraction, apply the relevant results | steps: apply the relevant results; compute}, group=SA41-P2-PHY-S3-G1-2]
What can the maximum value of d (in cm) be so that the light ray can still leave the other side of the semi cylinder? [Take : R = 10 cm]
\end{bankq}
\begin{bankq}[id=SA41-P2-PHY-S3-Q02, source={Score Advanced - Test 41}, paper=2, qno=2, type=numerical, subtopic={Refraction}, answer={1.20}, answer_src=key, difficulty=8, diff_parts={K3 S1 I2 C2 F1}, diff_note={concepts: standard problem in this chapter, refraction, apply the relevant results | steps: apply the relevant results; compute}, group=SA41-P2-PHY-S3-G1-2]
When the value of d chosen is such that TIR just takes place then time (in nanoseconds) for which light remains inside the cylinder is (take R = 9cm \& C = 3 $\times$ 10$^8$ m/s)
\end{bankq}

\section*{TOAS papers}
\begin{bankq}[id=TOAS02-P1-PHY-S1-Q04, source={TOAS - Test 02 (13 Feb 2022)}, paper=1, qno=4, type=single, subtopic={Refraction}, answer={B}, answer_src=key, difficulty=1, diff_parts={K1 S1 I0 C0 F0}, diff_note={concepts: refraction and total internal reflection | steps: Snell's law at the top face; condition at the vertical face for total reflection}]
In the figure shown, for an angle of incidence $i$ at the top surface, what is the minimum refractive index needed for total internal reflection at the vertical face?
\bfig{75.1}{TOAS02_p1_phy_s1_q04_1.png}
\begin{qoptions}
\opt{A}{$\sqrt{1-\sin^2i}$}
\opt{B}{$\sqrt{1+\sin^2i}$}
\opt{C}{$\sqrt{1-\cos^2i}$}
\opt{D}{$\sqrt{1+\cos^2i}$}
\end{qoptions}
\end{bankq}
\begin{bankq}[id=TOAS02-P1-PHY-S1-Q05, source={TOAS - Test 02 (13 Feb 2022)}, paper=1, qno=5, type=single, subtopic={Lenses}, answer={D}, answer_src=key, difficulty=2, diff_parts={K2 S1 I0 C0 F0}, diff_note={concepts: lens imaging, plane mirror reflection | steps: image of A by the lens; reflection of the image in the inclined mirror}]
A point $A$ is kept on the axis of a converging lens of focal length 10 cm as shown in the figure. A plane mirror $M$ is placed inclined at $45^\circ$ to the lens axis at a distance of 10 cm to the right of the lens. The co-ordinates of the image of $A$ formed by this combination are (take the centre of the lens as the origin and the lens axis as the axis of $x$). (Consider only two optical phenomenon)
\bfig{115.4}{TOAS02_p1_phy_s1_q05_1.png}
\begin{qoptions}
\opt{A}{($-10$ cm, 20 cm)}
\opt{B}{($-10$ cm, $-20$ cm)}
\opt{C}{(10 cm, $-20$ cm)}
\opt{D}{(10 cm, 20 cm)}
\end{qoptions}
\end{bankq}
\begin{bankq}[id=TOAS02-P1-PHY-S2-Q07, source={TOAS - Test 02 (13 Feb 2022)}, paper=1, qno=7, type=multi, subtopic={Refraction}, answer={A}, answer_src=key, difficulty=3, diff_parts={K1 S2 I0 C0 F1}, diff_note={concepts: refraction at a spherical surface | steps: plane face passes the beam undeviated; refraction at the curved face; second case with refraction at the curved face first and then the plane face}]
A parallel beam of light is incident normally on the flat surface of a hemisphere of radius 6cm and refractive index 1.5, placed in air as shown in figure (i). Assume paraxial ray approximation. Choose the correct option(s) :-
\bfig{130.6}{TOAS02_p1_phy_s2_q07_1.png}
\begin{qoptions}
\opt{A}{The rays are focussed at 12 cm from the point $P$ to the right, in the situation as shown in figure (i)}
\opt{B}{The rays are focussed at 16 cm from the point $P$ to the right, in the situation as shown in figure (i)}
\opt{C}{If the rays are incident at the curved surface (figure (ii)) then these are focused at distance 18 cm from point $P$ to the right}
\opt{D}{If the rays are incident at the curved surface (figure (ii)) then these are focused at distance 14 cm from point $P$ to the right}
\end{qoptions}
\end{bankq}
\begin{bankq}[id=TOAS02-P1-PHY-S2-Q08, source={TOAS - Test 02 (13 Feb 2022)}, paper=1, qno=8, type=multi, subtopic={Reflection}, answer={A}, answer_src=key, difficulty=3, diff_parts={K1 S1 I1 C0 F1}, diff_note={concepts: image formation by a spherical mirror | steps: object and image on opposite sides of the axis means real inverted image; mirror type and position | I: real image formed by a converging mirror}]
Two points $P$ and $Q$ lie on either side of an axis $XY$ as shown. It is desired to produce an image at $Q$ for a real object placed at $P$ using a spherical mirror, with $XY$ as the optic axis. The mirror must be
\bfig{129.4}{TOAS02_p1_phy_s2_q08_1.png}
\begin{qoptions}
\opt{A}{converging and positioned to the left of $P$}
\opt{B}{diverging and positioned to the left of $P$}
\opt{C}{converging and positioned to the right of $Q$}
\opt{D}{diverging and positioned to the right of $Q$}
\end{qoptions}
\end{bankq}
\begin{bankq}[id=TOAS02-P1-PHY-S3-Q01, source={TOAS - Test 02 (13 Feb 2022)}, paper=1, qno=1, type=numerical, subtopic={Lenses}, answer={3}, answer_src=key, difficulty=5, diff_parts={K2 S2 I1 C0 F1}, diff_note={concepts: lens maker formula, silvered lens as a mirror | steps: equivalent focal length with the plane face silvered; with the convex face silvered; solve both for the refractive index | I: power of lens used twice plus the mirror}]
When the plane surface of plano-convex lens is silvered, the lens behaves like a concave mirror of focal length 30 cm, but when the convex surface is silvered, it behaves like a concave mirror of focal length 10 cm. The refractive index of material of lens is $\dfrac{P}{2}$. What is the value of P?
\end{bankq}
\begin{bankq}[id=TOAS02-P1-PHY-S3-Q03, source={TOAS - Test 02 (13 Feb 2022)}, paper=1, qno=3, type=numerical, subtopic={Prisms}, answer={4}, answer_src=key, difficulty=7, diff_parts={K2 S3 I1 C1 F1}, diff_note={concepts: prism deviation, Snell's law | steps: emergent angle from the deviation formula; Snell's law at the first face; Snell's law at the second face; r1 + r2 = alpha; eliminate the angles; express in cos\^{}2 alpha}]
A ray of light passes through a prism in principle plane. The angle of prism is $\alpha$. Angle of deviation and angle of incidence both are equal to $2\alpha$. If $\mu$ is the refractive index of the prism then it can be found that $\cos^2\alpha=\dfrac{\mu^n-1}{2x}$ where $n$ is a positive integer. The value of $x$ is:
\end{bankq}
\begin{bankq}[id=TOAS02-P1-PHY-S3-Q04, source={TOAS - Test 02 (13 Feb 2022)}, paper=1, qno=4, type=numerical, subtopic={Reflection}, answer={3.46}, answer_src=key, difficulty=3, diff_parts={K1 S1 I1 C0 F1}, diff_note={concepts: plane mirror image geometry | steps: image of B behind the mirror; the guard sees B when the line from S to the image passes through the mirror edge | I: image-edge condition}]
Figure shows an overhead view of a corridor with a plane mirror $M$ mounted at one end. A burglar $B$ sneaks along the corridor directly towards the centre of the mirror. If $d=4$m, how far (in meter) from the mirror is the burglar when he will be first seen in the mirror by the security guard $S$ ?
\bfig{167.5}{TOAS02_p1_phy_s3_q04_1.png}
\end{bankq}
\begin{bankq}[id=TOAS02-P1-PHY-S3-Q06, source={TOAS - Test 02 (13 Feb 2022)}, paper=1, qno=6, type=numerical, subtopic={Lenses}, answer={2.50}, answer_src=key, difficulty=4, diff_parts={K2 S2 I0 C0 F1}, diff_note={concepts: diverging lens, plane mirror | steps: image formed by the lens first; reflection in the mirror; second refraction through the lens; position relative to the mirror}]
A diverging lens of focal length 10 cm is placed 10 cm in front of a plane mirror as shown in the figure. Light from a very far away source falls on the lens. The final image is at a distance "$V$" cm from the mirror. Write the value of $V$ (positive) :
\bfig{126.3}{TOAS02_p1_phy_s3_q06_1.png}
\end{bankq}
\begin{bankq}[id=TOAS03-P1-PHY-S1-Q01, source={TOAS - Test 03 (27 Feb 2022)}, paper=1, qno=1, type=single, subtopic={Lenses}, answer={A}, answer_src=key, difficulty=2, diff_parts={K1 S2 I0 C0 F0}, diff_note={concepts: lens combination magnification | steps: image of the first lens; object distance for the second lens; magnification of each lens and the product}]
Find the lateral magnification produced by the combination of lenses shown in the figure.
\bfig{124.6}{TOAS03_p1_phy_s1_q01_1.png}
\begin{qoptions}
\opt{A}{2}
\opt{B}{4}
\opt{C}{6}
\opt{D}{8}
\end{qoptions}
\end{bankq}
\begin{bankq}[id=TOAS03-P1-PHY-S1-Q03, source={TOAS - Test 03 (27 Feb 2022)}, paper=1, qno=3, type=single, subtopic={Refraction}, answer={C}, answer_src=key, difficulty=3, diff_parts={K1 S2 I1 C0 F0}, diff_note={concepts: total internal reflection at a curved surface | steps: critical angle sin C = 1/2; angle of incidence at the surface equals theta for the axial geometry; minimum theta | I: normal passes through the centre}]
The minimum value of $\theta$ so that the ray gets totally reflected from spherical surface of radius of curvature $R$ is $\left[\dfrac{\mu_2}{\mu_1}=\dfrac12\right]$
\bfig{108.7}{TOAS03_p1_phy_s1_q03_1.png}
\begin{qoptions}
\opt{A}{$30^\circ$}
\opt{B}{$60^\circ$}
\opt{C}{zero}
\opt{D}{not possible}
\end{qoptions}
\end{bankq}
\begin{bankq}[id=TOAS03-P1-PHY-S2-Q07, source={TOAS - Test 03 (27 Feb 2022)}, paper=1, qno=7, type=multi, subtopic={Lenses}, answer={A,B,C}, answer_src=key, difficulty=2, diff_parts={K1 S1 I0 C0 F1}, diff_note={concepts: lens equation and magnification | steps: m = v/u and the lens formula give m = 1 - v/f; read intercepts and slope}]
Consider the linear graph between magnification ($m$) and image position ($v$) for a lens. Which of the following statement about the straight line is correct ?
\begin{qoptions}
\opt{A}{Focal length of the lens is equal to intercept on $v$-axis}
\opt{B}{Magnitude of focal length of the lens is equal to inverse of slope of the line.}
\opt{C}{Magnitude of intercept on $m$-axis is equal to unity}
\opt{D}{None of these}
\end{qoptions}
\end{bankq}
\begin{bankq}[id=TOAS03-P1-PHY-S2-Q09, source={TOAS - Test 03 (27 Feb 2022)}, paper=1, qno=9, type=multi, subtopic={Reflection}, answer={A,C}, answer_src=key, difficulty=4, diff_parts={K1 S2 I1 C0 F1}, diff_note={concepts: image velocity in a plane mirror | steps: image velocity is the mirror reflection of the object velocity; relative speed over the circle; extremes of the relative speed; compare with the options | I: reflect the velocity vector}]
An object is moving in a circular path of radius 1 m around (0, 10 m) in $xy$ plane with constant speed 1 m/s. Speed of image w.r.t. object is $v$.
\bfig{113.4}{TOAS03_p1_phy_s2_q09_1.png}
\begin{qoptions}
\opt{A}{$v_{\max}=2$ m/s}
\opt{B}{$v_{\min}=1$ m/s}
\opt{C}{$v$ is minimum when object is closest or farthest from the mirror}
\opt{D}{$v$ is maximum when object crosses $y$--axis}
\end{qoptions}
\end{bankq}
\begin{bankq}[id=TOAS03-P1-PHY-S3-Q01, source={TOAS - Test 03 (27 Feb 2022)}, paper=1, qno=1, type=numerical, subtopic={Reflection}, answer={0.50}, answer_src=key, difficulty=4, diff_parts={K1 S2 I1 C0 F1}, diff_note={concepts: plane mirror sight lines | steps: image of A in the mirror; line from B to the image of A through the mirror edge; distance of B | I: sight line through the image}]
A man `A' stands at the position shown in the figure and a second man `B' approaches the mirror along the line perpendicular to it which passes through its centre. At what distance (in meter) from the mirror will `B' be at the moment when `A' and `B' first see each other in the mirror ?
\bfig{99.6}{TOAS03_p1_phy_s3_q01_1.png}
\end{bankq}
\begin{bankq}[id=TOAS03-P1-PHY-S3-Q03, source={TOAS - Test 03 (27 Feb 2022)}, paper=1, qno=3, type=numerical, subtopic={Refraction}, answer={1.72 to 1.73}, answer_src=key, difficulty=6, diff_parts={K2 S3 I1 C0 F1}, diff_note={concepts: refraction in a sphere, reflection at the back surface | steps: Snell's law at the entry; deviation of the transmitted ray; deviation after one internal reflection; the 1/3 condition; solve for n | I: deviation of a ray through a sphere}]
A ray is incident on a glass sphere as shown. The opposite surface of the sphere is partially silvered. If the net deviation of the ray transmitted at the partially silvered surface is $1/3$rd of the net deviation suffered by the ray coming out after one reflection, the refractive index of the sphere is $n$. Find the value of $n$.
\bfig{128.6}{TOAS03_p1_phy_s3_q03_1.png}
\end{bankq}
\begin{bankq}[id=TOAS03-P2-PHY-S1-Q02, source={TOAS - Test 03 (27 Feb 2022)}, paper=2, qno=2, type=multi, subtopic={Lenses}, answer={A,B,C,D}, answer_src=key, difficulty=3, diff_parts={K1 S2 I0 C0 F1}, diff_note={concepts: imaging by a lens-mirror system | steps: image position as a function of object position; direction of image motion; speeds compared}]
When object `O' moves towards a fixed lens mirror combination, select correct choice/choices:
\bfig{177.5}{TOAS03_p2_phy_s1_q02_1.png}
\begin{qoptions}
\opt{A}{Image moves towards negative $x$-axis.}
\opt{B}{Speed of image and object may be same}
\opt{C}{Image may move faster than object}
\opt{D}{Image may come closer to arrangement}
\end{qoptions}
\end{bankq}
\begin{bankq}[id=TOAS03-P2-PHY-S1-Q04, source={TOAS - Test 03 (27 Feb 2022)}, paper=2, qno=4, type=multi, subtopic={Lenses}, answer={A,B,C}, answer_src=key, difficulty=4, diff_parts={K1 S2 I1 C0 F1}, diff_note={concepts: lens in media of different refractive indices | steps: power as (2 n - 1 - mu)/R; character for mu > 2; focal length for mu = 4; same distances for mu = 2 | I: generalised lens maker formula}]
An equiconvex lens of focal length 20 cm having a glass of refractive index $\dfrac32$ is used in a medium air on left side and another medium having refractive index $\mu$ on the right side of lens as shown in figure.
\bfig{152.5}{TOAS03_p2_phy_s1_q04_1.png}
\begin{qoptions}
\opt{A}{Character of lens will be diverging for $\mu>2$}
\opt{B}{Focal length of lens will be 40 cm for $\mu=4$ and character will be diverging}
\opt{C}{Object distance and image distance with respect to lens will be same for $\mu=2$}
\opt{D}{Character of lens can't be explained}
\end{qoptions}
\end{bankq}
\begin{bankq}[id=TOAS03-P2-PHY-S1-Q05, source={TOAS - Test 03 (27 Feb 2022)}, paper=2, qno=5, type=multi, subtopic={Prisms}, answer={A,B}, answer_src=key, difficulty=3, diff_parts={K1 S1 I1 C0 F1}, diff_note={concepts: deviation through identical prisms in contact | steps: alternate prisms cancel the deviation; parity of n | I: alternate orientation cancels}, also_in={Nurture (General) - Test 01 (12 Jun 2021) P1 Q7}]
$n$ number of identical equilateral prisms are kept in contact as shown in figure. If deviation through a single prism is $\delta$, then ($n$, $m$ are integers)
\bfig{182.5}{TOAS03_p2_phy_s1_q05_1.png}
\begin{qoptions}
\opt{A}{if $n=2m$, deviation through $n$ prism is zero.}
\opt{B}{if $n=2m+1$, deviation through system of $n$ prisms is $\delta$.}
\opt{C}{if $n=2m$, deviation through system of $n$ prism is $\delta$.}
\opt{D}{if $n=2m+1$ deviation through system of $n$ prisms is zero}
\end{qoptions}
\end{bankq}
\begin{bankq}[id=TOAS03-P2-PHY-S1-Q06, source={TOAS - Test 03 (27 Feb 2022)}, paper=2, qno=6, type=multi, subtopic={Lenses}, answer={A,C}, answer_src=key, difficulty=6, diff_parts={K2 S3 I0 C1 F1}, diff_note={concepts: lens maker formula, sequential imaging | steps: radius of curvature of each lens; focal length of each lens with water in contact; image of the first lens; image of the second lens; size}]
Consider an arrangement of two thin equiconvex lenses made of glass ($\mu_g=3/2$) and each of focal length in air 10 cm are separated by a distance 10 cm. The space between the two lenses is filled with water of refractive index $\mu_W=\dfrac43$. A small linear object `AB' of size 0.20 cm is placed perpendicular to the optic axis at a distance 10 cm from the first lens in air as shown. Then choose the correct option(s) :
\bfig{140.6}{TOAS03_p2_phy_s1_q06_1.png}
\begin{qoptions}
\opt{A}{The position of the final image formed is 25 cm right to the second lens.}
\opt{B}{The position of the final image formed is 40 cm right to the second lens.}
\opt{C}{The size of the final image formed is 0.40 cm.}
\opt{D}{The size of the final image formed is 0.60 cm.}
\end{qoptions}
\end{bankq}
\begin{bankq}[id=TOAS03-P2-PHY-S2-Q03, source={TOAS - Test 03 (27 Feb 2022)}, paper=2, qno=3, type=numerical, subtopic={Refraction}, answer={0.37 to 0.40}, answer_src=key, difficulty=4, diff_parts={K1 S2 I1 C0 F1}, diff_note={concepts: refraction in a medium with continuously varying index | steps: Snell's invariant mu sin theta; horizontal ray at P gives mu(P) = sqrt3; solve for H | I: invariant of a stratified medium}]
A system of coordinates is drawn in a medium whose refractive index varies as $\mu=\dfrac2{1+y^2}$, where $0\leqslant y\leqslant1$. A ray of light is incident at origin at an angle $60^\circ$ with $y$-axis as shown in the figure. At point $P$ ray becomes parallel to $x$-axis. Find the value of `$H$'. (Take $\mu=2$ for $y<0$)
\bfig{124.2}{TOAS03_p2_phy_s2_q03_1.png}
\end{bankq}
\begin{bankq}[id=TOAS03-P2-PHY-S3-Q01, source={TOAS - Test 03 (27 Feb 2022)}, paper=2, qno=1, type=numerical, subtopic={Lenses}, answer={6}, answer_src=key, difficulty=3, diff_parts={K1 S1 I1 C0 F1}, diff_note={concepts: images by halves of a lens | steps: image positions from the two halves; their separation | I: displaced optical centres}]
Distance between two images formed by upper and lower part of the point object placed at 30 cm from given lens is $(60+x)$ cm, then the value of $x$ is :
\bfig{166.8}{TOAS03_p2_phy_s3_q01_1.png}
\end{bankq}
\begin{bankq}[id=TOAS03-P2-PHY-S3-Q03, source={TOAS - Test 03 (27 Feb 2022)}, paper=2, qno=3, type=numerical, subtopic={Prisms}, answer={3}, answer_src=key, difficulty=3, diff_parts={K1 S1 I1 C0 F1}, diff_note={concepts: minimum deviation of a prism | steps: mu = sin((A + delta)/2)/sin(A/2) with delta = A/2; solve for A | I: find A by trying standard angles}]
The minimum deviation suffered by a ray emerging through a prism $\left(\mu=\dfrac2{\sqrt3}\right)$ is equal to half the refracting angle of the prism. Find the value of $C$ if the refracting angle of the prism is $\left(\dfrac{2\pi}C\right)$.
\end{bankq}
\begin{bankq}[id=TOAS03-P2-PHY-S3-Q04, source={TOAS - Test 03 (27 Feb 2022)}, paper=2, qno=4, type=numerical, subtopic={Reflection}, answer={7}, answer_src=key, difficulty=2, diff_parts={K1 S1 I0 C0 F1}, diff_note={concepts: image distance in a plane mirror | steps: image is 3 m behind the mirror so 7 m from the camera}]
It is desired to photograph the image of an object placed at a distance of 3 m from the plane mirror. The camera which is at a distance of 4 m from the mirror should be focussed for a distance of (in m)
\end{bankq}
\begin{bankq}[id=TOAS03-P2-PHY-S3-Q06, source={TOAS - Test 03 (27 Feb 2022)}, paper=2, qno=6, type=numerical, subtopic={Reflection}, answer={9}, answer_src=key, difficulty=2, diff_parts={K1 S1 I0 C0 F1}, diff_note={concepts: reflection of rays | steps: right angle between incident and reflected rays gives incidence 45 degrees so h = 9 m}]
When a plane mirror is placed horizontally on level ground at a distance of 9 m from the foot of a tower, the ray coming from the top of the tower and its reflected ray from the mirror subtends an angle of $90^\circ$. What is the height of the tower ?
\end{bankq}
\begin{bankq}[id=TOAS04-P1-PHY-S1-Q01, source={TOAS - Test 04 (15 Mar 2022)}, paper=1, qno=1, type=single, subtopic={Refraction}, answer={A}, answer_src=key, difficulty=3, diff_parts={K1 S2 I1 C0 F0}, diff_note={concepts: apparent depth and related rates | steps: height h = V/(pi r\^{}2); apparent position shifts by (1 - 1/mu) h; differentiate with dr/dt = -K | I: differentiate the apparent depth}]
A cylinder is filled with a liquid of refractive index $\mu$. The radius of the cylinder is decreasing at a constant rate $K$. The volume of the liquid inside the container remains constant at $V$. The observer and the object $O$ are in a state of rest and at a distance $L$ from each other. The apparent velocity of the object as seen by the observer, is : (when radius of cylinder is $r$)
\bfig{53.4}{TOAS04_p1_phy_s1_q01_1.png}
\begin{qoptions}
\opt{A}{$\dfrac{(1-\mu)2KV}{\pi\mu r^3}$}
\opt{B}{$\dfrac{(1-\mu)2KV}{\pi\mu Lr^2}$}
\opt{C}{$\dfrac{(1-\mu)2K}{\mu}$}
\opt{D}{$\dfrac{(1-\mu)K}{2\mu}$}
\end{qoptions}
\end{bankq}
\begin{bankq}[id=TOAS04-P1-PHY-S2-Q08, source={TOAS - Test 04 (15 Mar 2022)}, paper=1, qno=8, type=multi, subtopic={Lenses}, answer={B,C,D}, answer_src=key, difficulty=4, diff_parts={K1 S2 I1 C0 F1}, diff_note={concepts: intensity of a defocused image | steps: image distance from the lens equation; spot radius on the screen from similar triangles; intensity variation with screen distance | I: spot size versus intensity}]
A point source of light $S$ is placed on the axis of the lens of focal length 20 cm as shown in the figure. A screen is placed normal to the axis of the lens at a distance $x$ from it. Treat all rays as paraxial. Then which of the following is/are correct regarding the intensity of the image produced on the screen of the point source ?
\bfig{148.5}{TOAS04_p1_phy_s2_q08_1.png}
\begin{qoptions}
\opt{A}{As $x$ is increased from zero, intensity continuously decreases.}
\opt{B}{As $x$ is increased from zero, intensity first increases then decreases.}
\opt{C}{Intensity at centre of screen for $x=90$ cm and $x=110$ cm is same.}
\opt{D}{Radius of bright circle obtained on screen is equal to 2 cm for $x=300$ cm.}
\end{qoptions}
\end{bankq}
\begin{bankq}[id=TOAS04-P1-PHY-S2-Q09, source={TOAS - Test 04 (15 Mar 2022)}, paper=1, qno=9, type=multi, subtopic={Reflection}, answer={A,B,C}, answer_src=key, difficulty=4, diff_parts={K1 S2 I1 C0 F1}, diff_note={concepts: images by the two halves of a mirror | steps: image position for the original mirror; shift of the images for lateral displacement; coincidence for rotation about O | I: halves behave as shifted mirrors}]
A concave mirror forms real image of a point object `$O$' lying on the optical axis at a distance of 60 cm from the pole of mirror. The focal length of mirror is 30 cm. Now, the mirror is cut into two halves symmetrically. In figure-1 both halves are shifted at a distance of 2 mm apart in a direction perpendicular to the optical axis. In figure-2 both halves are rotated (before being separated) about `$O$' with reference to original position. Let $I_1$ and $I_2$ are the image formed by both halves in both figure. Choose the correct statement(s).
\bfig{287.3}{TOAS04_p1_phy_s2_q09_1.png}

\begin{qoptions}
\opt{A}{In figure-2, $I_1$ and $I_2$ are coincide.}
\opt{B}{In figure-1, $I_1$ and $I_2$ are separated.}
\opt{C}{In figure-1, the distance between $I_1$ and $I_2$ is 8 mm.}
\opt{D}{In figure-2, the distance between $I_1$ and $I_2$ is 4 mm.}
\end{qoptions}
\end{bankq}
\begin{bankq}[id=TOAS04-P1-PHY-S2-Q10, source={TOAS - Test 04 (15 Mar 2022)}, paper=1, qno=10, type=multi, subtopic={Reflection}, answer={A,C}, answer_src=key, difficulty=3, diff_parts={K1 S1 I1 C0 F1}, diff_note={concepts: reflection from two inclined plane mirrors | steps: deviation after two reflections is 2(180 - 50) independent of theta; allowed range of theta | I: deviation independent of the incidence angle}]
Two plane mirrors are placed in such a way that they make an angle $50^\circ$ as shown in figure. A ray of light falls on one mirror and it reflects from each mirror once. For the given situation mark the correct statement(s).
\bfig{126.8}{TOAS04_p1_phy_s2_q10_1.png}
\begin{qoptions}
\opt{A}{The value of $\theta$ must be less than $40^\circ$.}
\opt{B}{$\theta$ can have any value ranging from $0^\circ$ to $90^\circ$.}
\opt{C}{Total deviation produced in light ray is independent of value $\theta$.}
\opt{D}{Laws of reflection are valid for only plane and spherical mirrors.}
\end{qoptions}
\end{bankq}
\begin{bankpassage}[id=TOAS04-P1-PHY-S3-G1-2, source={TOAS - Test 04 (15 Mar 2022)}, questions={1-2}]
An experiment is performed to calculate the focal length of a convex mirror using convex lens. The setup consists of two pins, optical bench, one convex lens and one convex mirror. When only convex lens is placed at a distance 60 cm from pin (object), its image is formed at a distance 30 cm from the lens as shown in figure-(i). Now convex mirror is also kept right to the lens without changing distance between pin (object) and lens. The mirror is kept at a distance 10 cm from lens such that position of image of object pin and position of object pin coincides as shown in figure-(ii)
\bfig{263.0}{TOAS04_p1_phy_s3_para1_2_1.png}

Now answer the following questions.
\end{bankpassage}
\begin{bankq}[id=TOAS04-P1-PHY-S3-Q01, source={TOAS - Test 04 (15 Mar 2022)}, paper=1, qno=1, type=numerical, subtopic={Lenses}, answer={20.00}, answer_src=key, difficulty=2, diff_parts={K1 S1 I0 C0 F1}, diff_note={concepts: lens equation from image data | steps: object and image distances for the convex lens; focal length}, group=TOAS04-P1-PHY-S3-G1-2]
Focal length of convex lens used (in cm) is :
\end{bankq}
\begin{bankq}[id=TOAS04-P1-PHY-S3-Q02, source={TOAS - Test 04 (15 Mar 2022)}, paper=1, qno=2, type=numerical, subtopic={Lenses}, answer={10.00}, answer_src=key, difficulty=2, diff_parts={K1 S1 I0 C0 F1}, diff_note={concepts: mirror equation from image data | steps: object and image distances for the convex mirror; focal length}, group=TOAS04-P1-PHY-S3-G1-2]
Focal length of convex mirror (in cm) is :
\end{bankq}
\begin{bankpassage}[id=TOAS04-P1-PHY-S3-G5-6, source={TOAS - Test 04 (15 Mar 2022)}, questions={5-6}]
A ray of light is incident normally on one face $30^\circ$--$60^\circ$--$90^\circ$ prism of refractive index $\mu_P=5/3$ immersed in water of refractive index $\mu_m=4/3$ as shown in the figure.
\bfig{165.6}{TOAS04_p1_phy_s3_para5_6_1.png}
\end{bankpassage}
\begin{bankq}[id=TOAS04-P1-PHY-S3-Q05, source={TOAS - Test 04 (15 Mar 2022)}, paper=1, qno=5, type=numerical, subtopic={Prisms}, answer={10.00}, answer_src=key, difficulty=3, diff_parts={K1 S2 I0 C0 F1}, diff_note={concepts: prism refraction | steps: refraction at the first face; angle at the second face; emergent angle}, group=TOAS04-P1-PHY-S3-G5-6]
The exit angle $\theta_2$ of the ray is $\sin^{-1}\left(\dfrac x{16}\right)$, find the value of $x$.
\end{bankq}
\begin{bankq}[id=TOAS04-P1-PHY-S3-Q06, source={TOAS - Test 04 (15 Mar 2022)}, paper=1, qno=6, type=numerical, subtopic={Prisms}, answer={2.50}, answer_src=key, difficulty=2, diff_parts={K1 S1 I0 C0 F1}, diff_note={concepts: total internal reflection at a prism face | steps: critical angle condition at P; refractive index of the surrounding medium}, group=TOAS04-P1-PHY-S3-G5-6]
The minimum refractive index of surrounding medium so that total internal reflection at point $P$ ceases is $\dfrac n{\sqrt3}$, find the value of $n$.
\end{bankq}
\begin{bankq}[id=TOAS04-P1-PHY-S4-Q07, source={TOAS - Test 04 (15 Mar 2022)}, paper=1, qno=7, type=numerical, subtopic={Lenses}, answer={35}, answer_src=key, difficulty=4, diff_parts={K2 S1 I1 C0 F1}, diff_note={concepts: lens maker formula, silvered surface | steps: object rays must meet the silvered surface normally; refraction at the first face gives the object distance | I: image coincides with the object when the ray hits the mirror normally}]
A biconvex thin lens is prepared from glass of refractive index $3/2$. The two bounding surfaces have equal radii of 35 cm each. One of the surfaces is silvered to make it reflecting. An object is placed infront of this lens in such a way that the image coincides with the object. If distance (in cm) of object from lens is $x$, then the value of $2x$ is
\end{bankq}
\begin{bankq}[id=TOAS04-P1-PHY-S4-Q09, source={TOAS - Test 04 (15 Mar 2022)}, paper=1, qno=9, type=numerical, subtopic={Reflection}, answer={60}, answer_src=key, difficulty=4, diff_parts={K1 S2 I1 C0 F1}, diff_note={concepts: image of an insect in a mirror beside a sphere | steps: tangent ray from the insect to the sphere; its reflection at the mirror; radius from similar triangles | I: limiting tangent ray}]
A circular plane mirror is kept on a horizontal surface, an opaque sphere of radius $R=20$ cm is placed symmetrically on the circular plane mirror. A small insect lies on the sphere at a vertical height of 36 cm from the circular mirror. Find the minimum radius (in cm) of mirror so that image of insect is formed by the mirror.
\end{bankq}
\begin{bankq}[id=TOAS04-P2-PHY-S1-Q01, source={TOAS - Test 04 (15 Mar 2022)}, paper=2, qno=1, type=multi, subtopic={Refraction}, answer={A,B,C}, answer_src=key, difficulty=6, diff_parts={K2 S2 I2 C0 F1}, diff_note={concepts: Snell's invariant, circular path geometry | steps: Snell's invariant mu sin theta; circular path geometry gives mu(y); maximum mu fixes the maximum angle | I: invariant and circular path}]
The refractive index of the medium with a certain region, $x>0$, $y>0$, changes with $y$. A thin light ray travelling in the $x$-direction in medium having refractive index $\mu_0=1$ strikes another medium of refractive index $\mu$ at right angles and moves through the medium along a circular arc of radius $R$ as shown in the figure. The material with the greatest known refractive index is diamond, but even the refractive index of this material does not reach the value $\mu_{\max}=2.5$. It is this limit that sets the maximum angular size of the arc the light ray can cover. Angular size of arc is the angle subtended by the arc at the centre.
\bfig{111.1}{TOAS04_p2_phy_s1_q01_1.png}
\begin{qoptions}
\opt{A}{The variation of refractive index $\mu$ with $y$ is given as $\mu=\dfrac{R}{R-y}$}
\opt{B}{The unit vector in the direction of refracted light at $y=\dfrac R2$ is $\dfrac12\hat i+\dfrac{\sqrt3}{2}\hat j$}
\opt{C}{The maximum angular size of the arc of light is $\theta_{\max}$, the value of $\cos\theta_{\max}$ is $2/5$}
\opt{D}{The maximum angular size of the arc of light is $\theta_{\max}$, the value of $\cos\theta_{\max}$ is $3/5$}
\end{qoptions}
\end{bankq}
\begin{bankq}[id=TOAS04-P2-PHY-S1-Q06, source={TOAS - Test 04 (15 Mar 2022)}, paper=2, qno=6, type=multi, subtopic={Lenses}, answer={A,C,D}, answer_src=key, difficulty=2, diff_parts={K1 S1 I0 C0 F1}, diff_note={concepts: image by a concave lens of a virtual object | steps: virtual object position relative to the focus; image nature and position}]
Consider the rays shown in the diagram as paraxial. The image of the virtual point object $O$ formed by the concave lens $LL'$ is
\bfig{156.5}{TOAS04_p2_phy_s1_q06_1.png}
\begin{qoptions}
\opt{A}{Virtual.}
\opt{B}{Real.}
\opt{C}{Located below the principal axis.}
\opt{D}{Located left of the lens.}
\end{qoptions}
\end{bankq}
\begin{bankpassage}[id=TOAS04-P2-PHY-S3-G5-6, source={TOAS - Test 04 (15 Mar 2022)}, questions={5-6}]
For a glass prism the plot of deviation ($\delta$) v/s angle of incidence ($i$) is as shown.
\bfig{118.0}{TOAS04_p2_phy_s3_para5_6_1.png}
\end{bankpassage}
\begin{bankq}[id=TOAS04-P2-PHY-S3-Q05, source={TOAS - Test 04 (15 Mar 2022)}, paper=2, qno=5, type=numerical, subtopic={Prisms}, answer={1.73}, answer_src=key, difficulty=3, diff_parts={K1 S2 I0 C0 F1}, diff_note={concepts: minimum deviation of a prism | steps: symmetric ray condition; refraction relation; refractive index}, group=TOAS04-P2-PHY-S3-G5-6]
Find the refractive index of the glass.
\end{bankq}
\begin{bankq}[id=TOAS04-P2-PHY-S3-Q06, source={TOAS - Test 04 (15 Mar 2022)}, paper=2, qno=6, type=numerical, subtopic={Prisms}, answer={41.00}, answer_src=key, difficulty=2, diff_parts={K1 S1 I0 C0 F1}, diff_note={concepts: prism refraction | steps: angle at the second face from the emergent condition; Snell's law at the first face}, group=TOAS04-P2-PHY-S3-G5-6]
Find the value of angle $i_1$ (in degree).
\end{bankq}
\begin{bankq}[id=TOAS04-P2-PHY-S4-Q08, source={TOAS - Test 04 (15 Mar 2022)}, paper=2, qno=8, type=numerical, subtopic={Reflection}, answer={12}, answer_src=key, difficulty=4, diff_parts={K1 S2 I1 C0 F1}, diff_note={concepts: images in three inclined plane mirrors | steps: images in each pair of mirrors; multiple reflections; total number of images | I: image counting for a triangle}]
Three plane mirrors are kept as shown in the figure. A point object ($O$) is kept at the centroid of the triangle seen in the figure. How many images will be formed ?
\bfig{103.0}{TOAS04_p2_phy_s4_q08_1.png}
\end{bankq}
\begin{bankq}[id=TOAS04-P2-PHY-S4-Q09, source={TOAS - Test 04 (15 Mar 2022)}, paper=2, qno=9, type=numerical, subtopic={Reflection}, answer={9}, answer_src=key, difficulty=5, diff_parts={K2 S2 I1 C0 F1}, diff_note={concepts: mirror formula, relative motion | steps: image positions of both ends; image velocities including the mirror motion; rate of change of length | I: differentiate with the moving mirror}]
A pencil ($AB$) of length 20 cm is moving along the principal axis of a concave mirror $MM'$, with a velocity 5 m/s approaching the mirror. The mirror itself is moving away from the pencil at a speed of 2 m/s. Radius of curvature of mirror is 40 cm. If the rate of change of length of the image of the pencil at the instant end $A$ is at a distance of 60 cm from the mirror is $\dfrac x4$ m/s, find $x$.
\bfig{158.4}{TOAS04_p2_phy_s4_q09_1.png}
\end{bankq}
\begin{bankq}[id=TOAS05-P1-PHY-S1-Q07, source={TOAS - Test 05 (03 Apr 2022)}, paper=1, qno=7, type=multi, subtopic={Lenses}, answer={B,C}, answer_src=key, difficulty=4, diff_parts={K1 S2 I1 C0 F1}, diff_note={concepts: refraction through a lens with water on one side | steps: radius of curvature from f in air; power with water on one side; image of the object | I: lens with different media on two sides}]
A long cylindrical tube containing water is closed by an equiconvex lens of focal length 10 cm in air. A point object `$O$' is placed along the axis of the tube outside it at a distance of 21 cm from the lens. Refractive index of material of lens is 1.5 and that of water is $\dfrac43$.
\bfig{137.3}{TOAS05_p1_phy_s1_q07_1.png}
\begin{qoptions}
\opt{A}{Final image is virtual in nature.}
\opt{B}{Final image is real in nature.}
\opt{C}{Final image is located at a distance 70 cm from lens inside the tube.}
\opt{D}{Final image is located at a distance 30 cm from lens inside the tube.}
\end{qoptions}
\end{bankq}
\begin{bankq}[id=TOAS05-P1-PHY-S1-Q08, source={TOAS - Test 05 (03 Apr 2022)}, paper=1, qno=8, type=multi, subtopic={Prisms}, answer={A,D}, answer_src=key, difficulty=3, diff_parts={K1 S2 I0 C0 F1}, diff_note={concepts: prism and critical angle | steps: critical angle from the refractive index; limits on r1; minimum incidence angle and maximum deviation}]
A prism with prism angle $60^\circ$ and refractive index $\sqrt{\dfrac73}$ is given.
\bfig{128.8}{TOAS05_p1_phy_s1_q08_1.png}
\begin{qoptions}
\opt{A}{Minimum possible angle of incidence so that the light ray is refracted from the second surface is $30^\circ$.}
\opt{B}{Minimum possible angle of incidence, so that the light ray is refracted from the second surface is $60^\circ$.}
\opt{C}{Maximum possible deviation produced by prism is $30^\circ$.}
\opt{D}{Maximum possible deviation produced by prism is $60^\circ$.}
\end{qoptions}
\end{bankq}
\begin{bankq}[id=TOAS05-P1-PHY-S2-Q11, source={TOAS - Test 05 (03 Apr 2022)}, paper=1, qno=11, type=single, subtopic={Prisms}, answer={D}, answer_src=key, difficulty=4, diff_parts={K2 S2 I1 C0 F0}, diff_note={concepts: prism refraction, total internal reflection | steps: minimum refractive index for TIR of ray 1; condition for ray 2; condition for ray 3 (minimum deviation); match | I: symmetric ray at minimum deviation}]
$ABC$ is a right-angled prism kept in air. A ray (1) is incident on the face $AB$ along the normal. Refractive index of the material of prism is the minimum value that will be required so that ray (1) undergoes total internal reflection at the face $AC$. Another ray (2) is incident on the face $AB$ such that it emerges from face $AC$ along the normal to $AC$. A third ray (3) falls on the face $BC$ and emerges from face $AC$ such that its angle of emergence is the same as that of incidence. Assuming light (1), (2) and (3) have the same wavelength, then match the following.
\bfig{80.8}{TOAS05_p1_phy_s2_q11_1.png}
{\small\setlength{\tabcolsep}{3pt}\renewcommand{\arraystretch}{1.25}\begin{center}\begin{tabular}{|L{0.75cm}|L{3.6cm}|L{0.7cm}|L{1.0cm}|}
\hline
\multicolumn{2}{|c|}{List-I} & \multicolumn{2}{c|}{List-II}\\
\hline
(I) & Refractive index of the material of prism is & (P) & 120\\
\hline
(II) & Angle of incidence in degree of ray (2) is & (Q) & 90\\
\hline
(III) & Deviation in degree suffered by ray (2) is & (R) & 2\\
\hline
(IV) & Deviation in degree suffered by ray (3) is & (S) & 60\\
\hline
\multicolumn{2}{|c|}{} & (T) & 1.5\\
\hline
\multicolumn{2}{|c|}{} & (U) & 4\\
\hline
\end{tabular}\end{center}}
\begin{qoptions}
\opt{A}{I $\to$ U;II $\to$ Q;III $\to$ S;IV $\to$ P}
\opt{B}{I $\to$ R;II $\to$ P;III $\to$ S;IV $\to$ P}
\opt{C}{I $\to$ U;II $\to$ S;III $\to$ P;IV $\to$ T}
\opt{D}{I $\to$ R;II $\to$ Q;III $\to$ S;IV $\to$ P}
\end{qoptions}
\end{bankq}
\begin{bankq}[id=TOAS05-P1-PHY-S2-Q12, source={TOAS - Test 05 (03 Apr 2022)}, paper=1, qno=12, type=single, subtopic={Prisms}, answer={B}, answer_src=key, difficulty=4, diff_parts={K2 S2 I1 C0 F0}, diff_note={concepts: prism refraction, total internal reflection | steps: minimum refractive index for TIR of ray 1; condition for ray 2; condition for ray 3 with the given angles; match | I: symmetric ray at minimum deviation}]
$ABC$ is a right-angled prism kept in air. A ray (1) is incident on the face $AB$ along the normal. Refractive index of the material of prism is the minimum value that will be required so that ray (1) undergoes total internal reflection at the face $AC$. Another ray (2) is incident on the face $AB$ such that it emerges from face $AC$ along the normal to $AC$. A third ray (3) falls on the face $BC$ and emerges from face $AC$ such that its angle of emergence is the same as that of incidence. Assuming light (1), (2) and (3) have the same wavelength, then match the following. [Take : $\sin^{-1}(0.29)\approx17^\circ$, $\sin15^\circ=0.258$]
\bfig{90.8}{TOAS05_p1_phy_s2_q12_1.png}
{\small\setlength{\tabcolsep}{3pt}\renewcommand{\arraystretch}{1.25}\begin{center}\begin{tabular}{|L{0.75cm}|L{3.6cm}|L{0.7cm}|L{1.0cm}|}
\hline
\multicolumn{2}{|c|}{List-I} & \multicolumn{2}{c|}{List-II}\\
\hline
(I) & Refractive index of the material of prism is & (P) & 120\\
\hline
(II) & Angle of incidence in degree of ray (2) is & (Q) & 90\\
\hline
(III) & Deviation in degree suffered by ray (2) is & (R) & 2\\
\hline
(IV) & Deviation in degree suffered by ray (3) is & (S) & 30\\
\hline
\multicolumn{2}{|c|}{} & (T) & 1.15\\
\hline
\multicolumn{2}{|c|}{} & (U) & 4\\
\hline
\end{tabular}\end{center}}
\begin{qoptions}
\opt{A}{I $\to$ U;II $\to$ Q;III $\to$ S;IV $\to$ P}
\opt{B}{I $\to$ T;II $\to$ Q;III $\to$ S;IV $\to$ U}
\opt{C}{I $\to$ U;II $\to$ S;III $\to$ P;IV $\to$ T}
\opt{D}{I $\to$ R;II $\to$ Q;III $\to$ S;IV $\to$ P}
\end{qoptions}
\end{bankq}
\begin{bankq}[id=TOAS05-P1-PHY-S3-Q03, source={TOAS - Test 05 (03 Apr 2022)}, paper=1, qno=3, type=numerical, subtopic={Lenses}, answer={48.00}, answer_src=key, difficulty=6, diff_parts={K2 S3 I1 C0 F1}, diff_note={concepts: lens coated with a layer, silvered | steps: refraction through A and B; reflection at the silvered surface; refraction back; relate object and image distances; solve for R | I: sequential imaging with the mirror}]
A thin plano convex lens $A$ has material of refractive index $\mu_A=1.8$ and it's curved surface has radius of curvature $R$. A thin layer of transparent material $B$ is laid over the curved surface of $A$. The refractive index of $B$ is 1.2 and the curved surface of $B$ that is not touching $A$ has radius of curvature $\dfrac R2$. This surface of $B$ is silvered. A point object `$O$' is kept at a distance of 10 cm on the principal axis of the system (above the plane surface) and it's image is formed at a distance of 40 cm above the plane surface. Find $R$ (in cm).
\bfig{102.7}{TOAS05_p1_phy_s3_q03_1.png}
\end{bankq}
\begin{bankq}[id=TOAS05-P1-PHY-S3-Q04, source={TOAS - Test 05 (03 Apr 2022)}, paper=1, qno=4, type=numerical, subtopic={Lenses}, answer={20.00}, answer_src=key, difficulty=5, diff_parts={K2 S2 I1 C0 F1}, diff_note={concepts: lens, a silvered plano-convex lens | steps: image by the first lens; the silvered lens as a mirror of the combined power; image falls back on the object; solve for d | I: image coincides when the rays retrace}]
A convex lens of focal length 20 cm and another plano-convex lens of focal length 40 cm are placed co-axially. The plane surface of the plano convex lens is silvered. An object $O$ is kept on the principal axis at a distance of 10 cm from the convex lens as shown in figure. Find the distance $d$ (in cm) between two lenses so that final image is formed on object itself.
\bfig{143.2}{TOAS05_p1_phy_s3_q04_1.png}
\end{bankq}
\begin{bankq}[id=TOAS05-P2-PHY-S1-Q06, source={TOAS - Test 05 (03 Apr 2022)}, paper=2, qno=6, type=single, subtopic={Reflection}, answer={B}, answer_src=key, difficulty=1, diff_parts={K1 S1 I0 C0 F0}, diff_note={concepts: image of an extended object on the axis | steps: image of each end by the mirror formula; length of the image}]
A linear object is placed along the axis of a mirror as shown in the figure. If $f$ is the focal length of the mirror then the length of image is
\bfig{104.3}{TOAS05_p2_phy_s1_q06_1.png}
\begin{qoptions}
\opt{A}{$\dfrac{2f}3$}
\opt{B}{$f$}
\opt{C}{$\dfrac f3$}
\opt{D}{$\dfrac{4f}3$}
\end{qoptions}
\end{bankq}
\begin{bankq}[id=TOAS05-P2-PHY-S1-Q13, source={TOAS - Test 05 (03 Apr 2022)}, paper=2, qno=13, type=single, subtopic={Microscope and telescope}, answer={B}, answer_src=key, difficulty=1, diff_parts={K1 S1 I0 C0 F0}, diff_note={concepts: magnifying power of a telescope | steps: M = f0/fe}]
If an astronomical telescope has objective and eye-piece of focal lengths 100 cm and 4 cm respectively, then the magnifying power of the telescope for the normal vision is :
\begin{qoptions}
\opt{A}{42}
\opt{B}{25}
\opt{C}{58}
\opt{D}{104}
\end{qoptions}
\end{bankq}
\begin{bankq}[id=TOAS05-P2-PHY-S1-Q15, source={TOAS - Test 05 (03 Apr 2022)}, paper=2, qno=15, type=single, subtopic={Lenses}, answer={B}, answer_src=key, difficulty=1, diff_parts={K1 S1 I0 C0 F0}, diff_note={concepts: two lenses giving a parallel beam | steps: d = f1 + f2}]
A convex lens $A$ of focal length 20 and a concave lens $B$ of focal length 5 cm are kept along the same axis with the distance $d$ between them. If a parallel beam of light falling on $A$ leaves $B$ as a parallel beam, then distance $d$ in cm will be :-
\begin{qoptions}
\opt{A}{25}
\opt{B}{15}
\opt{C}{30}
\opt{D}{50}
\end{qoptions}
\end{bankq}
\begin{bankq}[id=TOAS05-P2-PHY-S1-Q17, source={TOAS - Test 05 (03 Apr 2022)}, paper=2, qno=17, type=single, subtopic={Prisms}, answer={A}, answer_src=key, difficulty=1, diff_parts={K1 S1 I0 C0 F0}, diff_note={concepts: deviation by thin prisms | steps: deviation (n1 - n2) theta from the pair; solve}]
Two identical prisms with slightly different indices are located as shown. Angle $\theta$ is small. When a laser beam strikes one of the prisms perpendicular to the surface, the refracted ray is deviated by a small angle $\phi$. The difference between the indices of refraction of the prism in terms of $\theta$ and $\phi$ is :-
\bfig{102.4}{TOAS05_p2_phy_s1_q17_1.png}
\begin{qoptions}
\opt{A}{$\dfrac\phi\theta$}
\opt{B}{$\dfrac\theta\phi$}
\opt{C}{$\dfrac{2\theta}\phi$}
\opt{D}{$\dfrac{2\phi}\theta$}
\end{qoptions}
\end{bankq}
\begin{bankq}[id=TOAS05-P2-PHY-S1-Q18, source={TOAS - Test 05 (03 Apr 2022)}, paper=2, qno=18, type=single, subtopic={Prisms}, answer={B}, answer_src=key, difficulty=2, diff_parts={K1 S1 I1 C0 F0}, diff_note={concepts: deviation versus incidence graph | steps: minimum deviation point of the graph; relation to the prism angle | I: read the minimum deviation condition}]
Figure shows a graph of angle of deviation $\delta$ versus angle of incidence $i$ for a light ray striking a prism. Angle of prism is
\bfig{93.0}{TOAS05_p2_phy_s1_q18_1.png}
\begin{qoptions}
\opt{A}{$30^\circ$}
\opt{B}{$45^\circ$}
\opt{C}{$60^\circ$}
\opt{D}{$75^\circ$}
\end{qoptions}
\end{bankq}
\begin{bankq}[id=TOAS05-P2-PHY-S1-Q19, source={TOAS - Test 05 (03 Apr 2022)}, paper=2, qno=19, type=single, subtopic={Refraction}, answer={B}, answer_src=key, difficulty=1, diff_parts={K1 S1 I0 C0 F0}, diff_note={concepts: refraction in a sphere | steps: Snell's law gives r = 30; deviation 2(i - r)}]
A ray of light is incident on a glass sphere of refractive index $\sqrt3$ at an angle of incidence $60^\circ$ as shown in figure. The total deviation suffered by the ray is
\bfig{103.1}{TOAS05_p2_phy_s1_q19_1.png}
\begin{qoptions}
\opt{A}{$30^\circ$}
\opt{B}{$60^\circ$}
\opt{C}{$90^\circ$}
\opt{D}{$120^\circ$}
\end{qoptions}
\end{bankq}
\begin{bankq}[id=TOAS05-P2-PHY-S1-Q20, source={TOAS - Test 05 (03 Apr 2022)}, paper=2, qno=20, type=single, subtopic={Microscope and telescope}, answer={B}, answer_src=key, difficulty=1, diff_parts={K1 S1 I0 C0 F0}, diff_note={concepts: telescope in normal adjustment | steps: M = fo/fe gives fe = 2 cm; tube length fo + fe}]
In an astronomical telescope, objective is of focal length 20 cm. What should be tube length for magnifying power of 10, under normal magnification ?
\begin{qoptions}
\opt{A}{10 cm}
\opt{B}{22 cm}
\opt{C}{24 cm}
\opt{D}{30 cm}
\end{qoptions}
\end{bankq}
\begin{bankq}[id=TOAS05-P2-PHY-S2-Q02, source={TOAS - Test 05 (03 Apr 2022)}, paper=2, qno=2, type=numerical, subtopic={Lenses}, answer={48}, answer_src=key, difficulty=3, diff_parts={K2 S1 I0 C0 F1}, diff_note={concepts: lens formula, sequential imaging | steps: image by the first element; image by the second element}]
A point object is placed at distance of 20 cm from a lens as shown in figure. The distance of final image from the lens (in cm) is \underline{\hspace{1cm}}.
\bfig{147.0}{TOAS05_p2_phy_s2_q02_1.png}
\end{bankq}
\begin{bankq}[id=TOAS05-P2-PHY-S2-Q05, source={TOAS - Test 05 (03 Apr 2022)}, paper=2, qno=5, type=numerical, subtopic={Refraction}, answer={4}, answer_src=key, difficulty=3, diff_parts={K1 S1 I1 C0 F1}, diff_note={concepts: refraction at a curved surface | steps: spherical refraction relation for the ray; geometry fixes m | I: parallel emergent ray condition}]
A quarter cylinder of glass of radius $R$ and refractive index 1.5 is placed on a table. A point object $P$ is kept at a distance of $mR$ from it. Find the value of $3m$ for which a ray from $P$ will emerge parallel to the table as shown in the figure.
\bfig{101.1}{TOAS05_p2_phy_s2_q05_1.png}
\end{bankq}
\begin{bankq}[id=TOAS05-P2-PHY-S2-Q06, source={TOAS - Test 05 (03 Apr 2022)}, paper=2, qno=6, type=numerical, subtopic={Lenses}, answer={15}, answer_src=key, difficulty=4, diff_parts={K2 S2 I0 C0 F1}, diff_note={concepts: lens maker formula, mirror formula | steps: focal length of the lens in air; equivalent mirror focal length; object at the centre of curvature of the equivalent mirror}]
The convex surface of a thin concavo-convex lens of glass of refractive index 1.5 has a radius of curvature 20 cm. The concave surface has a radius of curvature 60 cm. The convex side is silvered and placed on a horizontal surface. At what distance (in cm) should a pin be placed on optical axis so that image is formed at same place.
\bfig{72.8}{TOAS05_p2_phy_s2_q06_1.png}
\end{bankq}
\begin{bankq}[id=TOAS05-P2-PHY-S2-Q08, source={TOAS - Test 05 (03 Apr 2022)}, paper=2, qno=8, type=numerical, subtopic={Lenses}, answer={40}, answer_src=key, difficulty=2, diff_parts={K1 S1 I0 C0 F1}, diff_note={concepts: real image with magnification 2 | steps: u + v = 60 and v = 2u; lens formula}]
The distance between the object and its real image from the convex lens is 60 cm and the height of image is two times the height of object. The focal length of the lens (in cm) is $f$. Then $3f$ will be :
\end{bankq}
\begin{bankq}[id=TOAS06-P1-PHY-S2-Q05, source={TOAS - Test 06 (17 Apr 2022)}, paper=1, qno=5, type=multi, subtopic={Lenses}, answer={A,B,C}, answer_src=key, difficulty=3, diff_parts={K1 S2 I0 C0 F1}, diff_note={concepts: thick lens relations | steps: Gaussian form of the lens equation; magnification relation; ratio of focal lengths in different media}]
Which of the following relations for a coaxial lens-system separating two media of refractive indices $\mu_1$ and $\mu_2$ are true ? (Here $u$ and $v$ are the distances of the object and the image from the first and the second principal points respectively, $f_1$ and $f_2$ the focal lengths of the system and $m$ is linear transverse magnification.)
\begin{qoptions}
\opt{A}{$\dfrac{f_1}{u}+\dfrac{f_2}{v}=1$}
\opt{B}{$m=-\dfrac{v}{u}\left(\dfrac{f_1}{f_2}\right)$}
\opt{C}{$\dfrac{f_1}{f_2}=-\dfrac{\mu_1}{\mu_2}$}
\opt{D}{$\dfrac{f_1}{u}-\dfrac{f_2}{v}=1$}
\end{qoptions}
\end{bankq}
\begin{bankq}[id=TOAS06-P1-PHY-S2-Q09, source={TOAS - Test 06 (17 Apr 2022)}, paper=1, qno=9, type=multi, subtopic={Lenses}, answer={C,D}, answer_src=key, difficulty=4, diff_parts={K1 S2 I1 C0 F1}, diff_note={concepts: imaging with a reflecting tube | steps: image by the lens of the source; images in the mirror walls; pattern of images on the screen | I: reflected images lie on a circle}]
A tube with reflecting inner walls has a point light source at the bottom of it (see figure). The inner diameter of the tube is $d=12$ mm, the tube's length is $\ell=60$ mm. A convex lens with a focal length $F=36$ mm is placed against the open end of the tube and at a distance $L=90$ mm from the tip of the tube there is a screen, which has a graph paper attached to it. The intersection point with the optical axis O is marked on the paper.
\bfig{212.0}{TOAS06_p1_phy_s2_q09_1.png}
\begin{qoptions}
\opt{A}{A point image is formed on the screen at point O.}
\opt{B}{The pattern of images obtained on the screen are parallel straight lines.}
\opt{C}{The pattern of images obtained on the screen are concentric circles.}
\opt{D}{The radius of concentric circles formed on the screen are $(18n)$ mm, where $n$ is an integer.}
\end{qoptions}
\end{bankq}
\begin{bankq}[id=TOAS06-P1-PHY-S3-Q01, source={TOAS - Test 06 (17 Apr 2022)}, paper=1, qno=1, type=numerical, subtopic={Lenses}, answer={24.00}, answer_src=key, difficulty=4, diff_parts={K1 S2 I1 C0 F1}, diff_note={concepts: lens equation with two object positions | steps: image distance for the first source; the other source must give the same image; separation | I: coincident images need opposite object sides}]
Two point light sources are positioned on different points of a convex lens' optical axis. The images of these light sources formed by the lens coincide. It is known that one light source is at a distance $a=18$ cm from the center of the lens. How far (in cm) is the other light source distanced from this light source? The focal length of the lens is $f=9$ cm.
\end{bankq}
\begin{bankq}[id=TOAS06-P1-PHY-S3-Q05, source={TOAS - Test 06 (17 Apr 2022)}, paper=1, qno=5, type=numerical, subtopic={Lenses}, answer={8.00}, answer_src=key, difficulty=6, diff_parts={K2 S2 I2 C0 F1}, diff_note={concepts: real image positions of a lens, SHM | steps: lens positions giving real images (quadratic); equal time interval condition fixes the amplitude; omega; impulse | I: four real images at equal time intervals}]
A point object is located at a distance of 100 cm from a screen. A lens of focal length 23 cm mounted on a movable frictionless stand is kept between the source and the screen. The stand is attached to a spring of natural length 50 cm and spring constant 800 N/m as shown. Mass of the stand with lens is 2 kg. How much minimum impulse $P$ (in kg m/s) should be imparted to the stand so that a real image of the object is formed on the screen four times during a complete oscillation. All these are separated by same time interval. (Neglect width of the stand.)
\bfig{141.6}{TOAS06_p1_phy_s3_q05_1.png}
\end{bankq}
\begin{bankq}[id=TOAS07-P1-PHY-S1-Q04, source={TOAS - Test 07 (01 May 2022)}, paper=1, qno=4, type=single, subtopic={Lenses}, answer={C}, answer_src=key, difficulty=2, diff_parts={K1 S2 I0 C0 F0}, diff_note={concepts: lens combination | steps: focal length of the convex lens from 15 and 30; combined focal length from the shifted image; focal length of the concave lens}]
A point object is placed at a distance of 15 cm from a convex lens. The image is formed on the other side of lens at a distance 30 cm from lens. When a concave lens is placed in contact with convex lens, image is shifted further by 30 cm. Calculate the focal length of concave lens :-
\begin{qoptions}
\opt{A}{10 cm}
\opt{B}{$-10$ cm}
\opt{C}{$-60$ cm}
\opt{D}{60 cm}
\end{qoptions}
\end{bankq}
\begin{bankpassage}[id=TOAS07-P1-PHY-S3-G5-6, source={TOAS - Test 07 (01 May 2022)}, questions={5-6}]
In figure $L$ is half part of an equiconvex lens of refractive index 1.5 whose surfaces have radius of curvature 40 cm and its right surface is silvered as shown in figure. Normal to its principal axis, a plane mirror $M$ is placed on right of lens. Distance between lens $L$ and mirror $M$ is $b$. A small object $O$ is placed on left of the lens. Images found by lens and mirror coincide with each other and tranverse length of final image formed by the lens is twice that of image formed by the mirror.
\bfig{103.8}{TOAS07_p1_phy_s3_para5_6_1.png}
\end{bankpassage}
\begin{bankq}[id=TOAS07-P1-PHY-S3-Q05, source={TOAS - Test 07 (01 May 2022)}, paper=1, qno=5, type=numerical, subtopic={Lenses}, answer={5.00}, answer_src=key, difficulty=3, diff_parts={K1 S2 I0 C0 F1}, diff_note={concepts: lens problem from the passage | steps: image formed by the first lens; second lens; value of a}, group=TOAS07-P1-PHY-S3-G5-6]
The value of $a$ (in cm) is
\end{bankq}
\begin{bankq}[id=TOAS07-P1-PHY-S3-Q06, source={TOAS - Test 07 (01 May 2022)}, paper=1, qno=6, type=numerical, subtopic={Lenses}, answer={2.50}, answer_src=key, difficulty=3, diff_parts={K1 S2 I0 C0 F1}, diff_note={concepts: lens problem from the passage | steps: image formed by the first lens; second lens; value of b}, group=TOAS07-P1-PHY-S3-G5-6]
The value of $b$ (in cm) is
\end{bankq}
\begin{bankq}[id=TOAS07-P1-PHY-S4-Q07, source={TOAS - Test 07 (01 May 2022)}, paper=1, qno=7, type=numerical, subtopic={Microscope and telescope}, answer={94}, answer_src=key, difficulty=4, diff_parts={K1 S2 I0 C1 F1}, diff_note={concepts: compound microscope | steps: eyepiece object distance for image at 25 cm; objective image distance = L - ue; objective object distance; product of magnifications}]
In a compound microscope the objective and eye piece have focal lengths of 0.95 cm and 5 cm respectively, and are kept at a distance of 20 cm. The last image is formed at a distance of 25 cm from eye piece. Calculate the magnitude of total magnification of the compound microscope.
\end{bankq}
\begin{bankq}[id=TOAS07-P1-PHY-S4-Q08, source={TOAS - Test 07 (01 May 2022)}, paper=1, qno=8, type=numerical, subtopic={Refraction}, answer={5}, answer_src=key, difficulty=4, diff_parts={K2 S2 I0 C0 F1}, diff_note={concepts: lens with water, silvered surface | steps: refraction at the water-glass system; the silvered face behaves as a mirror; image coincides with the object when the ray retraces; solve for R}]
Bottom of a glass beaker is made of a thin equi-convex lens having bottom side silvered as shown in figure. Water is filled in the beaker up to a height 4m.The image of a point object $O$, floating at middle point of beaker at the surface of water coincides with it. Find out the radius (in cm) of curvature of the lens. Given that refractive index of glass is $\dfrac32$ and that of water is $\dfrac43$.
\bfig{107.8}{TOAS07_p1_phy_s4_q08_1.png}
\end{bankq}
\begin{bankq}[id=TOAS07-P2-PHY-S1-Q05, source={TOAS - Test 07 (01 May 2022)}, paper=2, qno=5, type=multi, subtopic={Refraction}, answer={A,C}, answer_src=key, difficulty=3, diff_parts={K1 S1 I1 C0 F1}, diff_note={concepts: image formation at a spherical interface | steps: sign of the power for each pair of indices; real or virtual outcomes for real and virtual objects | I: sign analysis of the refraction formula}]
Two refracting media are separated by a spherical interface as shown in the figure. $PP'$ is the principal axis, $\mu_1$ and $\mu_2$ are the refractive indices of medium of incidence and medium of refraction respectively. Then :
\bfig{94.6}{TOAS07_p2_phy_s1_q05_1.png}
\begin{qoptions}
\opt{A}{if $\mu_2=1.5$ and $\mu_1=1.2$, then there cannot be a real image of real object.}
\opt{B}{if $\mu_2=1.5$ and $\mu_1=1.2$, then there cannot be a real image of virtual object.}
\opt{C}{if $\mu_2=1.2$ and $\mu_1=1.5$, then there cannot be a virtual image of virtual object.}
\opt{D}{if $\mu_2=1.2$ and $\mu_1=1.5$, then there cannot be a real image of real object.}
\end{qoptions}
\end{bankq}
\begin{bankq}[id=TOAS07-P2-PHY-S1-Q06, source={TOAS - Test 07 (01 May 2022)}, paper=2, qno=6, type=multi, subtopic={Reflection}, answer={A}, answer_src=key, difficulty=3, diff_parts={K1 S1 I1 C0 F1}, diff_note={concepts: image by a plane mirror | steps: a real image needs converging incident rays | I: virtual object}]
The light on reflection from a plane mirror can give a real image when :
\begin{qoptions}
\opt{A}{the convergent rays are incident on the mirror}
\opt{B}{the divergent rays are incident on the mirror}
\opt{C}{an object is placed very close to the mirror}
\opt{D}{an object is placed very far away from the mirror}
\end{qoptions}
\end{bankq}
\begin{bankpassage}[id=TOAS07-P2-PHY-S2-G7-8, source={TOAS - Test 07 (01 May 2022)}, questions={7-8}]
Two identical plane mirrors AB and CD of length $L$ each are placed parallel and a point source of light $O$ is placed midway between them as shown. All the shown dashed segments are either parallel or perpendicular to the mirrors. Answer the following questions based on this information.
\bfig{128.3}{TOAS07_p2_phy_s2_para7_8_1.png}
\end{bankpassage}
\begin{bankq}[id=TOAS07-P2-PHY-S2-Q07, source={TOAS - Test 07 (01 May 2022)}, paper=2, qno=7, type=single, subtopic={Reflection}, answer={B}, answer_src=key, difficulty=1, diff_parts={K1 S1 I0 C0 F0}, diff_note={concepts: reflection from two mirrors | steps: ray coverage between the mirrors; compare AB and CD}, group=TOAS07-P2-PHY-S2-G7-8]
Choose the correct statement :
\begin{qoptions}
\opt{A}{All the ray's incident on mirror AB directly from source $O$, after reflection will be incident on mirror CD}
\opt{B}{Some of the ray's incident on mirror AB directly from source $O$, after reflection will be incident on mirror CD}
\opt{C}{All the ray's incident on mirror CD directly from source $O$, after reflection will be incident on mirror AB}
\opt{D}{Some of the ray's incident on mirror CD directly from source $O$, after reflection will be incident on mirror AB}
\end{qoptions}
\end{bankq}
\begin{bankq}[id=TOAS07-P2-PHY-S2-Q08, source={TOAS - Test 07 (01 May 2022)}, paper=2, qno=8, type=single, subtopic={Reflection}, answer={B}, answer_src=key, difficulty=1, diff_parts={K1 S1 I0 C0 F0}, diff_note={concepts: number of images in inclined mirrors | steps: angle between the mirrors; count}, group=TOAS07-P2-PHY-S2-G7-8]
The total number of images of the object as formed by both the mirrors is :
\begin{qoptions}
\opt{A}{2}
\opt{B}{3}
\opt{C}{4}
\opt{D}{infinite}
\end{qoptions}
\end{bankq}
\begin{bankpassage}[id=TOAS07-P2-PHY-S2-G9-10, source={TOAS - Test 07 (01 May 2022)}, questions={9-10}]
A rectangular wire frame is kept on the principal axis of a concave mirror as shown in figure. $F$ \& $C$ are the focus \& centre of curvature of mirror respectively
\bfig{144.0}{TOAS07_p2_phy_s2_para9_10_1.png}
$f$ is the focal length of mirror. A small particle is moving along the length of wire frame with constant velocity. Based on above information answer the following questions.
\end{bankpassage}
\begin{bankq}[id=TOAS07-P2-PHY-S2-Q09, source={TOAS - Test 07 (01 May 2022)}, paper=2, qno=9, type=single, subtopic={Reflection}, answer={A}, answer_src=key, difficulty=1, diff_parts={K1 S1 I0 C0 F0}, diff_note={concepts: images in two mirrors | steps: image by each mirror; shape of the combined image}, group=TOAS07-P2-PHY-S2-G9-10]
The image of the given object is best shown by
\begin{qoptions}
\opt{A}{\bfigi{107.0}{TOAS07_p2_phy_s2_q09_A.png}}
\opt{B}{\bfigi{104.4}{TOAS07_p2_phy_s2_q09_B.png}}
\opt{C}{\bfigi{105.5}{TOAS07_p2_phy_s2_q09_C.png}}
\opt{D}{\bfigi{105.8}{TOAS07_p2_phy_s2_q09_D.png}}
\end{qoptions}
\end{bankq}
\begin{bankq}[id=TOAS07-P2-PHY-S2-Q10, source={TOAS - Test 07 (01 May 2022)}, paper=2, qno=10, type=single, subtopic={Reflection}, answer={D}, answer_src=key, difficulty=3, diff_parts={K1 S2 I0 C1 F0}, diff_note={concepts: image of a wire frame in a mirror | steps: mirror formula for each end; velocity of the image; total length of the image}, group=TOAS07-P2-PHY-S2-G9-10]
Mark the correct statement(s) :
\begin{qoptions}
\opt{A}{The velocity of image of particle is remaining constant}
\opt{B}{The velocity of image of particle is varying}
\opt{C}{The total length of image of wire frame is $\left(\dfrac{7+\sqrt{17}}{4}\right)f$}
\opt{D}{Both (B) and (C) are correct}
\end{qoptions}
\end{bankq}
\begin{bankq}[id=TOAS08-P1-PHY-S1-Q03, source={TOAS - Test 08 (15 May 2022)}, paper=1, qno=3, type=multi, subtopic={Refraction}, answer={A,B,C}, answer_src=key, difficulty=4, diff_parts={K1 S2 I1 C0 F1}, diff_note={concepts: refraction at a plane surface and apparent position | steps: Snell's law for the ray along 3i + 4j; apparent image on the line through the source and observer; coordinates of the image; ray along 4i + 3j | I: image lies on the line along the refracted ray}]
A point source $S$ is placed in water at origin $(0,0)$ and an observer $O$ is in air at $(7\text{m},7\text{m})$ as shown in figure. (Given : $\mu_{\text{Water}}=4/3$)
\bfig{122.0}{TOAS08_p1_phy_s1_q03_1.png}
\begin{qoptions}
\opt{A}{Ray starting from source along $3\hat i+4\hat j$ vector will reach observer after refraction.}
\opt{B}{Image of object seen by observer will be at line $4y-3x=7$.}
\opt{C}{Co-ordinate of image seen by observer is $\left(\dfrac{21}{16}\text{m},\dfrac{175}{64}\text{m}\right)$.}
\opt{D}{Ray starting from source along $4\hat i+3\hat j$ vector will reach observer after refraction.}
\end{qoptions}
\end{bankq}
\begin{bankpassage}[id=TOAS08-P1-PHY-S2-G9-10, source={TOAS - Test 08 (15 May 2022)}, questions={9-10}]
A thick plano-convex lens is made of material having refractive index $\mu=\dfrac32$. When it's plane surface is slivered and an object is placed 60 cm away from it's convex surface then final image forms at object it self. While when it's convex surface is silvered and object is placed at distance $\dfrac{20}{3}$cm away from plane surface then final image again forms at object itself.
\bfig{112.3}{TOAS08_p1_phy_s2_para9_10_1.png}
\end{bankpassage}
\begin{bankq}[id=TOAS08-P1-PHY-S2-Q09, source={TOAS - Test 08 (15 May 2022)}, paper=1, qno=9, type=single, subtopic={Lenses}, answer={D}, answer_src=key, difficulty=1, diff_parts={K1 S1 I0 C0 F0}, diff_note={concepts: lens maker formula from the passage | steps: data from the passage; radius of curvature}, group=TOAS08-P1-PHY-S2-G9-10]
Radius of curvature of convex surface is :
\begin{qoptions}
\opt{A}{10 cm}
\opt{B}{15 cm}
\opt{C}{25 cm}
\opt{D}{30 cm}
\end{qoptions}
\end{bankq}
\begin{bankq}[id=TOAS08-P1-PHY-S2-Q10, source={TOAS - Test 08 (15 May 2022)}, paper=1, qno=10, type=single, subtopic={Lenses}, answer={C}, answer_src=key, difficulty=1, diff_parts={K1 S1 I0 C0 F0}, diff_note={concepts: lens thickness from the passage | steps: data from the passage; thickness}, group=TOAS08-P1-PHY-S2-G9-10]
Thickness of lens is
\begin{qoptions}
\opt{A}{10 cm}
\opt{B}{15 cm}
\opt{C}{20 cm}
\opt{D}{25 cm}
\end{qoptions}
\end{bankq}
\begin{bankq}[id=TOAS08-P1-PHY-S3-Q03, source={TOAS - Test 08 (15 May 2022)}, paper=1, qno=3, type=numerical, subtopic={Reflection}, answer={21.00}, answer_src=key, difficulty=4, diff_parts={K2 S2 I0 C0 F1}, diff_note={concepts: reflection in a plane mirror, refraction in a lens | steps: image by the inclined mirror; position relative to the lens; image by the lens}]
A plane mirror $M_1$ is placed at angle $30^\circ$ with $x$-axis and a lens $L_1$ of focal length 15 cm is placed 30 cm away from it as shown in figure. A source $O$ is placed at $[0,\,10\sqrt3\text{cm}]$. If co-ordinate of image formed after reflection from plane mirror then refraction from lens is $(x,y)$ then value of $\sqrt3\dfrac xy$ is.
\bfig{185.4}{TOAS08_p1_phy_s3_q03_1.png}
\end{bankq}
\begin{bankq}[id=TOAS08-P1-PHY-S3-Q08, source={TOAS - Test 08 (15 May 2022)}, paper=1, qno=8, type=numerical, subtopic={Refraction}, answer={40.00}, answer_src=key, difficulty=3, diff_parts={K1 S2 I0 C0 F1}, diff_note={concepts: refraction through a block | steps: refraction at the first face; critical angle check at the next face; deviation between incident and emergent rays}]
A monochromatic light ray incident on block $ABCD$ having R.I $\mu=\dfrac{2}{\sqrt3}$ at angle of incidence $i=\sin^{-1}\left[\dfrac{1}{\sqrt3}\right]$ as shown in figure. Find angle between incident and emergent ray (in degree)
\bfig{114.6}{TOAS08_p1_phy_s3_q08_1.png}
\end{bankq}
\begin{bankq}[id=TOAS08-P2-PHY-S1-Q03, source={TOAS - Test 08 (15 May 2022)}, paper=2, qno=3, type=multi, subtopic={Prisms}, answer={A,B,C}, answer_src=key, difficulty=4, diff_parts={K1 S2 I1 C0 F1}, diff_note={concepts: deviation graph of a prism | steps: minimum deviation from the graph gives mu; maximum deviation at grazing emergence; deviation at normal incidence | I: read the minimum deviation condition}]
For an equilateral glass prism graph of deviation $\delta$ versus angle of incidence $i$ is shown in figure. Choose incorrect option(s) :-
\bfig{104.1}{TOAS08_p2_phy_s1_q03_1.png}
\begin{qoptions}
\opt{A}{Maximum possible deviation produced by prism is $60^\circ+\sin^{-1}\left(\sqrt{\dfrac32}-\dfrac12\right)$}
\opt{B}{Maximum possible deviation produced by prism is $30^\circ+\sin^{-1}\left(\sqrt{\dfrac32}+\dfrac12\right)$.}
\opt{C}{Refraction index of prism is $\sqrt{\dfrac32}$.}
\opt{D}{For normal incidence on prism total deviation suffered by ray is $60^\circ$.}
\end{qoptions}
\end{bankq}
\begin{bankq}[id=TOAS08-P2-PHY-S1-Q05, source={TOAS - Test 08 (15 May 2022)}, paper=2, qno=5, type=multi, subtopic={Reflection}, answer={A,C}, answer_src=key, difficulty=3, diff_parts={K1 S2 I0 C0 F1}, diff_note={concepts: spherical mirror with a virtual object | steps: magnification fixes the image position; mirror formula; type of mirror and focal length}]
An image of a virtual object is formed by a spherical mirror and height of image $I$ is half of height of object $O$ from principal axis. Select correct options :-
\bfig{133.2}{TOAS08_p2_phy_s1_q05_1.png}
\begin{qoptions}
\opt{A}{Mirror is on the right side of $I$ and it is convex.}
\opt{B}{Focal length of mirror is 20 cm.}
\opt{C}{Focal length of mirror is $\dfrac{80}{3}$ cm.}
\opt{D}{Mirror is concave and on right side of object.}
\end{qoptions}
\end{bankq}
\begin{bankq}[id=TOAS08-P2-PHY-S2-Q10, source={TOAS - Test 08 (15 May 2022)}, paper=2, qno=10, type=single, subtopic={Refraction}, answer={D}, answer_src=key, difficulty=2, diff_parts={K1 S2 I0 C0 F0}, diff_note={concepts: refraction at a water surface | steps: critical and refraction angles; apparent depth and displacement; line of sight}]
There is a point source of light at a depth 8m below the surface of water. An observer in air sees the image of the source. Line of sight of observer makes an angle $53^\circ$ with the normal to surface. All quantities in List--II are in SI unit. (use : $\sqrt7=2.65$, $\mu$ of water $=4/3$) ?
\begin{center}\small\renewcommand{\arraystretch}{1.1}\setlength{\extrarowheight}{3.5pt}\setlength{\tabcolsep}{4pt}
\begin{tabular}{|c|l|c|c|}
\hline
 & \multicolumn{1}{c|}{List-I} & & List-II\\[3pt]
\hline
(P) & \parbox[t]{4.2cm}{Apparent depth of image from surface.} & (1) & 9.05\\[3pt]
(Q) & \parbox[t]{4.2cm}{Horizontal distance between object and its image.} & (2) & 12.02\\[3pt]
(R) & \parbox[t]{4.2cm}{Radius of circle of illuminance.} & (3) & 3.38\\[3pt]
(S) & \parbox[t]{4.2cm}{Percentage of power coming out of the water assuming no absorption of light by water and 100\% refraction within circle of illuminance.} & (4) & 16.34\\[3pt]
 & & (5) & 2.53\\[3pt]
\hline
\end{tabular}
\end{center}
\begin{qoptions}
\opt{A}{$P\to3;Q\to1;R\to5;S\to2$}
\opt{B}{$P\to2;Q\to1;R\to4;S\to3$}
\opt{C}{$P\to4;Q\to5;R\to1;S\to4$}
\opt{D}{$P\to5;Q\to3;R\to1;S\to4$}
\end{qoptions}
\end{bankq}
\begin{bankq}[id=TOAS08-P2-PHY-S3-Q01, source={TOAS - Test 08 (15 May 2022)}, paper=2, qno=1, type=numerical, subtopic={Refraction}, answer={5.60}, answer_src=key, difficulty=6, diff_parts={K1 S2 I2 C1 F1}, diff_note={concepts: refraction in a gradient index slab | steps: Snell's invariant for grazing incidence; relation between dy/dx and mu(x); integrate to reach the upper surface | I: invariant and path integral}]
A light ray incident at graving incidence on a long rectangular slab having thickness $d$ at origin. Determine $x$ coordinate of point $A$ where ray meets the upper surface of slab air interface. $\mu$ of slab varies as $\mu=\sqrt{1+e^{x/d}}$; $d=4$m and take $\ln2=0.7$
\bfig{114.4}{TOAS08_p2_phy_s3_q01_1.png}
\end{bankq}
\begin{bankq}[id=TOAS08-P2-PHY-S3-Q03, source={TOAS - Test 08 (15 May 2022)}, paper=2, qno=3, type=numerical, subtopic={Lenses}, answer={13.50}, answer_src=key, difficulty=4, diff_parts={K2 S2 I0 C0 F1}, diff_note={concepts: lens with water, a silvered surface | steps: focal length of the lens with water on top; mirror action of the silvered surface; image coincides with the object}]
Bottom of a container is made of an equi-convex lens having focal length 60 cm in air and $\mu=1.5$. At what height (cm) above water surface an object should be placed so that final image of object coincides with object. Lower surface of lens is silvered. ($\mu_{\text{water}}=4/3$)
\bfig{81.8}{TOAS08_p2_phy_s3_q03_1.png}
\end{bankq}
\begin{bankq}[id=TOAS09-P1-PHY-S2-Q10, source={TOAS - Test 09 (30 May 2022)}, paper=1, qno=10, type=multi, subtopic={Refraction}, answer={B,D}, answer_src=key, difficulty=4, diff_parts={K1 S2 I1 C0 F1}, diff_note={concepts: apparent velocity of an object seen through a liquid surface | steps: apparent depth rate for a moving surface; speed seen directly; speed seen after reflection from the silvered bottom | I: image in the mirror at the same depth}]
A bird is flying downward with constant velocity 6 cm/s towards a fixed container filled with a liquid of refractive index $\dfrac43$. A fish is moving upward with constant speed 4 cm/s towards the surface of liquid. Surface of liquid is also moving downward with constant velocity 2 cm/s shown in the figure. The bottom of container is silvered. Then :-
\bfig{76.2}{TOAS09_p1_phy_s2_q10_1.png}
\begin{qoptions}
\opt{A}{Speed of fish as seen by bird only after refraction is 9 cm/sec.}
\opt{B}{Speed of fish as seen by bird only after refraction is 8.5 cm/sec.}
\opt{C}{Speed of fish as seen by bird after reflection and then refraction is 3 cm/sec.}
\opt{D}{Speed of fish as seen by bird after reflection and then refraction is 2.5 cm/sec.}
\end{qoptions}
\end{bankq}
\begin{bankpassage}[id=TOAS09-P1-PHY-S3-G3-4, source={TOAS - Test 09 (30 May 2022)}, questions={3-4}]
A converging lens of power $P_1=+5$D and a concave mirror of power $P_2=+16$D are placed at a separation 80 cm such that their principal axis coincide. A point object is placed left side of the lens on its principal axes at a distance $X_1$ from the lens such that image of the object coincides with the object itself. When the mirror is further shifted right side by a distance $X_2$, the image of the object further coincides with the object.
\bfig{95.5}{TOAS09_p1_phy_s3_para3_4_1.png}
\end{bankpassage}
\begin{bankq}[id=TOAS09-P1-PHY-S3-Q03, source={TOAS - Test 09 (30 May 2022)}, paper=1, qno=3, type=numerical, subtopic={Lenses}, answer={26.67}, answer_src=key, difficulty=2, diff_parts={K1 S1 I0 C0 F1}, diff_note={concepts: lens problem from the passage | steps: image positions from the lens equation; value of X1}, group=TOAS09-P1-PHY-S3-G3-4]
The value of $X_1$ (in cm) is :
\end{bankq}
\begin{bankq}[id=TOAS09-P1-PHY-S3-Q04, source={TOAS - Test 09 (30 May 2022)}, paper=1, qno=4, type=numerical, subtopic={Lenses}, answer={12.50}, answer_src=key, difficulty=2, diff_parts={K1 S1 I0 C0 F1}, diff_note={concepts: lens problem from the passage | steps: image positions from the lens equation; value of X2}, group=TOAS09-P1-PHY-S3-G3-4]
The value of $X_2$ (in cm) is:
\end{bankq}
\begin{bankq}[id=TOAS09-P1-PHY-S4-Q09, source={TOAS - Test 09 (30 May 2022)}, paper=1, qno=9, type=numerical, subtopic={Refraction}, answer={6}, answer_src=key, difficulty=2, diff_parts={K1 S1 I0 C0 F1}, diff_note={concepts: apparent velocity through parallel slabs | steps: apparent shift through each slab; sum the velocity contributions}]
Find the velocity (in cm/s) of the car C as appeared to the observer shown in the figure. There are several parallel slabs of different medium placed between the car and the observer.
\bfig{337.2}{TOAS09_p1_phy_s4_q09_1.png}
\end{bankq}
\begin{bankq}[id=TOAS09-P2-PHY-S1-Q20, source={TOAS - Test 09 (30 May 2022)}, paper=2, qno=20, type=single, subtopic={Refraction}, answer={D}, answer_src=key, difficulty=2, diff_parts={K1 S1 I1 C0 F0}, diff_note={concepts: deviation at a water-air interface | steps: deviation before and after the critical angle; behaviour at total internal reflection | I: graph with a break at the critical angle}]
A ray of light incidents at an angle of incidence $\theta$ in water at water-air interface. Graph between angle of deviation $(\delta)$ and angle of incidence $(\theta)$ will be-
\bfig{110.0}{TOAS09_p2_phy_s1_q20_1.png}
\begin{qoptions}
\opt{A}{\bfigi{76.1}{TOAS09_p2_phy_s1_q20_A.png}}
\opt{B}{\bfigi{76.2}{TOAS09_p2_phy_s1_q20_B.png}}
\opt{C}{\bfigi{75.6}{TOAS09_p2_phy_s1_q20_C.png}}
\opt{D}{\bfigi{79.8}{TOAS09_p2_phy_s1_q20_D.png}}
\end{qoptions}
\end{bankq}
\begin{bankq}[id=TOAS09-P2-PHY-S2-Q05, source={TOAS - Test 09 (30 May 2022)}, paper=2, qno=5, type=numerical, subtopic={Refraction}, answer={2}, answer_src=solved, working={Vertical boundary: $\mu_1\sin i=\mu_2\sin r$ with $\sin i=1/\sqrt2$, $\sin r=\sqrt3/2$ gives $(\mu_1/\mu_2)^2=3/2$ and the ratio $=2$}, difficulty=2, diff_parts={K1 S1 I0 C0 F1}, diff_note={concepts: Snell's law with direction vectors | steps: angles from the unit vectors; ratio of indices}]
A ray incident on medium boundary as shown in figure. Unit vector in direction of incident ray is given as $\dfrac1{\sqrt2}\hat i+\dfrac1{\sqrt2}\hat j$ \& unit vector of ray in medium (2) is given by $\dfrac12\hat i+\dfrac{\sqrt3}2\hat j$. The ratio of
\[\dfrac{\left(\frac{\mu_1}{\mu_2}\right)^2+\frac12}{\left(\frac{\mu_1}{\mu_2}\right)^2-\frac12}\]
is
\bfig{113.4}{TOAS09_p2_phy_s2_q05_1.png}
\end{bankq}
\begin{bankq}[id=TOAS09-P2-PHY-S2-Q06, source={TOAS - Test 09 (30 May 2022)}, paper=2, qno=6, type=numerical, subtopic={Prisms}, answer={3}, answer_src=solved, working={Refraction at the prism face: $\sin r=1.5\sin30^\circ=\frac34$; deviation after the vertical mirror is $210^\circ-\sin^{-1}\frac34$ so $x=3$}, difficulty=4, diff_parts={K2 S2 I0 C0 F1}, diff_note={concepts: prism refraction, reflection | steps: refraction at the prism face; reflection from the mirror; total deviation}]
A ray incident normal to the surface of prism as shown in figure. Angle of deviation of ray from the mirror w.r.t incident ray is $210^\circ-\sin^{-1}\dfrac x4$, find the value of x.
\bfig{128.3}{TOAS09_p2_phy_s2_q06_1.png}
\end{bankq}
\begin{bankq}[id=TOAS11-P1-PHY-S1-Q03, source={TOAS - Test 11 (27 Jun 2022)}, paper=1, qno=3, type=single, subtopic={Refraction}, answer={B}, answer_src=key, difficulty=2, diff_parts={K1 S2 I0 C0 F0}, diff_note={concepts: refraction through a sphere | steps: image from the first surface; object for the second surface; condition for parallel emergent rays}]
A point object O is placed in front of a sphere of radius $R$ and $\mu=\dfrac32$ at distance $\dfrac RN$. Find value of $N$ so that emergent rays are parallel to $x$-axis. Consider paraxial rays only:
\bfig{143.7}{TOAS11_p1_phy_s1_q03_1.png}
\begin{qoptions}
\opt{A}{1}
\opt{B}{2}
\opt{C}{3}
\opt{D}{4}
\end{qoptions}
\end{bankq}
\begin{bankq}[id=TOAS11-P2-PHY-S3-Q06, source={TOAS - Test 11 (27 Jun 2022)}, paper=2, qno=6, type=numerical, subtopic={Lenses}, answer={90.00}, answer_src=key, difficulty=4, diff_parts={K2 S2 I0 C0 F1}, diff_note={concepts: lens in water, a mirror | steps: focal length of the lens in water; refraction at the lens; reflection from the mirror; image through the lens again}]
A thin equiconvex lens of glass of refractive index $\mu=\dfrac32$ and of focal length 0.3 m in air is sealed into an opening at one end of tank filled with water $\left(\mu=\dfrac43\right)$. On the opposite side of the lens, a mirror is placed inside the tank on the tank wall perpendicular to the lens axis as shown in figure. The separation between the lens and the mirror is 0.8 m. A small object is placed outside the tank in front of the lens at a distance of 0.9 m from the lens along its axis. Find the position (relative to the lens) of final image formed by the system (in cm).
\bfig{183.0}{TOAS11_p2_phy_s3_q06_1.png}
\end{bankq}
\begin{bankq}[id=TOAS12-P1-PHY-S1-Q02, source={TOAS - Test 12 (10 Jul 2022)}, paper=1, qno=2, type=multi, subtopic={Refraction}, answer={A,C,D}, answer_src=key, difficulty=3, diff_parts={K1 S2 I0 C0 F1}, diff_note={concepts: refraction through concentric spheres | steps: incident tangentially gives the first refraction angle; angle at the inner sphere; total deviation; options}]
A ray of light tangentialy incident on outer sphere.
\bfig{92.0}{TOAS12_p1_phy_s1_q02_1.png}
\begin{qoptions}
\opt{A}{Net deviation suffered by ray is $240^\circ$.}
\opt{B}{Ray will never enter inside inner sphere.}
\opt{C}{Ray is incident tangentially on inner sphere.}
\opt{D}{Angle made by ray with normal inside inner sphere is $30^\circ$.}
\end{qoptions}
\end{bankq}
\begin{bankq}[id=TOAS12-P1-PHY-S1-Q03, source={TOAS - Test 12 (10 Jul 2022)}, paper=1, qno=3, type=multi, subtopic={Lenses}, answer={A,B,D}, answer_src=key, difficulty=2, diff_parts={K1 S1 I0 C0 F1}, diff_note={concepts: magnification versus image distance | steps: m = v/u and the lens formula give m = v/f - 1 so slope is 1/f and the m intercept is -1}]
Variation of magnification of convex lens vs image distance for real objects is shown. Choose correct statement.
\bfig{67.8}{TOAS12_p1_phy_s1_q03_1.png}
\begin{qoptions}
\opt{A}{$f$ is equal to intercept of V-axis.}
\opt{B}{$f$ is equal to $\dfrac{1}{\text{slope}}$ of graph.}
\opt{C}{$f$ is equal to slope of graph}
\opt{D}{magnitude of intercept on $m$ axis is 1.}
\end{qoptions}
\end{bankq}
\begin{bankq}[id=TOAS12-P1-PHY-S2-Q02, source={TOAS - Test 12 (10 Jul 2022)}, paper=1, qno=2, type=numerical, subtopic={Refraction}, answer={1.33}, answer_src=key, difficulty=2, diff_parts={K1 S1 I0 C0 F1}, diff_note={concepts: Snell's law with direction vectors | steps: sin theta is the x-component; 1.5 a = 2 c}]
x-z plane seprates 2 medium A \& B having $\mu_1=1.5$ \& $\mu_2=2$. A ray of light passes through two medium from A to B. If unit vector of ray are $a\hat i+b\hat j$ \& $c\hat i+d\hat j$ in A \& B. then value of $a/c$ is
\end{bankq}
\begin{bankq}[id=TOAS12-P2-PHY-S1-Q04, source={TOAS - Test 12 (10 Jul 2022)}, paper=2, qno=4, type=single, subtopic={Lenses}, answer={D}, answer_src=key, difficulty=1, diff_parts={K1 S1 I0 C0 F0}, diff_note={concepts: lens in two different liquids | steps: focal lengths of each half with the surrounding media; type of beam from each half}]
An equi-convex lens is made of a material of refractive index $\mu_1$ and is placed inside two liquids of refractive indices $\mu_2$ and $\mu_3$ as shown with $\mu_2>\mu_1>\mu_3$. A wide, parallel beam of light is incident on the lens from the left. The lens will give rise to :
\bfig{135.2}{TOAS12_p2_phy_s1_q04_1.png}
\begin{qoptions}
\opt{A}{A single convergent beam.}
\opt{B}{Two different convergent beams.}
\opt{C}{Two different divergent beams}
\opt{D}{A convergent and divergent beam.}
\end{qoptions}
\end{bankq}
\begin{bankpassage}[id=TOAS12-P2-PHY-S2-G9-10, source={TOAS - Test 12 (10 Jul 2022)}, questions={9-10}]
A point object is placed at a distance 50 cm from the pole of a concave mirror. The radius of curvature of mirror is 30 cm. Point object O oscillates with amplitude of 1mm perpendicular to the principal axis as shown.
\bfig{97.0}{TOAS12_p2_phy_s2_para9_10_1.png}
\end{bankpassage}
\begin{bankq}[id=TOAS12-P2-PHY-S2-Q09, source={TOAS - Test 12 (10 Jul 2022)}, paper=2, qno=9, type=multi, subtopic={Reflection}, answer={A}, answer_src=key, difficulty=3, diff_parts={K1 S2 I0 C0 F1}, diff_note={concepts: image amplitude from magnification | steps: image distance by the mirror formula; transverse magnification; amplitude}, group=TOAS12-P2-PHY-S2-G9-10]
The amplitude of image is :-
\begin{qoptions}
\opt{A}{$\dfrac37$ mm}
\opt{B}{$\dfrac27$ mm}
\opt{C}{$\dfrac43$ mm}
\opt{D}{$\dfrac{11}{7}$ mm}
\end{qoptions}
\end{bankq}
\begin{bankq}[id=TOAS12-P2-PHY-S2-Q10, source={TOAS - Test 12 (10 Jul 2022)}, paper=2, qno=10, type=multi, subtopic={Reflection}, answer={A}, answer_src=key, difficulty=3, diff_parts={K1 S1 I1 C0 F1}, diff_note={concepts: phase of the image motion | steps: transverse magnification is negative so the image moves oppositely | I: real inverted image}, group=TOAS12-P2-PHY-S2-G9-10]
Phase difference between motion of object and image when object crosses the principal axis is :-
\begin{qoptions}
\opt{A}{$\pi$}
\opt{B}{0}
\opt{C}{$\dfrac\pi2$}
\opt{D}{None of these}
\end{qoptions}
\end{bankq}
\begin{bankq}[id=TOAS12-P2-PHY-S3-Q03, source={TOAS - Test 12 (10 Jul 2022)}, paper=2, qno=3, type=numerical, subtopic={Lenses}, answer={26.00}, answer_src=key, difficulty=3, diff_parts={K1 S2 I0 C0 F1}, diff_note={concepts: off-axis parallel beam through two lenses | steps: image by the convex lens at its focus; offset between axes; refraction at the concave lens; coordinates}]
If the optical axis of convex and concave lenses are separated by a distance 5 mm shown in the figure. If the coordinates of final image formed by the combination of Lenses is $(x,y)$ when the Parallel Beam of Light is incident on Convex Lens. If origin is at the Optical centre of convex lens, find $x+4y$ (in cm).
\bfig{129.8}{TOAS12_p2_phy_s3_q03_1.png}
\end{bankq}
\begin{bankq}[id=TOAS13-P1-PHY-S3-Q03, source={TOAS - Test 13 (24 Jul 2022)}, paper=1, qno=3, type=numerical, subtopic={Lenses}, answer={5.00}, answer_src=key, difficulty=7, diff_parts={K2 S3 I1 C1 F1}, diff_note={concepts: lens and mirror imaging, relative motion | steps: focal length of the lens from the radii; image positions through lens then mirror then lens at time t; condition for final image at the object; solve for the first time; solve for the second time; read N}]
An equiconvex lens of refractive index $\dfrac32$ and of radius of curvature 16 cm is placed co-axially at a distance of 90 cm in front of concave mirror of radius of curvature 60 cm. The surrounding medium is air. A point object O is placed on the common axis of the lens and the mirror at a distance 90 cm left of the lens. At an instant ($t=0$) the object, the lens and the mirror start moving with a speed 2 cm/s, 3 cm/s and 8 cm/s respectively in the direction shown in the figure. If the time after which the object and its final image coincide for the second time is $\dfrac{58}{N}$ second then $N$.
\bfig{113.1}{TOAS13_p1_phy_s3_q03_1.png}
\end{bankq}
\begin{bankq}[id=TOAS13-P2-PHY-S1-Q04, source={TOAS - Test 13 (24 Jul 2022)}, paper=2, qno=4, type=single, subtopic={Lenses}, answer={A}, answer_src=key, difficulty=1, diff_parts={K1 S1 I0 C0 F0}, diff_note={concepts: equivalent focal length of two lenses | steps: 1/F = 1/f1 + 1/f2 - d/(f1 f2) must be negative so d > 30 cm}]
Two convex lenses of focal length 10 cm and 20 cm respectively placed coaxially and are separated by some distance `d'. The whole system behaves like a concave lens for rays coming from infinity. One of the possible value of `d' is :
\begin{qoptions}
\opt{A}{40 cm}
\opt{B}{25 cm}
\opt{C}{15 cm}
\opt{D}{20 cm}
\end{qoptions}
\end{bankq}
\begin{bankq}[id=TOAS13-P2-PHY-S2-Q06, source={TOAS - Test 13 (24 Jul 2022)}, paper=2, qno=6, type=numerical, subtopic={Microscope and telescope}, answer={36.00}, answer_src=key, difficulty=2, diff_parts={K1 S1 I0 C0 F1}, diff_note={concepts: astronomical telescope | steps: M = fo/fe = 6 and fo + fe = 42}]
An astronomical telescope has angular magnification of magnitude 6 for distant objects. The separation between object and eye piece is 42 cm and final image is formed at infinity. Find focal length of objective (in cm).
\end{bankq}
\begin{bankq}[id=TOAS14-P1-PHY-S1-Q06, source={TOAS - Test 14 (07 Aug 2022)}, paper=1, qno=6, type=multi, subtopic={Prisms}, answer={A,C}, answer_src=key, difficulty=5, diff_parts={K2 S2 I1 C0 F1}, diff_note={concepts: refraction, total internal reflection in a prism | steps: refraction angle at the first face; geometry of the two reflections; relate mu to theta; heights ratio}]
When a ray parallel to base of right angled prism incident as shown in figure after two total internal reflection it emerge parallel to incident ray then :
\bfig{171.7}{TOAS14_p1_phy_s1_q06_1.png}
\begin{qoptions}
\opt{A}{Refractive index of prism is $\dfrac{\cos\theta}{\cos3\theta}$.}
\opt{B}{Refractive index of prism is $\dfrac{\sin3\theta}{\sin\theta}$.}
\opt{C}{$\dfrac{h_2}{h_1}=\dfrac{\sin\theta}{\sin3\theta}$}
\opt{D}{$\dfrac{h_2}{h_1}=\dfrac{\cos3\theta}{\cos\theta}$}
\end{qoptions}
\end{bankq}
\begin{bankq}[id=TOAS15-P1-PHY-S2-Q10, source={TOAS - Test 15 (22 Aug 2022)}, paper=1, qno=10, type=multi, subtopic={Lenses}, answer={A,D}, answer_src=key, difficulty=6, diff_parts={K2 S2 I1 C1 F1}, diff_note={concepts: lens, mirror combination with a water layer | steps: apparent depth through the water; image by the lens; condition for retracing the object; solve x}]
In given situation find possible values of distance $x$ between lens and mirror so that final image forms at object itself. Refractive index of water is $4/3$.
\bfig{179.6}{TOAS15_p1_phy_s2_q10_1.png}
\begin{qoptions}
\opt{A}{96 cm}
\opt{B}{36 cm}
\opt{C}{66 cm}
\opt{D}{156 cm}
\end{qoptions}
\end{bankq}
\begin{bankq}[id=TOAS15-P1-PHY-S3-Q05, source={TOAS - Test 15 (22 Aug 2022)}, paper=1, qno=5, type=numerical, subtopic={Refraction}, answer={2.00}, answer_src=key, difficulty=5, diff_parts={K2 S2 I1 C0 F1}, diff_note={concepts: refraction, total internal reflection at a hexagonal body | steps: angles at AB; angle at the next faces; critical angle condition at CD; minimum index}]
A transparent hexagonal body ABCDEF is shown find minimum value of refractive index of body so that ray incidenting at surface AB do not reach observer through surface CD.
\bfig{84.4}{TOAS15_p1_phy_s3_q05_1.png}
\end{bankq}
\begin{bankq}[id=TOAS15-P2-PHY-S3-Q02, source={TOAS - Test 15 (22 Aug 2022)}, paper=2, qno=2, type=numerical, subtopic={Reflection}, answer={5}, answer_src=key, difficulty=6, diff_parts={K2 S2 I1 C1 F1}, diff_note={concepts: lens image, tilted plane mirror | steps: image by the lens; virtual object for the mirror; reflection about a tilted line; solve for d}]
A convex lens of focal length 1.5m is placed in a system of coordinate axis such that its optical centre is at origin and principal axis coinciding with the x-axis. An object and a plane mirror are arranged on the principal axis as shown in figure. Find the value of $d$ (in m) so that y-coordinate of final image (after refraction and reflection) is 0.3m. (Take $\tan\theta=0.3$)
\bfig{156.5}{TOAS15_p2_phy_s3_q02_1.png}
\end{bankq}
\begin{bankq}[id=TOAS16-P2-PHY-S1-Q06, source={TOAS - Test 16 (04 Sep 2022)}, paper=2, qno=6, type=multi, subtopic={Refraction}, answer={B,C,D}, answer_src=key, difficulty=5, diff_parts={K1 S2 I1 C1 F1}, diff_note={concepts: refraction at a spherical surface | steps: object distance from the apparent depth; the apparent position from the other side; no total internal reflection check}]
A transparent sphere of refractive index 1.5 has radius $R$. There is a point light source $S$ inside the sphere at some distance from center. A person observes the light source along diametric line passing through $S$. Source appears at $\dfrac47R$ distance from surface.
\begin{qoptions}
\opt{A}{Source is at $\dfrac{2R}3$ distance from center of sphere.}
\opt{B}{When person observes the source from diametrically opposite point it appears $\dfrac{8R}5$ distance from surface.}
\opt{C}{All light from source passes the spherical boundary without any total internal reflection.}
\opt{D}{Source is at $\dfrac R3$ distance from center of sphere}
\end{qoptions}
\end{bankq}
\begin{bankpassage}[id=TOAS16-P2-PHY-S3-G3-4, source={TOAS - Test 16 (04 Sep 2022)}, questions={3-4}]
Large number of identical thin prisms are placed on $x$-axis. Each has refractive index 1.5 and angle of prism $\dfrac\pi2\times10^{-3}$rad. A light ray incidents parallel to $x$-axis along line $y=y_0$. Take thickness of each prism $t=1$ mm. Take $\ln2=0.7$.
\bfig{134.9}{TOAS16_p2_phy_s3_para3_4_1.png}
\end{bankpassage}
\begin{bankq}[id=TOAS16-P2-PHY-S3-Q03, source={TOAS - Test 16 (04 Sep 2022)}, paper=2, qno=3, type=numerical, subtopic={Prisms}, answer={0.785}, answer_src=solved, working={Each thin prism deviates by $(\mu-1)A=\tfrac\pi4\times10^{-3}$ rad; 1000 prisms in 1 m give $\theta=\pi/4$}, difficulty=4, diff_parts={K1 S2 I1 C0 F1}, diff_note={concepts: thin prism deviation and accumulation | steps: deviation of one prism; number of prisms in 1 m; sum}, group=TOAS16-P2-PHY-S3-G3-4]
When ray reaches $x=1$m angle made by ray with $x$-axis is $\theta$ rad. Value of $\theta$ is \ldots\ldots
\end{bankq}
\begin{bankq}[id=TOAS16-P2-PHY-S3-Q04, source={TOAS - Test 16 (04 Sep 2022)}, paper=2, qno=4, type=numerical, subtopic={Prisms}, answer={0.44 to 0.45}, answer_src=key, difficulty=7, diff_parts={K2 S2 I2 C1 F1}, diff_note={concepts: continuous deviation of a ray, integration of the path | steps: angle as a function of x; dy/dx = tan theta; integrate with the logarithm}, group=TOAS16-P2-PHY-S3-G3-4]
When ray reaches $x=1$m its $y$ displacement magnitude in meter is \ldots\ldots\ldots
\end{bankq}
\begin{bankq}[id=TOAS17-P1-PHY-S1-Q02, source={TOAS - Test 17 (18 Sep 2022)}, paper=1, qno=2, type=multi, subtopic={Reflection}, answer={B,C,D}, answer_src=key, difficulty=4, diff_parts={K1 S2 I1 C0 F1}, diff_note={concepts: multiple reflections in a square of mirrors | steps: unfold the reflections; condition to retrace to the hole for each angle; check each option}]
Four plane mirrors form a square with each side of length $L$. There is a small hole at one of the corners as shown in the figure. A ray of light passed through the hole at an angle $\theta$ and can come out of same hole. The cross-section of mirror configuration and the ray of light lies on the same plane. Which of the following statements are correct.
\bfig{134.5}{TOAS17_p1_phy_s1_q02_1.png}
\begin{qoptions}
\opt{A}{A ray of light will come out for $\theta=45^\circ$ only if $\ell=\dfrac L2$.}
\opt{B}{A ray of light will come out for $\theta=90^\circ$ for $0<\ell<L$.}
\opt{C}{A ray of light will come out for $\theta=\tan^{-1}\left(\dfrac1{2n}\right)$, where $n$ is a natural number for $\ell=\dfrac L2$.}
\opt{D}{A ray of light will come out for $\theta=45^\circ$ for $0<\ell<L$.}
\end{qoptions}
\end{bankq}
\begin{bankpassage}[id=TOAS17-P1-PHY-S2-G11-13, source={TOAS - Test 17 (18 Sep 2022)}, questions={11-13}]
Column-1 represent distance of real object from optical instrument-1 given in Column-2. 2$^{\text{nd}}$ optical instrument place in contact with 1$^{\text{st}}$ is given column-3.\par
\begin{center}\footnotesize\setlength{\tabcolsep}{1pt}
\begin{tabular}{|L{0.7cm}L{0.8cm}|L{0.6cm}L{1.9cm}|L{0.5cm}L{1.9cm}|}
\hline
\multicolumn{2}{|c|}{Column-1}&\multicolumn{2}{c|}{Column-2}&\multicolumn{2}{c|}{Column-3}\\\hline
(I)&15 cm&(i)&\bfigi{55.6}{TOAS17_p1_phy_s2_para11_13_1.png}&(P)&\bfigi{35.1}{TOAS17_p1_phy_s2_para11_13_5.png}\\\hline
(II)&10 cm&(ii)&\bfigi{55.6}{TOAS17_p1_phy_s2_para11_13_2.png}&(Q)&\bfigi{60.5}{TOAS17_p1_phy_s2_para11_13_6.png}\\\hline
(III)&20 cm&(iii)&\bfigi{55.6}{TOAS17_p1_phy_s2_para11_13_3.png}&(R)&\bfigi{61.2}{TOAS17_p1_phy_s2_para11_13_7.png}\\\hline
(IV)&30 cm&(iv)&\bfigi{51.4}{TOAS17_p1_phy_s2_para11_13_4.png}&(S)&\bfigi{52.9}{TOAS17_p1_phy_s2_para11_13_8.png}\\\hline
\end{tabular}\end{center}
\end{bankpassage}
\begin{bankq}[id=TOAS17-P1-PHY-S2-Q11, source={TOAS - Test 17 (18 Sep 2022)}, paper=1, qno=11, type=single, subtopic={Lenses}, answer={C}, answer_src=key, difficulty=5, diff_parts={K2 S2 I1 C1 F0}, diff_note={concepts: combination of two optical elements in contact, retracing | steps: power of the combination; image position; condition to coincide with the object}, group=TOAS17-P1-PHY-S2-G11-13]
In which of following case final image will form on object itself.
\begin{qoptions}
\opt{A}{(II), (i), (Q)}
\opt{B}{(I), (i), (P)}
\opt{C}{(I), (i), (R)}
\opt{D}{(III), (iii), (R)}
\end{qoptions}
\end{bankq}
\begin{bankq}[id=TOAS17-P1-PHY-S2-Q12, source={TOAS - Test 17 (18 Sep 2022)}, paper=1, qno=12, type=single, subtopic={Lenses}, answer={A}, answer_src=key, difficulty=5, diff_parts={K2 S2 I1 C1 F0}, diff_note={concepts: combination of two optical elements in contact, image distance | steps: focal length of the pair; image distance; separation of object and image}, group=TOAS17-P1-PHY-S2-G11-13]
In which case distance between object and it's final virtual image will be 24 cm.
\begin{qoptions}
\opt{A}{(II), (iv), (P)}
\opt{B}{(III), (ii), (Q)}
\opt{C}{(IV), (iii), (S)}
\opt{D}{(I), (i), (R)}
\end{qoptions}
\end{bankq}
\begin{bankq}[id=TOAS17-P1-PHY-S2-Q13, source={TOAS - Test 17 (18 Sep 2022)}, paper=1, qno=13, type=single, subtopic={Lenses}, answer={D}, answer_src=key, difficulty=5, diff_parts={K2 S2 I1 C1 F0}, diff_note={concepts: combination of two optical elements, magnification | steps: focal length of the pair; condition for real image of equal size; object distance}, group=TOAS17-P1-PHY-S2-G11-13]
In which case real image of same size as object forms.
\begin{qoptions}
\opt{A}{(I), (ii), (R)}
\opt{B}{(II), (iii), (P)}
\opt{C}{(III), (iv), (R)}
\opt{D}{(IV), (ii), (S)}
\end{qoptions}
\end{bankq}
\begin{bankq}[id=TOAS17-P2-PHY-S2-Q01, source={TOAS - Test 17 (18 Sep 2022)}, paper=2, qno=1, type=numerical, subtopic={Refraction}, answer={30}, answer_src=key, difficulty=4, diff_parts={K1 S2 I1 C0 F1}, diff_note={concepts: refraction at the boundary of two slabs | steps: angle in the first medium; Snell law at the diagonal; deviation}]
Refractive indices of 2 different media with separating boundary at the diagonal of rectangular glass slab are shown. Then angle of deviation (in degree) of the ray in going from $n_1$ to $n_2$ as shown in the figure is :
\bfig{139.9}{TOAS17_p2_phy_s2_q01_1.png}
\end{bankq}
\begin{bankq}[id=TOAS17-P2-PHY-S2-Q04, source={TOAS - Test 17 (18 Sep 2022)}, paper=2, qno=4, type=numerical, subtopic={Reflection}, answer={16}, answer_src=key, difficulty=3, diff_parts={K1 S1 I1 C0 F1}, diff_note={concepts: image by a silvered lens | steps: power of the silvered lens; condition to form the image on the object}]
Radii of curvature of a concavo-convex lens (refractive index = 1.5) are 40 cm (concave side) and 20 cm (convex side) as shown. The convex side is silvered. The distance $x$ on the principal axis where object is placed so that its image is created on the object itself, is equal to (in cm) :
\bfig{112.4}{TOAS17_p2_phy_s2_q04_1.png}
\end{bankq}
\begin{bankq}[id=TOAS18-P1-PHY-S2-Q08, source={TOAS - Test 18 (30 Oct 2022)}, paper=1, qno=8, type=single, subtopic={Refraction}, answer={D}, answer_src=solved, working={Lens deviates by $h/f=\pi/100$ rad $=1.8^\circ$ (R); the ray totally reflects at the $\mu_2/\mu_3$ face ($2\sin60^\circ/\sqrt2>1$) and returns, deviation $60^\circ$ (Q); prism with $\mu=\sqrt3$ at $60^\circ$ is symmetric, deviation $60^\circ$ (Q); concave mirror: $v=-30$ cm, deviation $180^\circ+\beta-\alpha=179.4^\circ$ (T)}, difficulty=7, diff_parts={K3 S3 I1 C1 F0}, diff_note={concepts: thin lens deviation, total internal reflection, prism, mirror | steps: lens deviation h/f; critical angle in the slab and the reflected direction; symmetric prism deviation; mirror image distance; deviation of the reflected ray}]
List-I contain some optical arrangement and list-II contains deviation produced after all possible optical event in degree then select correct answer using following codes.
\begin{center}\footnotesize\setlength{\tabcolsep}{2pt}
\begin{tabular}{|L{0.7cm}L{4.4cm}|L{0.6cm}L{1.3cm}|}\hline
\multicolumn{2}{|c|}{List-I}&\multicolumn{2}{c|}{List-II}\\\hline
(I)&\bfigi{92.5}{TOAS18_p1_phy_s2_q08_1.png}&(P)&0\\\hline
(II)&\bfigi{110.6}{TOAS18_p1_phy_s2_q08_2.png}&(Q)&$60^\circ$\\\hline
(III)&\bfigi{79.4}{TOAS18_p1_phy_s2_q08_3.png}&(R)&$1.8^\circ$\\\hline
(IV)&\bfigi{87.2}{TOAS18_p1_phy_s2_q08_4.png}&(S)&$178.2^\circ$\\\hline
&&(T)&$179.4^\circ$\\\hline
\end{tabular}\end{center}
\begin{qoptions}
\opt{A}{I $\to$ R;II $\to$ P;III $\to$ Q;IV $\to$ R}
\opt{B}{I $\to$ R;II $\to$ Q;III $\to$ Q;IV $\to$ S}
\opt{C}{I $\to$ Q;II $\to$ P;III $\to$ S;IV $\to$ R}
\opt{D}{I $\to$ R;II $\to$ Q;III $\to$ Q;IV $\to$ T}
\end{qoptions}
\end{bankq}
\begin{bankq}[id=TOAS18-P2-PHY-S2-Q07, source={TOAS - Test 18 (30 Oct 2022)}, paper=2, qno=7, type=multi, subtopic={Lenses}, answer={A,B,C}, answer_src=key, difficulty=2, diff_parts={K1 S1 I0 C0 F1}, diff_note={concepts: lens magnification | steps: sign of magnification and image type; power of the lens}]
A real object is kept perpendicular to principal axis of a lens in air. With respect to different cases of image formation by lens, choose \textbf{CORRECT} option(s) :
\begin{qoptions}
\opt{A}{If magnification is $-2.5$ then image is real \& power of lens is positive.}
\opt{B}{If magnification is 0.5 then image is virtual \& power of lens is negative.}
\opt{C}{If length of image is equal to object then image is real \& power of lens is positive.}
\opt{D}{If length of image is equal to object then image is virtual \& power of lens is negative.}
\end{qoptions}
\end{bankq}
\begin{bankq}[id=TOAS18-P2-PHY-S3-Q02, source={TOAS - Test 18 (30 Oct 2022)}, paper=2, qno=2, type=numerical, subtopic={Refraction}, answer={3}, answer_src=key, difficulty=5, diff_parts={K1 S2 I1 C1 F1}, diff_note={concepts: ray in a medium of varying refractive index | steps: Snell invariant; slope of the path; refractive index at the point}]
A ray of light travelling in air is incident at $30^\circ$ on one surface of a slab of variable refractive index. If light travels along path $y=4x^2$ then value of $\mu$ at $y=\dfrac12$ inside the slab is $\dfrac x2$. Find $x$.
\bfig{119.5}{TOAS18_p2_phy_s3_q02_1.png}
\end{bankq}

\section*{TAAS papers}
\begin{bankq}[id=TAAS01-P1-PHY-S2-Q06, source={TAAS - Test 01 (18 Jun 2023)}, paper=1, qno=6, type=multi, subtopic={Reflection}, answer={A}, answer_src=key, difficulty=2, diff_parts={K1 S1 I0 C0 F1}, diff_note={concepts: images of virtual objects in mirrors | steps: mirror equation for a virtual object; sign of the image distance for each mirror}]
Which of the following options are correct for a small object placed very near to principal axis of a spherical mirror.
\begin{qoptions}
\opt{A}{A virtual object always have a real image in concave mirror.}
\opt{B}{A virtual object always have a virtual image in concave mirror.}
\opt{C}{A virtual object always have a real image in convex mirror.}
\opt{D}{A virtual object always have a virtual image in convex mirror.}
\end{qoptions}
\end{bankq}
\begin{bankq}[id=TAAS01-P1-PHY-S2-Q12, source={TAAS - Test 01 (18 Jun 2023)}, paper=1, qno=12, type=multi, subtopic={Reflection}, answer={A,B}, answer_src=key, difficulty=3, diff_parts={K1 S1 I1 C0 F1}, diff_note={concepts: images by spherical mirrors | steps: position of image relative to the axis and the mirror; check each case}]
$I$ is the image of a point object $O$ formed by spherical mirror then which of the following statements is/are correct ?
\begin{qoptions}
\opt{A}{If $O$ and $I$ are on same side of the principal axis, then they have to be on opposite side of the mirror.}
\opt{B}{If $O$ and $I$ are on opposite sides of the principal axis, then they have to be on same side of the mirror.}
\opt{C}{If $O$ and $I$ are on opposite side of the principal axis, then they can be on opposite side of the mirror as well.}
\opt{D}{If $O$ and $I$ are on same side of the principal axis, then they have to be on same side of the mirror.}
\end{qoptions}
\end{bankq}

\section*{TNAS papers}
\begin{bankq}[id=TNAS01-P1-PHY-S1-Q01, source={TNAS - Test 01 (05 Feb 2023)}, paper=1, qno=1, type=single, subtopic={Reflection}, answer={C}, answer_src=key, difficulty=4, diff_parts={K1 S2 I1 C1 F0}, diff_note={concepts: rotation of a mirror | steps: deviation by two mirrors; condition for horizontal; angle}]
By what angle should $M_2$ be rotated, so that the light ray after reflection from both the mirrors become horizontal :-
\bfig{108.3}{TNAS01_p1_phy_s1_q01_1.png}
\begin{qoptions}
\opt{A}{$5^\circ$ clockwise}
\opt{B}{$95^\circ$ clockwise}
\opt{C}{$5^\circ$ anticlockwise}
\opt{D}{$15^\circ$ anticlockwise}
\end{qoptions}
\end{bankq}
\begin{bankq}[id=TNAS01-P1-PHY-S1-Q02, source={TNAS - Test 01 (05 Feb 2023)}, paper=1, qno=2, type=single, subtopic={Refraction}, answer={C}, answer_src=key, difficulty=4, diff_parts={K1 S2 I1 C1 F0}, diff_note={concepts: refraction in vector form | steps: angle of incidence; Snell law; unit vector}]
A ray of light moving along vector $\hat i-2\hat j$ in denser medium undergoes refraction at an interface of two media, which is $x$-$z$ plane. The refractive index for $y>0$ is $2$ while for $y<0$ it is $\dfrac{\sqrt5}{2}$. The unit vector along the refracted ray is :-
\begin{qoptions}
\opt{A}{$\dfrac{-3\hat i-5\hat j}{\sqrt{34}}$}
\opt{B}{$\dfrac{-4\hat i+5\hat j}{5}$}
\opt{C}{$\dfrac{-3\hat i-4\hat j}{5}$}
\opt{D}{$\dfrac{4\hat i-3\hat j}{5}$}
\end{qoptions}
\end{bankq}
\begin{bankq}[id=TNAS01-P1-PHY-S1-Q03, source={TNAS - Test 01 (05 Feb 2023)}, paper=1, qno=3, type=single, subtopic={Refraction}, answer={D}, answer_src=key, difficulty=4, diff_parts={K1 S2 I1 C1 F0}, diff_note={concepts: apparent depth and mirror | steps: apparent position; mirror image; final distance}]
A bird is flying over a swimming pool at a height of $2$ m from the water surface. If the bottom is perfectly plane reflecting surface and depth of swimming pool is $1$m, then the distance of final image of bird from the bird itself is $\mu_w=4/3$.
\begin{qoptions}
\opt{A}{$\dfrac{11}{3}$m}
\opt{B}{$\dfrac{23}{3}$m}
\opt{C}{$\dfrac{11}{4}$m}
\opt{D}{$\dfrac{11}{2}$m}
\end{qoptions}
\end{bankq}
\begin{bankq}[id=TNAS01-P1-PHY-S1-Q04, source={TNAS - Test 01 (05 Feb 2023)}, paper=1, qno=4, type=single, subtopic={Reflection}, answer={B}, answer_src=key, difficulty=2, diff_parts={K1 S1 I0 C1 F0}, diff_note={concepts: image velocity of a moving mirror | steps: normal component relation; tangential unchanged}]
A plane mirror is moving with velocity $4\hat i+4\hat j+8\hat k$. A point object in front of the mirror moves with a velocity $3\hat i+4\hat j+5\hat k$. Here, $\hat k$ is along the normal to the plane mirror and facing towards the object. The velocity of the image is :-
\begin{qoptions}
\opt{A}{$-3\hat i-4\hat j+5\hat k$}
\opt{B}{$3\hat i+4\hat j+11\hat k$}
\opt{C}{$-4\hat i+5\hat j+11\hat k$}
\opt{D}{$7\hat i+9\hat j+3\hat k$}
\end{qoptions}
\end{bankq}
\begin{bankq}[id=TNAS01-P1-PHY-S1-Q20, source={TNAS - Test 01 (05 Feb 2023)}, paper=1, qno=20, type=single, subtopic={Reflection}, answer={B}, answer_src=key, difficulty=5, diff_parts={K1 S2 I2 C1 F0}, diff_note={concepts: image velocity in a concave mirror | steps: image position; longitudinal and transverse magnification; velocity components}]
In a three dimensional coordinate system (OXYZ), a concave mirror of radius of curvature $40$ cm is placed at $x=80$ cm. An object placed at origin is given a velocity $\vec v_0=\left(9\hat i+6\hat j+3\hat k\right)$ cm/s. Find out the magnitude of velocity of its image.
\bfig{79.4}{TNAS01_p1_phy_s1_q20_1.png}
\begin{qoptions}
\opt{A}{$\sqrt{46}$ cm/s}
\opt{B}{$\sqrt6$ cm/s}
\opt{C}{$\dfrac{\sqrt{14}}{3}$ cm/s}
\opt{D}{None of these}
\end{qoptions}
\end{bankq}
\begin{bankq}[id=TNAS01-P1-PHY-S2-Q06, source={TNAS - Test 01 (05 Feb 2023)}, paper=1, qno=6, type=numerical, subtopic={Reflection}, answer={1}, answer_src=key, difficulty=6, diff_parts={K1 S2 I2 C1 F1}, diff_note={concepts: mirror image speed | steps: focal length; speed ratio; value}]
When an object is kept at a distance of $30$ cm from a concave mirror, the image is formed at a distance of $10$ cm. If the object is moved toward the mirror with a speed of $9$ m/s, the speed with which image moves is \underline{\hspace{0.8cm}} m/s.
\end{bankq}
\begin{bankq}[id=TNAS01-P1-PHY-S2-Q09, source={TNAS - Test 01 (05 Feb 2023)}, paper=1, qno=9, type=numerical, subtopic={Refraction}, answer={7}, answer_src=key, difficulty=4, diff_parts={K1 S1 I1 C1 F1}, diff_note={concepts: apparent shift in layered media | steps: shift of each layer; sum}]
A tank contains three layers of immiscible liquids. The first layer is of water with refractive index $4/3$ and thickness $4$ cm. The second layer is of oil with refractive index $3/2$ and thickness $6$ cm while third layer is of liquid with refractive index $2$ and thickness $8$ cm. The shift in the position of bottom when viewed from the above as shown is \underline{\hspace{0.8cm}} cm.
\bfig{108.2}{TNAS01_p1_phy_s2_q09_1.png}
\end{bankq}
\begin{bankq}[id=TNAS01-P1-PHY-S2-Q10, source={TNAS - Test 01 (05 Feb 2023)}, paper=1, qno=10, type=numerical, subtopic={Refraction}, answer={6}, answer_src=key, difficulty=6, diff_parts={K1 S2 I2 C1 F1}, diff_note={concepts: image by a mirror at an incline | steps: mirror image geometry; coordinates; distance}]
Distance between image and object placed at origin is `$n$' cm, then $n$ is
\bfig{108.3}{TNAS01_p1_phy_s2_q10_1.png}
\end{bankq}
\begin{bankq}[id=TNAS01-P2-PHY-S1-Q04, source={TNAS - Test 01 (05 Feb 2023)}, paper=2, qno=4, type=multi, subtopic={Reflection}, answer={A,C}, answer_src=key, difficulty=5, diff_parts={K1 S2 I1 C1 F1}, diff_note={concepts: mirror maze | steps: virtual paths; trace; visible monsters}]
Figure is an overhead view of a mirror maze based on floor sections that are equilateral triangles. Every wall within the maze is mirrored. If you stand at entrance $x$, which of the maze monsters $a$, $b$ and $c$ hiding in the maze can you see along the virtual hall ways extending from entrance $x$ ;
\bfig{132.2}{TNAS01_p2_phy_s1_q04_1.png}
\begin{qoptions}
\opt{A}{a}
\opt{B}{b}
\opt{C}{c}
\opt{D}{None}
\end{qoptions}
\end{bankq}
\begin{bankq}[id=TNAS01-P2-PHY-S1-Q06, source={TNAS - Test 01 (05 Feb 2023)}, paper=2, qno=6, type=multi, subtopic={Refraction}, answer={A,C}, answer_src=key, difficulty=6, diff_parts={K1 S2 I2 C1 F1}, diff_note={concepts: apparent distance through layers | steps: shift in each medium; sum; from both sides}]
A fish watcher at point $P$ watches a fish through a glass wall of a fish tank. The watcher is level with the fish ; the index of refraction of the glass is $8/5$, and that of the water is $4/3$. The distances are $d_1=8.0$ cm, $d_2=3.0$ cm, and $d_3=6.8$ cm.
\bfig{124.1}{TNAS01_p2_phy_s1_q06_1.png}
\begin{qoptions}
\opt{A}{The apparent separation between watcher and fish as observed by the watcher is $20$ cm.}
\opt{B}{The apparent separation between watcher and fish as observed by the watcher is $13.2$ cm.}
\opt{C}{The apparent separation between watcher and fish as observed by the fish is $15$ cm.}
\opt{D}{The apparent separation between watcher and fish as observed by the fish is $7$ cm.}
\end{qoptions}
\end{bankq}
\begin{bankq}[id=TNAS02-P1-PHY-S1-Q02, source={TNAS - Test 02 (07 May 2023)}, paper=1, qno=2, type=single, subtopic={Refraction}, answer={C}, answer_src=key, difficulty=4, diff_parts={K1 S2 I1 C1 F0}, diff_note={concepts: width of a refracted beam | steps: angle of refraction; width ratio; expression}]
A beam of light of width $t$ is incident on an air-water boundary at an angle of incidence $45^\circ$. The width of the beam in water is (Refractive index of water $=\mu$)
\begin{qoptions}
\opt{A}{$(\mu-1)t$}
\opt{B}{$\dfrac{\sqrt{2\mu^2-1}}{\sqrt2\,\mu}\cdot t$}
\opt{C}{$\dfrac{\sqrt{2\mu^2-1}}{\mu}\cdot t$}
\opt{D}{$\dfrac{\sqrt{\mu^2-1}}{\mu}\cdot t$}
\end{qoptions}
\end{bankq}
\begin{bankpassage}[id=TNAS02-P1-PHY-S3-G1-2, source={TNAS - Test 02 (07 May 2023)}, questions={1-2}]
\textbf{Rear view mirrors}\\
Rear view mirrors in automobiles are generally convex. Suppose a car A moves with a constant speed of 40.0 kilometer per hour on a straight level road and is followed by another car B moving with the constant speed 60.0 kilometer per hour. At a given instant of time, we denote :\\
$x$ : distance of the car B from the mirror of car A,\\
$y$ : distance of the car B from A as seen by the driver of A in the mirror,\\
$v_x$ : speed of approach of B relative to A and\\
$v_y$ : speed of approach of B as seen in the mirror of A.
\end{bankpassage}
\begin{bankq}[id=TNAS02-P1-PHY-S3-Q01, source={TNAS - Test 02 (07 May 2023)}, paper=1, qno=1, type=numerical, subtopic={Reflection}, answer={4}, answer_src=solved, working={Convex mirror $f=R/2$; $v_y=\left(\dfrac f{x+f}\right)^2v_x=\dfrac{R^2}{4(x+0.5R)^2}v_x$, so $n=4$}, difficulty=4, diff_parts={K1 S1 I1 C1 F1}, diff_note={concepts: image speed in a convex mirror | steps: differentiate mirror formula; factor}, group=TNAS02-P1-PHY-S3-G1-2]
If expression for $v_y$ in terms of $x$, $v_x$ and the radius of curvature $R$ of the convex mirror used as the rear view mirrors in the car A is $v_y=\dfrac{R^2}{n(x+0.5R)^2}v_x$. Find $n$.
\end{bankq}
\begin{bankq}[id=TNAS02-P1-PHY-S3-Q02, source={TNAS - Test 02 (07 May 2023)}, paper=1, qno=2, type=numerical, subtopic={Reflection}, answer={2.22}, answer_src=solved, working={Relative approach speed $v_x=20$ km/h; with $R=2$ m, $x=2$ m the factor is $\dfrac{4}{4(3)^2}=\dfrac19$, so $v_y=2.22$ km/h}, difficulty=3, diff_parts={K1 S1 I0 C1 F1}, diff_note={concepts: image speed | steps: apply the factor to the relative speed}, group=TNAS02-P1-PHY-S3-G1-2]
If $R=2.0$ m, what is the speed of approach of B in kilometer per hour as seen by the driver of A in the mirror for $x=2.0$ m ?
\end{bankq}
\begin{bankq}[id=TNAS07-P1-PHY-S1-Q06, source={TNAS - Test 07 (20 Aug 2023)}, paper=1, qno=6, type=single, subtopic={Microscope and telescope}, answer={C}, answer_src=key, difficulty=4, diff_parts={K1 S2 I1 C1 F0}, diff_note={concepts: telescope adjusted for a near point | steps: normal adjustment lengths; focal lengths from the magnifying power; lens formula for the final image; change in tube length}]
When a telescope is adjusted for parallel light, the distance of the objective from the eyepiece is 100 cm for normal adjustment. The magnifying power of the telescope in this case is 9. If an old man cannot see beyond 90 cm and wishes to use the telescope then he will have to reduce the tube length by
\begin{qoptions}
\opt{A}{1.5 cm}
\opt{B}{0.5 cm}
\opt{C}{1 cm}
\opt{D}{2 cm}
\end{qoptions}
\end{bankq}
\begin{bankq}[id=TNAS07-P1-PHY-S1-Q07, source={TNAS - Test 07 (20 Aug 2023)}, paper=1, qno=7, type=single, subtopic={Prisms}, answer={D}, answer_src=key, difficulty=4, diff_parts={K1 S2 I1 C1 F0}, diff_note={concepts: refraction through a prism | steps: refraction at the first face; angle at the second face; critical angle check; angle of emergence}]
A light ray is incident on a prism of prism angle $60^\circ$ and refractive index $\sqrt3$ as shown in figure. The angle of emergence $e$ is given by :
\bfig{109.4}{TNAS07_p1_phy_s1_q07_1.png}
\begin{qoptions}
\opt{A}{$30^\circ$}
\opt{B}{$45^\circ$}
\opt{C}{$90^\circ$}
\opt{D}{$60^\circ$}
\end{qoptions}
\end{bankq}
\begin{bankq}[id=TNAS07-P1-PHY-S1-Q08, source={TNAS - Test 07 (20 Aug 2023)}, paper=1, qno=8, type=single, subtopic={Lenses}, answer={C}, answer_src=key, difficulty=1, diff_parts={K1 S1 I0 C0 F0}, diff_note={concepts: correction of chromatic aberration | steps: achromatic combination}]
Chromatic aberration of a lens can be corrected by :
\begin{qoptions}
\opt{A}{providing different suitable curvature of its two surfaces}
\opt{B}{proper polishing of its two surface}
\opt{C}{suitably combining it with another lens}
\opt{D}{None of these}
\end{qoptions}
\end{bankq}
\begin{bankq}[id=TNAS07-P1-PHY-S1-Q09, source={TNAS - Test 07 (20 Aug 2023)}, paper=1, qno=9, type=single, subtopic={Reflection}, answer={D}, answer_src=key, difficulty=3, diff_parts={K1 S2 I1 C0 F0}, diff_note={concepts: image formation by a concave mirror | steps: position of the object; nature of the image; match}]
For a concave mirror. The correct matching of items on the left with those on the right is\par
{\footnotesize\setlength{\tabcolsep}{2pt}\begin{tabular}{|c|p{0.3\linewidth}|c|p{0.2\linewidth}|}
\hline
\multicolumn{2}{|c|}{\textbf{List-I}} & \multicolumn{2}{c|}{\textbf{List-II}}\\ \hline
(P) & Between centre of curvature and focus & (1) & enlarged, virtual, erect\\ \hline
(Q) & Beyond the centre of curvature (away from the mirror) & (2) & enlarged, real, inverted\\ \hline
(R) & Between focus and pole & (3) & same size, real, inverted\\ \hline
(S) & At the centre of curvature & (4) & diminished, real, inverted\\ \hline
\end{tabular}}\par
\begin{qoptions}
\opt{A}{(P)-(2); (Q)-(3) ; (R)-(1) ; (S)-(4)}
\opt{B}{(P)-(3); (Q)-(2) ; (R)-(1) ; (S)-(4)}
\opt{C}{(P)-(2) ; (Q)-(1) ; (R)-(3) ; (S)-(4)}
\opt{D}{(P)-(2) ; (Q)-(4) ; (R)-(1) ; (S)-(3)}
\end{qoptions}
\end{bankq}
\begin{bankq}[id=TNAS07-P1-PHY-S1-Q10, source={TNAS - Test 07 (20 Aug 2023)}, paper=1, qno=10, type=single, subtopic={Lenses}, answer={B}, answer_src=key, difficulty=4, diff_parts={K1 S2 I1 C1 F0}, diff_note={concepts: correction of long sightedness | steps: object at the book distance; image at the near point; lens power}]
A far sighted person has a near point 60 cm away. What should be the power of a lens he should use for eye glasses so that he can read a book at a distance of 25 cm?
\begin{qoptions}
\opt{A}{1.33D}
\opt{B}{2.33D}
\opt{C}{$-1.33$D}
\opt{D}{$-2.33$D}
\end{qoptions}
\end{bankq}
\begin{bankq}[id=TNAS07-P1-PHY-S1-Q11, source={TNAS - Test 07 (20 Aug 2023)}, paper=1, qno=11, type=single, subtopic={Lenses}, answer={C}, answer_src=key, difficulty=7, diff_parts={K1 S3 I2 C2 F0}, diff_note={concepts: image formed by a tilted lens | steps: object distance along the axis; lens formula; image position; project on the horizontal axis; coordinate}]
As shown in figure, a thin convex lens of focal length 20 cm is inclined at an angle of $53^\circ$ with $x$-axis. A point object is placed on the $x$-axis at a distance of 100 cm from the centre of the lens. Taking the origin as the centre of lens, find the $x$-coordinate of the image.
\bfig{133.0}{TNAS07_p1_phy_s1_q11_1.png}
\begin{qoptions}
\opt{A}{26.7 cm}
\opt{B}{25 cm}
\opt{C}{33.3 cm}
\opt{D}{20 cm}
\end{qoptions}
\end{bankq}
\begin{bankq}[id=TNAS07-P1-PHY-S1-Q12, source={TNAS - Test 07 (20 Aug 2023)}, paper=1, qno=12, type=single, subtopic={Reflection}, answer={A}, answer_src=key, difficulty=7, diff_parts={K2 S2 I2 C2 F0}, diff_note={concepts: rays reflected from a spherical mirror, spherical aberration | steps: angle of each ray; intersection point with the axis; difference}]
Two rays are incident on a spherical concave mirror of radius $R=5$ cm and rays are parallel to a line $AB$ which passes through the centre $C$. The rays are at perpendicular distances 3 cm and 4 cm from $AB$ respectively as shown. Find the distance between the points at which these rays intersect the line $AB$ after reflection :
\bfig{106.3}{TNAS07_p1_phy_s1_q12_1.png}
\begin{qoptions}
\opt{A}{$\dfrac{25}{24}$ cm}
\opt{B}{$\dfrac{24}{25}$ cm}
\opt{C}{zero}
\opt{D}{$\dfrac13$ cm}
\end{qoptions}
\end{bankq}
\begin{bankq}[id=TNAS07-P2-PHY-S2-Q07, source={TNAS - Test 07 (20 Aug 2023)}, paper=2, qno=7, type=single, subtopic={Reflection}, answer={C}, answer_src=key, difficulty=4, diff_parts={K1 S2 I1 C1 F0}, diff_note={concepts: graphs of image position against object position | steps: lens and mirror formulas; shape of each graph; match}]
Match the following graphs (in list-II) between position of object ($u$) and position of image ($v$).\par
{\small\setlength{\tabcolsep}{3pt}\begin{tabular}{|c|p{0.3\linewidth}|c|c|}
\hline
 & \textbf{List--I} & & \textbf{List--II}\\ \hline
(I) & Concave mirror & (P) & \bfigi{57.8}{TNAS07_p2_phy_s2_q07_1.png}\\ \hline
(II) & Convex mirror & (Q) & \bfigi{61.2}{TNAS07_p2_phy_s2_q07_2.png}\\ \hline
(III) & Convex lens & (R) & \bfigi{64.7}{TNAS07_p2_phy_s2_q07_3.png}\\ \hline
(IV) & Concave lens & (S) & \bfigi{69.2}{TNAS07_p2_phy_s2_q07_4.png}\\ \hline
 & & (T) & \bfigi{58.3}{TNAS07_p2_phy_s2_q07_5.png}\\ \hline
\end{tabular}}\par
\begin{qoptions}
\opt{A}{$\mathrm I\to S;\allowbreak \mathrm{II}\to P;\allowbreak \mathrm{III}\to R;\allowbreak \mathrm{IV}\to Q$}
\opt{B}{$\mathrm I\to R;\allowbreak \mathrm{II}\to Q;\allowbreak \mathrm{III}\to S;\allowbreak \mathrm{IV}\to P$}
\opt{C}{$\mathrm I\to P;\allowbreak \mathrm{II}\to S;\allowbreak \mathrm{III}\to Q;\allowbreak \mathrm{IV}\to R$}
\opt{D}{$\mathrm I\to Q;\allowbreak \mathrm{II}\to R;\allowbreak \mathrm{III}\to P;\allowbreak \mathrm{IV}\to S$}
\end{qoptions}
\end{bankq}
\begin{bankq}[id=TNAS07-P2-PHY-S2-Q08, source={TNAS - Test 07 (20 Aug 2023)}, paper=2, qno=8, type=single, subtopic={Reflection}, answer={A}, answer_src=key, difficulty=7, diff_parts={K1 S3 I2 C2 F0}, diff_note={concepts: moving object and image velocity components | steps: image distance; longitudinal magnification; axial velocity; transverse velocity; path of the image; line coefficients; match}]
An object lies on the principal axis of a convex mirror of focal length $f=1.0$ m at a distance $3f$ from it. At time $t=0$, the object starts moving with velocity components $v_0=1$ m/s along the principal axis (towards the mirror) and $v_0=1$ m/s perpendicular to the principal axis. The image moves along the line $ax+by=c$ (where $a$, $b$, $c$ are relatively prime integers and $x$ \& $y$ are in meters.) Match the quantities given in List--I with their magnitudes (in SI units) in List--II. We are only concerned with the time in which the rays from the object to the mirror are paraxial. Origin as well as $x$ \& $y$ axis are considered as per sign convention.
\bfig{158.7}{TNAS07_p2_phy_s2_q08_1.png}
{\small\setlength{\tabcolsep}{3pt}\begin{tabular}{|c|p{0.4\linewidth}|c|c|}
\hline
 & \textbf{List--I} & & \textbf{List--II}\\ \hline
(I) & Axial component of velocity of image at time $t=0$ & (P) & $1/3$\\ \hline
(II) & Transverse component of velocity of image at time $t=0$ & (Q) & $1/16$\\ \hline
(III) & The value of $1/a$ & (R) & $1/4$\\ \hline
(IV) & The value of $1/c$ & (S) & $1/6$\\ \hline
 & & (T) & $1/8$\\ \hline
\end{tabular}}\par
\begin{qoptions}
\opt{A}{$\mathrm I\to Q;\allowbreak \mathrm{II}\to R;\allowbreak \mathrm{III}\to R;\allowbreak \mathrm{IV}\to P$}
\opt{B}{$\mathrm I\to P;\allowbreak \mathrm{II}\to R;\allowbreak \mathrm{III}\to S;\allowbreak \mathrm{IV}\to Q$}
\opt{C}{$\mathrm I\to Q;\allowbreak \mathrm{II}\to S;\allowbreak \mathrm{III}\to S;\allowbreak \mathrm{IV}\to P$}
\opt{D}{$\mathrm I\to P;\allowbreak \mathrm{II}\to R;\allowbreak \mathrm{III}\to Q;\allowbreak \mathrm{IV}\to S$}
\end{qoptions}
\end{bankq}
\begin{bankq}[id=TNAS07-P2-PHY-S3-Q01, source={TNAS - Test 07 (20 Aug 2023)}, paper=2, qno=1, type=numerical, subtopic={Refraction}, answer={1.50}, answer_src=key, difficulty=5, diff_parts={K1 S2 I1 C1 F1}, diff_note={concepts: magnification of a hemispherical lens | steps: refraction at the curved surface; apparent size; ratio}]
If you place the flat side of a hemispherical glass object (lens) against a paper, it is possible to see a magnified image of the paper's surface around the center of the lens. How many times is the image bigger than the object if you are observing at a distance that is much bigger than the dimensions of the lens? The refractive index of glass is $n=1.5$.
\end{bankq}

\section*{Star Review (Other) papers}
\begin{bankq}[id=STAR03-P1-PHY-S2-Q07, source={Star Review (Other) - Test 03 (04 Dec 2023)}, paper=1, qno=7, type=multi, subtopic={Refraction}, answer={B,C,D}, answer_src=key, difficulty=6, diff_parts={K1 S2 I2 C1 F1}, diff_note={concepts: refraction and total reflection at three media | steps: apply Snell law at the first boundary; condition for total reflection at the second; deviation; check statements}]
A light ray is incident on a boundary as shown.\par
\bfig{139.6}{STAR03_p1_phy_s2_q07_1.png}
\begin{qoptions}
\opt{A}{If $\mu_2<\sqrt3$, light will enter medium-3 with deviation $30^\circ$.}
\opt{B}{If $\mu_2>\sqrt3$, light will enter medium-3 with deviation $30^\circ$.}
\opt{C}{If $\mu_2>\sqrt3$, light will never return to medium-1 again.}
\opt{D}{If $\mu_2<\sqrt3$, light will never return to medium-1.}
\end{qoptions}
\end{bankq}
\begin{bankq}[id=STAR03-P2-PHY-S1-Q04, source={Star Review (Other) - Test 03 (04 Dec 2023)}, paper=2, qno=4, type=multi, subtopic={Refraction}, answer={B,D}, answer_src=key, difficulty=6, diff_parts={K1 S2 I2 C1 F1}, diff_note={concepts: apparent depth and velocity for a fish and a bird | steps: apparent height seen from water; apparent depth seen from air; relative velocities; check each}]
A fish, $F$ in the pond, is at a depth of $0.8$ m from water surface and is moving vertically upwards with velocity $2\ \mathrm{ms^{-1}}$. At the same instant, a bird $B$ is at a height of $6$ m from water surface and is moving downwards with velocity $3\ \mathrm{ms^{-1}}$. At this instant both are on the same vertical lines as shown in the figure. Which of the following statement(s) is(are) correct?\par
\bfig{129.3}{STAR03_p2_phy_s1_q04_1.png}
\begin{qoptions}
\opt{A}{Height of $B$, observed by $F$ (from itself) is equal to $8.00$ m.}
\opt{B}{Depth of $F$, observed by $B$ (from itself) is equal to $6.60$ m.}
\opt{C}{Velocity of $B$, observed by $F$ (relative to itself) is equal to $5.00\ \mathrm{ms^{-1}}$.}
\opt{D}{Velocity of $F$, observed by $B$ (relative to itself) if equal to $4.50\ \mathrm{ms^{-1}}$.}
\end{qoptions}
\end{bankq}
\end{document}
